% ---------------------------------------------------------------
% TEMPLATE TCC2 - SENAC SANTO AMARO
% ---------------------------------------------------------------
% Trabalho de Conclusão de Curso 2
% Elaborado para apresentação do trabalho final do TCC
% Estrutura: Introdução, Revisão Bibliográfica, Metodologia,
%            Desenvolvimento, Resultados e Análise, Conclusão
% ---------------------------------------------------------------

\documentclass[
    12pt,                % tamanho da fonte
    openright,           % capítulos começam em página ímpar
    oneside,             % para impressão em apenas um lado do papel
    a4paper,             % tamanho do papel
    english,             % idioma adicional
    brazil               % idioma principal
]{abntex2}

% ---------------------------------------------------------------
% PACOTES
% ---------------------------------------------------------------
\usepackage{lmodern}                % Fonte Latin Modern
\usepackage[T1]{fontenc}            % Seleção de códigos de fonte
\usepackage[utf8]{inputenc}         % Codificação do documento
\usepackage{lastpage}               % Usado pela Ficha catalográfica
\usepackage{indentfirst}            % Indenta o primeiro parágrafo
\usepackage{color}                  % Controle de cores
\usepackage{graphicx}               % Inclusão de gráficos
\usepackage{microtype}              % Melhorias de justificação
\usepackage{float}                  % Para posicionamento de figuras
\usepackage{booktabs}               % Tabelas profissionais
\usepackage{multirow}               % Células mescladas em tabelas
\usepackage{longtable}              % Tabelas longas
\usepackage{amsmath,amssymb}        % Símbolos matemáticos
\usepackage{pgfgantt}               % Diagramas de Gantt
\usepackage[brazilian,hyperpageref]{backref}  % Páginas com citações
\usepackage[alf]{abntex2cite}       % Citações ABNT


% ---------------------------------------------------------------
% CONFIGURAÇÕES DO PDF
% ---------------------------------------------------------------
\makeatletter
\hypersetup{
    pdftitle={\@title},
    pdfauthor={\@author},
    pdfsubject={TCC2 - SENAC Santo Amaro},
    pdfcreator={LaTeX with abnTeX2},
    pdfkeywords={tcc2}{senac}{santo amaro}{trabalho de conclusão},
    colorlinks=true,        % false: links em caixas; true: links coloridos
    linkcolor=blue,         % cor dos links internos
    citecolor=blue,         % cor dos links para bibliografia
    filecolor=magenta,      % cor dos links para arquivos
    urlcolor=blue,
    bookmarksdepth=4
}
\makeatother

% ---------------------------------------------------------------
% CONFIGURAÇÕES DE APARÊNCIA DO PDF FINAL
% ---------------------------------------------------------------
\definecolor{blue}{RGB}{41,5,195}

% ---------------------------------------------------------------
% TAMANHO DA FONTE DAS LEGENDAS (ABNT) -- 10 pt
% ---------------------------------------------------------------
% O título (\caption) e a linha de fonte (\legend) de figuras e
% tabelas devem ficar em 10 pt (o corpo do texto é 12 pt). Por
% padrão o memoir/abntex2 deixa essas legendas em 12 pt; os comandos
% abaixo reduzem o título e a fonte para 10 pt.
\renewcommand{\ABNTEXfontereduzida}{\fontsize{10}{12}\selectfont}
\captionnamefont{\ABNTEXfontereduzida}   % rótulo "Figura N"/"Tabela N"
\captiontitlefont{\ABNTEXfontereduzida}  % texto do título da legenda
\makeatletter
% Redefine \legend (usado em "Fonte: ...") para sair em 10 pt
% (a definição original do memoir fixa \normalsize = 12 pt).
\renewcommand{\legend}[1]{%
  \M@gettitle{#1}%
  \memlegendinfo{#1}%
  \par
  \begingroup
    \@parboxrestore
    \if@minipage
      \@setminipage
    \fi
    \ABNTEXfontereduzida
    \captiondelim{\mbox{}}%
    \@makecaption{}{\ignorespaces #1}\par
  \endgroup}
\makeatother

% ---------------------------------------------------------------
% INFORMAÇÕES DO DOCUMENTO
% ---------------------------------------------------------------
\titulo{Detecção de Ataques DDoS em Redes de Computadores utilizando Algoritmos Supervisionados de Machine Learning: uma Análise Comparativa entre Random Forest e Support Vector Machine}
\autor{Bianca Lima Rodrigues e Bruno José Oliveira Gomes}
\local{São Paulo}
\data{2026}
\orientador{Prof. Afonso Cesar Lelis Brandão}
% \coorientador{Prof. Nome do Coorientador} % Descomente se houver coorientador

\instituicao{%
  SENAC -- Serviço Nacional de Aprendizagem Comercial
  \par
  Unidade Santo Amaro
  \par
  Curso de Ciência da Computação
}

\tipotrabalho{Trabalho de Conclusão de Curso II}

\preambulo{Trabalho de Conclusão de Curso II apresentado ao SENAC Santo Amaro como requisito parcial para conclusão do curso de Ciência da Computação.}


% ---------------------------------------------------------------
% COMPILA O ÍNDICE
% ---------------------------------------------------------------
\makeindex

\makeatletter
\def\fps@figure{tbp}
\makeatother


\floatplacement{figure}{H}
\floatplacement{table}{H}

% ---------------------------------------------------------------
% INÍCIO DO DOCUMENTO
% ---------------------------------------------------------------
\begin{document}

% Seleciona o idioma do documento
\selectlanguage{brazil}

% Retira espaço extra obsoleto entre as frases
\frenchspacing

% ---------------------------------------------------------------
% ELEMENTOS PRÉ-TEXTUAIS
% ---------------------------------------------------------------
\pretextual

% ---------------------------------------------------------------
% CAPA
% ---------------------------------------------------------------
\imprimircapa

% ---------------------------------------------------------------
% FOLHA DE ROSTO
% ---------------------------------------------------------------
\imprimirfolhaderosto

% ---------------------------------------------------------------
% RESUMO
% ---------------------------------------------------------------
\begin{resumo}
Este documento apresenta o Trabalho de Conclusão de Curso (TCC2) desenvolvido no SENAC Santo Amaro, com base na estrutura recomendada por Raul Sidnei Wazlawick em \textit{Metodologia de Pesquisa para Ciência da Computação}. O foco está em evidenciar o problema, os objetivos, o método, o desenvolvimento, os resultados e a contribuição científica do trabalho, evitando textos meramente descritivos. O documento está organizado em capítulos que apresentam a introdução ao tema, a revisão bibliográfica que fundamenta a pesquisa, a metodologia adotada, o desenvolvimento da solução, os resultados e a respectiva análise crítica, seguidos da conclusão e das referências bibliográficas.

\vspace{\onelineskip}

\noindent
\textbf{Palavras-chave}: palavra1. palavra2. palavra3. palavra4. palavra5.
\end{resumo}

% ---------------------------------------------------------------
% ABSTRACT (RESUMO EM LÍNGUA ESTRANGEIRA) -- OBRIGATÓRIO NO TCC2
% ---------------------------------------------------------------
% Versão em inglês do resumo. Deve conter o mesmo conteúdo do resumo
% em português (problema, objetivo, método, resultados e conclusão)
% e as palavras-chave traduzidas (Keywords). Pela ABNT NBR 14724, o
% resumo em língua estrangeira vem logo após o resumo em português.
\begin{resumo}[Abstract]
 \begin{otherlanguage*}{english}
   This document presents the undergraduate final project (TCC2) developed at SENAC Santo Amaro, based on the structure recommended by Raul Sidnei Wazlawick in \textit{Metodologia de Pesquisa para Ciência da Computação}. The focus is on highlighting the problem, the objectives, the method, the development, the results and the scientific contribution of the work, avoiding merely descriptive texts. The document is organized into chapters presenting the introduction to the topic, the literature review that grounds the research, the adopted methodology, the development of the solution, the results and their critical analysis, followed by the conclusion and the bibliographic references.

   \vspace{\onelineskip}

   \noindent
   \textbf{Keywords}: keyword1. keyword2. keyword3. keyword4. keyword5.
 \end{otherlanguage*}
\end{resumo}

% ---------------------------------------------------------------
% LISTA DE ILUSTRAÇÕES (FIGURAS)
% ---------------------------------------------------------------
% Gerada automaticamente a partir das \caption{} das figuras.
% Toda figura com legenda passa a constar nesta lista.
\pdfbookmark[0]{\listfigurename}{lof}
\listoffigures*
\cleardoublepage

% ---------------------------------------------------------------
% LISTA DE TABELAS
% ---------------------------------------------------------------
% Gerada automaticamente a partir das \caption{} das tabelas.
\pdfbookmark[0]{\listtablename}{lot}
\listoftables*
\cleardoublepage

% ---------------------------------------------------------------
% LISTA DE ABREVIATURAS E SIGLAS
% ---------------------------------------------------------------
% Esta lista deve vir antes do início do trabalho (elementos pré-textuais).
% Em ordem alfabética, relacione APENAS as abreviaturas e siglas que
% realmente aparecem no texto. Edite a tabela conforme o seu trabalho.
\pdfbookmark[0]{Lista de Abreviaturas e Siglas}{loas}
\chapter*{Lista de Abreviaturas e Siglas}

\noindent A lista a seguir reúne, em ordem alfabética, as abreviaturas e siglas utilizadas neste trabalho, acompanhadas de seus respectivos significados. Substitua, acrescente ou remova linhas conforme a necessidade do seu projeto.

\vspace{\onelineskip}

\begin{center}
\begin{tabular}{@{}p{3.2cm}p{0.6\textwidth}@{}}
\toprule
\textbf{Abreviatura/Sigla} & \textbf{Significado} \\
\midrule
ABNT  & Associação Brasileira de Normas Técnicas \\
API   & \textit{Application Programming Interface} (Interface de Programação de Aplicações) \\
HTTP  & \textit{Hypertext Transfer Protocol} \\
REST  & \textit{Representational State Transfer} \\
SQL   & \textit{Structured Query Language} \\
TCC   & Trabalho de Conclusão de Curso \\
UML   & \textit{Unified Modeling Language} \\
\bottomrule
\end{tabular}
\end{center}

\cleardoublepage

% ---------------------------------------------------------------
% SUMÁRIO
% ---------------------------------------------------------------
\pdfbookmark[0]{\contentsname}{toc}
\tableofcontents*
\cleardoublepage

% ---------------------------------------------------------------
% ELEMENTOS TEXTUAIS
% ---------------------------------------------------------------
\textual

% ===============================================================
% CAPÍTULO 1 - INTRODUÇÃO
% ===============================================================
\chapter{Introdução}
\label{cap:introducao}
\section{Introdução}

Segundo o relatório \textit{Cyber Security Report}, da Check Point Research \cite{CheckPoint2026}, as organizações sofreram, em média, 1.968 ataques cibernéticos semanais em 2025, um aumento de 70\% em relação a 2023. Além disso, em fevereiro de 2026, a média de ataques por organização cresceu 9,6\% em comparação ao mesmo período do ano anterior, evidenciando o aumento das ameaças às redes de computadores.

A Figura~\ref{fig:ataques-latam} apresenta a evolução da quantidade de ataques cibernéticos registrados nos cinco países mais visados da América Latina ao longo de 2025. Observa-se que o Brasil se destaca em relação aos demais países, mantendo o maior volume de ocorrências durante todo o período analisado. Esses dados evidenciam a relevância do cenário brasileiro no contexto das ameaças cibernéticas e reforçam a necessidade de pesquisas voltadas ao desenvolvimento de mecanismos eficientes para a detecção de ataques em redes de computadores.

\begin{figure}
    \centering
    \caption{Quantidade de ataques cibernéticos nos cinco países mais visados da América Latina em 2025}
    \label{fig:ataques-latam}
    \includegraphics[width=0.90\textwidth]{assets/ataquesLATAM.png}
    \legend{Fonte: \citeonline{NETSCOUT2025}}
\end{figure}

Dentre essas ameaças, destacam-se os ataques de negação de serviço distribuído (DDoS), que representam um risco significativo diante da crescente dependência de serviços digitais. Nesse cenário, a disponibilidade dos sistemas é essencial para a continuidade das operações das organizações e o atendimento aos usuários, uma vez que interrupções podem gerar prejuízos financeiros e impedir o acesso a serviços. Os ataques DDoS exploram essa vulnerabilidade ao sobrecarregar servidores e recursos de rede, dificultando o acesso de usuários legítimos. Segundo a NETSCOUT ATLAS \cite{NETSCOUT2025}, mais de 9,1 milhões de ataques DDoS foram registrados mundialmente no primeiro semestre de 2026.

Atualmente, a detecção de ataques DDoS pode ser realizada por meio de Sistemas de Detecção de Intrusão (IDS), que monitoram o tráfego da rede em busca de ameaças conhecidas ou comportamentos suspeitos \cite{najafimehr2023}. Entretanto, essas técnicas podem apresentar dificuldades para identificar ataques desconhecidos e gerar falsos positivos \cite{najafimehr2023,mahjabin2017}.

Nesse contexto, o Aprendizado de Máquina (\textit{Machine Learning}) surge como uma alternativa para aprimorar a detecção, permitindo que modelos aprendam a identificar padrões associados a atividades maliciosas. Conforme apontam Apostu et al. \cite{apostu2026}, essa abordagem representa uma transição dos métodos tradicionais, pautados em regras preestabelecidas, para modelos que empregam técnicas de inteligência artificial. Diante disso, este trabalho propõe uma análise comparativa entre dois algoritmos supervisionados de aprendizado de máquina: Random Forest (RF) e Support Vector Machine (SVM), avaliando seu desempenho na detecção de ataques DDoS.

% ---------------------------------------------------------------
\section{Contextualização do problema}
\label{sec:contexto}

\subsection{Contexto}

A crescente presença da tecnologia nas atividades cotidianas e nos processos organizacionais tem intensificado a dependência de serviços digitais em setores como comércio, educação, saúde e administração pública. Nesse cenário, a disponibilidade das redes e infraestruturas computacionais torna-se fundamental para a continuidade dessas atividades. Entretanto, a ampliação da conectividade também expõe instituições públicas e privadas a ameaças cibernéticas, especialmente ataques que comprometem a disponibilidade de seus serviços, como os ataques distribuídos de negação de serviço (\textit{Distributed Denial-of-Service} -- DDoS) \cite{ouhssini2024,thaqi2026}.

Esses ataques consistem na sobrecarga intencional de servidores, aplicações ou redes por meio do envio massivo de tráfego proveniente de múltiplas fontes comprometidas, com o objetivo de interromper a disponibilidade dos serviços \cite{ibm2026}. Sua execução pode envolver computadores, telefones celulares e dispositivos de Internet das Coisas (IoT), cuja crescente presença nos ambientes conectados amplia a quantidade de equipamentos que podem ser explorados pelos atacantes \cite{ibm2026}.

A ocorrência desses ataques pode provocar interrupções de serviços essenciais, prejuízos financeiros e operacionais, danos à reputação das organizações e impactos diretos aos usuários \cite{ouhssini2024}. Casos como o ataque do tipo \textit{Distributed Reflective Denial of Service} (DRDoS) direcionado à infraestrutura da Amazon Web Services (AWS) demonstram o potencial dessas ameaças de afetar milhões de usuários \cite{subramani2021}. Além disso, a frequência e a complexidade dos ataques têm aumentado com a evolução das ferramentas de Inteligência Artificial (IA) e \textit{Machine Learning} (ML), que permitem aos atacantes identificar vulnerabilidades e adaptar suas estratégias para contornar mecanismos de mitigação \cite{ibm2026,thaqi2026}.

Diante da evolução dos ataques DDoS e das limitações dos mecanismos tradicionais de detecção, torna-se necessário desenvolver soluções mais dinâmicas e adaptáveis, capazes de identificar ameaças em tempo real e acompanhar a evolução dos padrões de tráfego malicioso \cite{thaqi2026}. As técnicas de \textit{Machine Learning} apresentam potencial nesse cenário, pois permitem analisar grandes volumes de tráfego e identificar padrões que indiquem comportamentos maliciosos, contribuindo para superar limitações dos métodos tradicionais e aprimorar a proteção das redes de computadores \cite{ouhssini2024}.
% ---------------------------------------------------------------
\section{Formulação do problema}
\label{sec:problema}

Embora técnicas de \textit{Machine Learning} apresentem potencial para aprimorar a detecção de ataques DDoS, diferentes algoritmos podem apresentar desempenhos distintos quanto à capacidade de identificação de ameaças e ao custo computacional. Diante disso, questiona-se: \textbf{qual dos algoritmos, entre Random Forest e Support Vector Machine (SVM), apresenta melhor desempenho na detecção de ataques DDoS, considerando métricas de classificação e custo computacional, em cenários de classificação binária e multiclasse?}

% ---------------------------------------------------------------
\section{Objetivos e hipótese}
\label{sec:objetivos}

O objetivo deste trabalho visa identificar qual o melhor desempenho entre dois modelos diferentes de aprendizado de máquina, no quesito detecção de ataques DDoS em redes de computadores.

\subsection{Objetivo geral}

Comparar e analisar dois modelos diferentes supervisionados de Machine Learning: Random Forest e Support Vector Machine, avaliando suas eficácias na identificação de ataques DDoS em redes de computadores.

\subsection{Objetivos específicos}

\begin{enumerate}
    \item Implementar e treinar os algoritmos Random Forest e Support Vector Machine (SVM), estabelecendo uma base experimental para a comparação de seu desempenho na detecção de ataques DDoS;
    \item Avaliar o desempenho preditivo dos modelos antes e após a otimização de hiperparâmetros, utilizando técnicas de busca de parâmetros e validação cruzada (\textit{Cross-Validation}) para identificar as configurações mais adequadas;
    \item Medir o desempenho dos modelos nas tarefas de classificação binária e multiclasse de fluxos de rede, considerando a eficiência computacional, o tempo de treinamento e os resultados obtidos;
    \item Comparar os resultados experimentais com os estudos presentes no estado da arte, considerando métricas de avaliação como acurácia, precisão, revocação, F1-score e MCC.
\end{enumerate}


\subsection{Hipótese ou critério de validação}

Parte-se da hipótese de que os algoritmos Random Forest e Support Vector Machine (SVM) apresentam diferenças de desempenho na detecção de ataques DDoS, considerando as métricas de acurácia, precisão, revocação, F1-score e MCC, além do custo computacional. Espera-se que essas diferenças variem conforme o tipo de classificação, binária ou multiclasse, e que a otimização de hiperparâmetros contribua para aprimorar os resultados preditivos, embora seus efeitos possam variar entre os modelos. A comparação permitirá identificar as vantagens e limitações de cada algoritmo quanto à eficácia na detecção de ataques e à eficiência computacional, considerando também os tempos de treinamento e predição. 

% ---------------------------------------------------------------
\section{Justificativa}
\label{sec:justificativa}

Com a evolução da infraestrutura da inteligência artificial e o aumento da complexidade dos ambientes computacionais, os ataques cibernéticos, especialmente os ataques distribuídos de negação de serviço (DDoS), têm se tornado mais sofisticados, frequentes e difíceis de mitigar, exigindo soluções cada vez mais rápidas e inteligentes para sua detecção \cite{thaqi2026}. Diferentemente das abordagens tradicionais baseadas em assinaturas, os modelos de aprendizado de máquina possibilitam a análise de grandes volumes de dados e a identificação de padrões anômalos, permitindo maior capacidade de adaptação frente à constante evolução dos ataques cibernéticos \cite{ouhssini2024}.

\begin{itemize}
    \item \textbf{Relevância acadêmica:} Sob a perspectiva acadêmica, este trabalho contribui para o avanço das pesquisas na interseção entre segurança da informação e inteligência artificial, explorando a aplicação de técnicas de aprendizado de máquina na detecção de ataques DDoS. Além disso, busca ampliar a compreensão sobre o comportamento e a eficácia de diferentes algoritmos de classificação, considerando suas vantagens, limitações e custos computacionais. Essa investigação é relevante diante da necessidade de aprimorar os mecanismos de defesa em ambientes computacionais, conforme apontado por Ouhssini et al. \cite{ouhssini2024}, que destacam a importância de desenvolver soluções mais robustas para a detecção desses ataques. Dessa forma, o estudo também contribui para o conhecimento científico ao fornecer uma análise comparativa que pode auxiliar na escolha de modelos de aprendizado de máquina mais adequados a diferentes cenários de detecção de ameaças.

    \item \textbf{Relevância prática/profissional:} Do ponto de vista prático e profissional, o desenvolvimento de modelos inteligentes para detecção de intrusão pode proporcionar benefícios significativos, como a identificação de padrões maliciosos, a antecipação de ataques e a redução do tempo de resposta a incidentes \cite{thaqi2026}. Essas soluções podem contribuir para prevenir interrupções operacionais, preservar a disponibilidade dos serviços e fortalecer a segurança das infraestruturas de rede \cite{ouhssini2024}. Dessa forma, podem favorecer a continuidade das atividades organizacionais, reduzir os impactos causados por ataques DDoS e proporcionar serviços digitais mais seguros e confiáveis para os usuários.

    \item \textbf{Lacuna identificada:} Embora os modelos de linguagem de grande porte (LLMs) representem uma abordagem emergente para a segurança de redes, sua aplicação na detecção de ataques DDoS ainda apresenta desafios. Esses modelos possuem capacidade de analisar dados de tráfego de rede e identificar padrões associados a atividades maliciosas, utilizando técnicas como \textit{few-shot learning} e \textit{fine-tuning} para aprimorar sua classificação \cite{guastalla2024}. Entretanto, seu elevado consumo computacional, energético e de memória pode dificultar a implementação em ambientes com recursos limitados \cite{argerich2024}. Nesse contexto, algoritmos tradicionais de \textit{Machine Learning}, como Random Forest e Support Vector Machine (SVM), representam alternativas relevantes para tarefas específicas de classificação e detecção de anomalias. Dessa forma, justifica-se a realização deste estudo, que busca comparar o desempenho desses algoritmos e contribuir para a identificação de abordagens eficientes e adequadas à detecção de ataques DDoS.
\end{itemize}

% ===============================================================
% CAPÍTULO 2 - REVISÃO BIBLIOGRÁFICA / TRABALHOS RELACIONADOS
% ===============================================================
\chapter{Revisão Bibliográfica / Trabalhos Relacionados}
\label{cap:revisao}

% ---------------------------------------------------------------
\section{Metodologia da pesquisa bibliográfica}
\label{sec:metodologia_pesquisa_bib}

A pesquisa bibliográfica foi conduzida com o objetivo de reunir estudos relevantes sobre sistemas de detecção de intrusão em redes (IDS), ataques distribuídos de negação de serviço (DDoS) e a aplicação de técnicas de \textit{Machine Learning} supervisionado para classificação de tráfego malicioso.

\begin{itemize}
    \item \textbf{Bases de dados consultadas:} Foram consultadas bases de dados reconhecidas nas áreas de computação e segurança da informação, incluindo Google Scholar, IEEE Xplore, Papers with Code, Directory of Open Access Journals (DOAJ), SciELO e Semantic Scholar, buscando garantir diversidade e qualidade das fontes analisadas.

    \item \textbf{Palavras-chave utilizadas:} As buscas foram realizadas com palavras-chave em português e inglês relacionadas ao tema da pesquisa. Entre os principais termos utilizados estão: ``machine learning'', ``aprendizado de máquina'', ``supervised learning'', ``aprendizado supervisionado'', ``intrusion detection system'', ``sistema de detecção de intrusão'', ``DDoS attack detection'', ``detecção de ataques DDoS'', ``network anomaly detection'', ``detecção de anomalias em redes'', ``Random Forest'', ``Support Vector Machine'', ``SVM'', ``artificial intelligence'', ``inteligência artificial'', ``cybersecurity'', ``network security'', ``CICDDoS2019'' e ``CICDDoS2019 dataset''. A combinação desses termos permitiu localizar estudos que abordassem diretamente a aplicação de algoritmos supervisionados na identificação de tráfego malicioso em redes computacionais.

    \item \textbf{Exemplos de strings de busca utilizadas:}
    \begin{itemize}
        \item Google Scholar: \texttt{("DDoS attack detection" OR "detecção de ataques DDoS") AND ("machine learning" OR "aprendizado de máquina") AND ("supervised learning" OR "aprendizado supervisionado") AND ("Random Forest" OR "Support Vector Machine" OR SVM) AND ("network security" OR cybersecurity)}
        \item IEEE Xplore: \texttt{("intrusion detection system") AND ("network anomaly detection") AND (cybersecurity) AND ("CICDDoS2019")}
    \end{itemize}

    \item \textbf{Critérios de inclusão:} Foram selecionados artigos científicos, capítulos de livros e publicações técnicas em português e inglês que apresentassem relação direta com o problema investigado, especialmente estudos envolvendo \textit{Machine Learning} aplicado à detecção de intrusões e ataques DDoS, com ênfase nos algoritmos Random Forest e Support Vector Machine. Foram priorizados trabalhos com validação experimental, uso de conjuntos de dados reconhecidos na literatura, como o CICDDoS2019, e apresentação de métricas de avaliação como acurácia, precisão, \textit{recall}, F1-score e MCC.

    \item \textbf{Critérios de exclusão:} Foram descartados trabalhos duplicados, estudos que não abordavam diretamente os modelos de \textit{Machine Learning} escolhidos, artigos com foco exclusivo em técnicas de aprendizado profundo sem relação com os algoritmos analisados neste trabalho, além de materiais desatualizados ou com metodologia insuficientemente descrita. Dessa forma, a pesquisa bibliográfica permitiu construir uma base teórica consistente para fundamentar o desenvolvimento experimental e a análise comparativa proposta neste estudo.

    \item \textbf{Período de abrangência:} Foram considerados artigos publicados entre 2020 e 2026, salvo artigos e livros seminais ou fundamentais.
\end{itemize}

\begin{figure}
    \centering
    \caption{Critérios de exclusão metodologia da pesquisa bibliográfica}
    \label{fig:filtro-criterios}
    \includegraphics[width=0.75\textwidth]{assets/filtroCriterios.png}
    \legend{Fonte: Elaborado pelo autor (2026)}
\end{figure}

% ---------------------------------------------------------------
\section{Fundamentos teóricos}
\label{sec:fundamentos}
\subsubsection{Redes de computadores}

Segundo Tanenbaum e Wetherall \cite{tanenbaum2011}, um conjunto de dispositivos de computação autônomos interconectados é denominado ``redes de computadores'', estando eles interconectados quando é possível realizar a troca de dados entre eles. A transmissão se dá por meio de protocolos de comunicação, organizados em camadas no modelo TCP/IP. As aplicações de redes de computadores podem ser categorizadas em Aplicações Comerciais, cujo objetivo é que os dados estejam disponíveis para qualquer pessoa na rede. Já a Aplicação Doméstica visa abranger usuários em casa, oferecendo conectividade externa e acesso a uma vasta gama de recursos. Usuários Móveis são a categoria que utiliza a rede em trânsito (em terra, mar ou ar), por meio de redes sem fio. Por fim, a categoria relacionada a Questões Sociais é focada em cidadãos comuns que distribuem e visualizam conteúdos entre si \cite{tanenbaum2011}.

Os dispositivos interconectados realizam a troca de informações por meio de protocolos de comunicação, definidos como um conjunto de regras e padrões que permitem a comunicação adequada entre dispositivos em uma rede. Esses protocolos são organizados em modelos de referência que estruturam as funções de comunicação em uma pilha de camadas. O modelo OSI é considerado um modelo teórico, utilizado como base conceitual para compreender as responsabilidades de cada camada no processo de comunicação. Em contrapartida, o modelo TCP/IP possui caráter prático, sendo amplamente utilizado na Internet e nas redes de computadores atuais para viabilizar a comunicação e a troca de dados entre diferentes sistemas \cite{tanenbaum2011}.

\subsubsection{Modelo TCP/IP}

Segundo Stevens \cite{stevens1994}, o modelo TCP/IP consiste em um conjunto de protocolos que possibilita a comunicação e a troca de dados entre dispositivos conectados em rede, independentemente de diferenças de hardware, fabricante ou sistema operacional. Seu nome deriva de dois de seus principais protocolos, o TCP (\textit{Transmission Control Protocol}) e o IP (\textit{Internet Protocol}), embora a suíte seja composta por diversos outros protocolos que, em conjunto, formam a base da comunicação na Internet. Sua arquitetura é organizada em quatro camadas: \textbf{Aplicação, Transporte, Internet e Interface de Rede}, cada uma responsável por funções específicas durante a comunicação.

A camada de Aplicação fornece serviços utilizados pelas aplicações, abrangendo protocolos como HTTP, FTP, SMTP e DNS. A camada de Transporte realiza a comunicação fim a fim, principalmente por meio do TCP, que oferece uma comunicação orientada à conexão e mecanismos de controle de erros e retransmissão, e do UDP, que prioriza uma comunicação mais simples e rápida, sem garantia de entrega. Na sequência, a camada de Internet é responsável pelo endereçamento e roteamento dos pacotes, tendo o IP como seu principal protocolo, enquanto a camada de Interface de Rede realiza a transmissão dos dados pelo meio de comunicação, utilizando tecnologias como Ethernet \cite{stevens1994,esr2020}.

Durante a transmissão, os dados percorrem essas camadas em um processo de \textbf{encapsulamento}, no qual são adicionadas informações necessárias à comunicação. No dispositivo de destino, ocorre o processo inverso, denominado desencapsulamento. Essa organização em camadas promove padronização e interoperabilidade entre diferentes tecnologias e sistemas, características que contribuíram para a adoção do TCP/IP como principal arquitetura de comunicação da Internet \cite{stevens1994,esr2020}.

Essas camadas e seus respectivos protocolos são relevantes para a compreensão dos ataques DDoS, uma vez que diferentes modalidades de ataque exploram serviços e mecanismos de comunicação associados a elas \cite{sharafaldin2019}.

\paragraph{IP -- Internet Protocol}

O IP (\textit{Internet Protocol}) é o principal protocolo da camada de Internet do modelo TCP/IP, responsável pelo \textbf{endereçamento e encaminhamento de pacotes entre dispositivos e redes}. Por meio dos endereços IP de origem e destino, os pacotes podem ser identificados e direcionados até o dispositivo correto, possibilitando a comunicação entre equipamentos localizados na mesma rede ou em redes distintas \cite{tanenbaum2011,stevens1994}. Durante essa comunicação, os pacotes transportam informações que permitem identificar características do tráfego gerado na rede. Dessa forma, compreender o funcionamento do protocolo IP contribui para a análise do tráfego de rede e, consequentemente, para a identificação de comportamentos e padrões relacionados a ataques de negação de serviço \cite{sharafaldin2019}.

\paragraph{TCP -- Transmission Control Protocol}

O TCP (\textit{Transmission Control Protocol}) é um protocolo da camada de transporte da arquitetura TCP/IP, responsável por estabelecer uma comunicação confiável entre dispositivos conectados em uma rede. Diferentemente do protocolo IP, o TCP oferece mecanismos para garantir a entrega dos dados, como controle de fluxo, retransmissão de segmentos perdidos e reorganização de informações recebidas fora de ordem \cite{tanenbaum2011}.

A comunicação TCP ocorre por meio de conexões ponto a ponto, identificadas por \textit{sockets}, compostos pelo endereço IP e pelo número de porta dos dispositivos envolvidos. Os dados são divididos em segmentos numerados sequencialmente, permitindo que o dispositivo receptor os reorganize e confirme seu recebimento por meio de mensagens de reconhecimento (ACKs) \cite{tanenbaum2011}.

O estabelecimento de uma conexão TCP utiliza o mecanismo denominado \textit{Three-Way Handshake}, composto por três etapas: inicialmente, o cliente envia um segmento SYN solicitando a conexão; em seguida, o servidor responde com SYN-ACK; por fim, o cliente envia um ACK, confirmando o recebimento da resposta. Após esse processo, a conexão é estabelecida e a transmissão de dados pode ser iniciada. Esse mecanismo é particularmente relevante para o estudo de ataques DDoS, como o SYN Flood, que explora o processo de estabelecimento de conexões TCP \cite{sharafaldin2019}.

\subsubsection{Cibersegurança}

Segundo Schneier \cite{schneier1996}, a proteção de sistemas computacionais, dados e comunicações tornou-se essencial diante do crescimento da utilização de tecnologias digitais e da necessidade de garantir a confidencialidade, a integridade e a disponibilidade das informações. Em sua abordagem, o autor discute mecanismos de proteção voltados à segurança de computadores, à segurança de dados e à segurança das comunicações, destacando a importância da proteção contra acessos não autorizados, ataques e demais ameaças aos sistemas computacionais.

Embora o termo cibersegurança não seja utilizado diretamente por Schneier na obra publicada em 1996, atualmente esses conceitos são amplamente englobados por essa área. A cibersegurança pode ser compreendida como uma área da segurança da informação voltada à proteção de sistemas, redes, dispositivos e informações digitais contra ameaças cibernéticas, atuando na preservação da privacidade e da disponibilidade dos dados em ambientes computacionais \cite{stallings2022}.

\paragraph{Segurança de redes}

Segundo Stallings \cite{stallings2022}, a cibersegurança e a segurança de redes estão diretamente relacionadas, uma vez que a segurança de redes representa um dos principais mecanismos de proteção no ambiente digital. Enquanto a cibersegurança possui um escopo mais amplo, voltado à proteção de sistemas, aplicações, dispositivos e informações digitais contra ameaças cibernéticas, a segurança de redes concentra-se especificamente na proteção da infraestrutura de comunicação e do tráfego de dados entre dispositivos conectados.

Nesse contexto, mecanismos como criptografia, autenticação, controle de acesso, firewalls e sistemas de detecção de intrusão desempenham papel fundamental na prevenção de ataques e na preservação da confidencialidade, integridade e disponibilidade das informações. Segundo Tanenbaum e Wetherall \cite{tanenbaum2011} e Stallings \cite{stallings2022}, a segurança de redes representa um dos maiores desafios da computação moderna, considerando que atividades sensíveis são constantemente processadas e transmitidas entre dispositivos.

\subsubsection{Sistemas de detecção de intrusão (IDS)}

Um Sistema de Detecção de Intrusão (\textit{Intrusion Detection System} -- IDS) é uma tecnologia fundamental para o fortalecimento da segurança em ambientes computacionais, atuando como uma camada complementar aos mecanismos preventivos tradicionais, como firewalls e sistemas de autenticação \cite{stallings2020}. De acordo com Stallings \cite{stallings2020}, o IDS desempenha o papel de uma linha de defesa focada no monitoramento contínuo das atividades em redes e sistemas, com o propósito de analisar eventos e emitir avisos em tempo real sobre tentativas de acesso não autorizado, identificando intrusos. Seu funcionamento baseia-se na premissa de que o comportamento de um invasor difere do comportamento de um usuário legítimo de formas que podem ser quantificadas e mensuradas \cite{stallings2020}.

Entre os principais métodos empregados pelos sistemas IDS estão:

\begin{itemize}
    \item \textbf{Detecção baseada em assinaturas:} identifica ataques por meio da comparação do tráfego com padrões e regras de ameaças conhecidas. Apresenta boa precisão e baixa taxa de falsos alarmes, mas possui limitações na identificação de ataques inéditos (\textit{zero-day}) \cite{stallings2020}.

    \item \textbf{Detecção baseada em anomalias:} identifica atividades que se desviam do comportamento normal da rede, permitindo reconhecer ameaças desconhecidas. Para isso, utiliza perfis de comportamento e tráfego como referência para detectar alterações suspeitas \cite{stallings2020}.
\end{itemize}

A detecção por anomalias relaciona-se ao Aprendizado de Máquina (\textit{Machine Learning}), cujos modelos podem aprender padrões de tráfego e identificar desvios sem depender exclusivamente de regras predefinidas. Entretanto, essa abordagem exige equilíbrio entre a identificação de ameaças e a ocorrência de falsos positivos e falsos negativos \cite{gopalakrishnan2023}.

Além de detectar incidentes, os IDS contribuem para a coleta de evidências sobre os ataques e o aprimoramento das estratégias de proteção da infraestrutura computacional \cite{stallings2020}.

\subsubsection{Ataques de negação de serviço (DoS e DDoS)}

\paragraph{DoS -- Denial of Service}

Um ataque DoS (\textit{Denial of Service} -- negação de serviço) é um tipo de ataque cibernético caracterizado por tentativas explícitas de sobrecarregar um servidor, computador ou serviço de rede, tornando seus recursos indisponíveis para usuários legítimos, conforme descrito por Ouhssini et al. \cite{ouhssini2024}. Em muitos casos, esses ataques são utilizados como forma de extorsão (\textit{ransom}), na qual os criminosos exigem pagamentos para interromper o ataque e restabelecer o funcionamento do serviço \cite{ouhssini2024}.

Nesse tipo de ataque, o invasor envia um grande volume de requisições ou pacotes ao alvo, causando uma sobrecarga que força o sistema a consumir excessivamente recursos como processamento, memória e largura de banda. Como consequência, o sistema deixa de responder adequadamente às solicitações legítimas, causando lentidão ou até mesmo a indisponibilidade total do serviço. Os alvos mais comuns desse tipo de ataque são servidores web, nos quais o objetivo do atacante é tornar páginas e serviços inacessíveis \cite{ouhssini2024}.

\paragraph{DDoS -- Distributed Denial of Service}

Segundo Ouhssini et al. \cite{ouhssini2024}, o ataque DDoS (\textit{Distributed Denial of Service} -- negação de serviço distribuída) representa uma evolução do ataque DoS. Nesse caso, o invasor utiliza múltiplos dispositivos comprometidos, conhecidos como ``zumbis'' ou \textit{bots}, infectados por malwares (códigos maliciosos), para inundar simultaneamente o alvo com tráfego falso e solicitações excessivas. Esses dispositivos infectados formam uma rede chamada \textit{botnet}, controlada por um computador mestre operado pelo atacante. Dessa forma, milhares de máquinas podem acessar continuamente um mesmo servidor ou serviço ao mesmo tempo, esgotando seus recursos e impossibilitando o atendimento de novos usuários \cite{ouhssini2024,thaqi2026}.

Segundo Ouhssini et al. \cite{ouhssini2024}, os ataques DDoS podem ser classificados em duas categorias principais: ataques no nível de rede/transporte e ataques no nível de aplicação. Os ataques de rede/transporte têm como objetivo esgotar a largura de banda ou os recursos de conectividade da vítima, incluindo técnicas como UDP Flood, ICMP Flood e TCP SYN Flood. Destaca-se o ataque SYN Flood, que explora falhas específicas na implementação do protocolo TCP para consumir recursos do servidor. Nesse ataque, o invasor envia inúmeras solicitações falsas de \textit{handshake} TCP utilizando endereços IP falsificados, fazendo com que o servidor tente responder a conexões que nunca serão concluídas, consumindo seus recursos e sobrecarregando o sistema \cite{thaqi2026}.

Já os ataques de aplicação têm como alvo serviços web e aplicações HTTP (\textit{HyperText Transfer Protocol} -- protocolo de transferência de hipertexto), buscando sobrecarregar recursos como CPU, memória, banco de dados e conexões simultâneas. Entre os exemplos mais comuns estão os ataques de inundação HTTP (\textit{HTTP Flood}), nos quais um grande número de requisições força o servidor a executar mais operações do que consegue suportar, causando lentidão ou indisponibilidade do serviço \cite{ouhssini2024,thaqi2026}. A Figura 3 ilustra o funcionamento de um ataque distribuído de negação de serviço (\textit{Distributed Denial of Service} -- DDoS).


\begin{figure}
    \centering
    \caption{Arquitetura de um ataque DDoS}
    \label{fig:ataque-ddos}
    \includegraphics[width=0.6\textwidth]{assets/ataqueDDoS.png}
    \legend{Fonte: \citeonline{Apposite2025}}
\end{figure}


\subsubsection{CICDDoS2019}

O conjunto de dados CICDDoS2019, desenvolvido pelo Canadian Institute for Cybersecurity (CIC) da University of New Brunswick, foi criado com o objetivo de fornecer uma base moderna e realista para pesquisas relacionadas à detecção de ataques distribuídos de negação de serviço (DDoS) utilizando técnicas de aprendizado de máquina e mineração de dados \cite{sharafaldin2019}. Esse conjunto de dados foi projetado para representar cenários atuais de tráfego de rede, combinando atividades benignas com múltiplos tipos de ataques DDoS, permitindo uma avaliação mais precisa de modelos de detecção de intrusão \cite{sharafaldin2019}.

O CICDDoS2019 disponibiliza os dados nos formatos PCAP (\textit{Packet Capture}, ou captura de pacotes), que armazena o tráfego bruto da rede, e CSV, contendo 80 atributos de tráfego de rede extraídos dos fluxos de comunicação pela ferramenta CICFlowMeter \cite{sharafaldin2019}. Os arquivos PCAP registram os pacotes transmitidos durante a comunicação em rede, permitindo analisar as informações do tráfego original. A ferramenta CICFlowMeter processa essas capturas e extrai características dos fluxos de comunicação, como duração, quantidade e tamanho dos pacotes, gerando arquivos CSV estruturados. Esse formato facilita a utilização dos dados em análises estatísticas e no treinamento de modelos de aprendizado de máquina \cite{sharafaldin2019}.

O conjunto de dados foi desenvolvido em um ambiente experimental controlado (\textit{testbed}), composto por duas redes distintas — rede de ataque e rede de vítima —, integrando diferentes sistemas operacionais (Windows e Ubuntu), servidores, roteadores, \textit{switches} e dispositivos de segurança, como \textit{firewalls}, com o objetivo de simular cenários reais de redes corporativas \cite{sharafaldin2019}. Dessa forma, o CICDDoS2019 constitui uma base relevante para o treinamento e a validação de modelos de detecção de intrusão, permitindo a identificação de padrões associados a ataques DDoS e contribuindo para o desenvolvimento de soluções mais eficazes para a segurança de redes \cite{sharafaldin2019}.

Considerando os conjuntos de treinamento e teste, o CICDDoS2019 contempla 14 classes de tráfego: \texttt{Benign}, \texttt{DrDoS-DNS}, \texttt{DrDoS-LDAP}, \texttt{DrDoS-MSSQL}, \texttt{DrDoS-NetBIOS}, \texttt{DrDoS-NTP}, \texttt{DrDoS-SNMP}, \texttt{DrDoS-SSDP}, \texttt{DrDoS-TFTP}, \texttt{SYN}, \texttt{UDP}, \texttt{UDP-Lag}, \texttt{WebDDoS} e \texttt{PortScan} \cite{sharafaldin2019}. A classe \texttt{Benign} representa o tráfego legítimo de rede, gerado a partir de perfis de interação humana que simulam atividades realizadas por usuários por meio de protocolos como HTTP, HTTPS, FTP, SSH e e-mail. As demais classes representam diferentes comportamentos associados a ataques DDoS e atividades maliciosas. A classe \texttt{PortScan}, especificamente, foi executada apenas durante a captura do conjunto de teste, sendo utilizada como uma classe desconhecida, ausente do treinamento dos modelos \cite{sharafaldin2019}.

A separação entre os conjuntos de treinamento e teste foi realizada pelos próprios autores do CICDDoS2019 com base em capturas efetuadas em dois dias distintos. O conjunto de treinamento foi coletado em 12 de janeiro de 2019, entre 10h30 e 17h15, enquanto o conjunto de teste foi capturado em 11 de março de 2019, entre 9h40 e 17h35 \cite{sharafaldin2019}. Dessa maneira, o conjunto de dados já disponibiliza uma divisão experimental e temporal, na qual os registros de janeiro são destinados ao treinamento e os registros de março à avaliação dos modelos.

Essa separação foi adotada para permitir a avaliação dos algoritmos em dados coletados em períodos diferentes, contribuindo para analisar sua capacidade de generalização diante de condições distintas de tráfego e captura. Além disso, a presença exclusiva da classe \texttt{PortScan} no conjunto de teste possibilita investigar o comportamento dos classificadores diante de uma ameaça não apresentada durante o treinamento, aproximando a avaliação de situações em que sistemas de detecção de intrusão precisam lidar com padrões de ataques desconhecidos \cite{sharafaldin2019}.

A Figura~4 apresenta a distribuição percentual dos diferentes tipos de tráfego presentes no conjunto de dados CICDDoS2019. Observa-se que o conjunto é composto por tráfego legítimo e múltiplas classes de ataques DDoS, refletindo diferentes padrões de comunicação em rede \cite{sharafaldin2019}. Essa diversidade contribui para a avaliação dos modelos de aprendizado de máquina, permitindo analisar sua capacidade de distinguir atividades normais de comportamentos maliciosos associados a ataques distribuídos de negação de serviço \cite{sharafaldin2019}.

\begin{figure}
    \centering
    \caption{Distribuição de classes do dataset}
    \label{fig:distribuicao-dataset}
    \includegraphics[width=0.6\textwidth]{assets/distribuicao_classes_dataset_completo.png}
    \legend{Fonte: Elaborado pelo autor}
\end{figure}

\paragraph{Desbalanceamento de classes e \textit{undersampling}}

O desbalanceamento de classes ocorre quando as categorias presentes em um conjunto de dados possuem quantidades diferentes de amostras, fazendo com que determinadas classes sejam mais representadas que outras. Essa característica pode influenciar o processo de treinamento de modelos de aprendizado de máquina e, consequentemente, seu desempenho preditivo. Para lidar com esse cenário, podem ser empregadas técnicas de reamostragem que alteram a distribuição das classes utilizadas durante o treinamento, entre elas o \textit{oversampling} e o \textit{undersampling} \cite{yang2024}.

O \textit{random undersampling} consiste na redução aleatória de amostras da classe majoritária, buscando diminuir a diferença de proporção entre as classes presentes no conjunto utilizado para treinamento. Diferentemente do \textit{oversampling}, que aumenta a quantidade de observações da classe minoritária, o \textit{undersampling} modifica o conjunto de treinamento por meio da remoção de observações da classe predominante. Entretanto, a aplicação dessas técnicas deve ser avaliada de acordo com as características dos dados e do problema analisado, uma vez que o balanceamento das classes não implica necessariamente melhora no desempenho preditivo do modelo. Nos experimentos realizados por Yang et al. \cite{yang2024}, tanto o \textit{random undersampling} quanto o \textit{random oversampling} não produziram, de maneira geral, aumento do desempenho discriminativo dos modelos avaliados.

\subsection{IA (Inteligência Artificial)}

Nos fundamentos da obra de Russell e Norvig \cite{russell2010}, a Inteligência Artificial (IA) pode ser compreendida como o campo da ciência da computação dedicado ao desenvolvimento de sistemas capazes de simular capacidades cognitivas humanas, como raciocínio, aprendizagem, percepção, tomada de decisão e resolução de problemas. Sob essa perspectiva, a IA envolve a criação de agentes inteligentes capazes de perceber o ambiente em que estão inseridos e agir de maneira racional para atingir objetivos específicos. Entre suas principais aplicações estão a automação de processos, o reconhecimento de padrões, o processamento de linguagem natural e a análise preditiva, sendo seu uso cada vez mais relevante em áreas como segurança da informação e detecção de intrusões \cite{russell2010}. Nesse contexto, técnicas de aprendizado de máquina (\textit{Machine Learning}), um dos principais ramos da IA, permitem que sistemas identifiquem comportamentos anômalos em grandes volumes de dados, contribuindo para a detecção mais rápida e precisa de ameaças cibernéticas \cite{russell2010}.

\subsubsection{\textit{Machine Learning}}

Como uma das principais áreas derivadas da Inteligência Artificial (IA), o \textit{Machine Learning} (ML) surge como uma abordagem voltada ao desenvolvimento de sistemas capazes de aprender a partir de dados e aprimorar seu desempenho de forma progressiva, sem depender exclusivamente de regras previamente programadas \cite{russell2010,mitchell1997}. De acordo com Russell e Norvig \cite{russell2010}, a capacidade de adaptação e aprendizado constitui um dos elementos centrais para a construção de agentes inteligentes, tornando o aprendizado automatizado um componente fundamental da evolução da IA. Nesse contexto, o aprendizado de máquina permite que modelos computacionais identifiquem padrões, realizem previsões e apoiem processos de tomada de decisão em cenários complexos.

Segundo Mitchell \cite{mitchell1997}, o aprendizado de máquina baseia-se no desenvolvimento de algoritmos capazes de melhorar seu desempenho a partir da experiência obtida por meio dos dados. Dessa forma, o ML utiliza métodos estatísticos e computacionais para extrair conhecimento de grandes volumes de informações, sendo amplamente aplicado em problemas que exigem análise contínua, identificação de padrões e apoio à tomada de decisão. Além disso, essa área abrange diferentes tipos de aprendizado, que variam de acordo com a natureza dos dados e com os objetivos específicos de cada aplicação, sendo seu entendimento essencial para a utilização dessas técnicas em sistemas de detecção de intrusão e resposta a ameaças cibernéticas \cite{mitchell1997,russell2010}.

\paragraph{Algoritmos de Aprendizado de Máquina}

Segundo Russell e Norvig \cite{russell2010}, o processo de aprendizado ocorre a partir dos diferentes tipos de \textit{feedback} recebidos pelo agente, os quais determinam as principais categorias de aprendizado de máquina. De modo geral, o \textit{Machine Learning} utiliza conjuntos de dados para treinar modelos computacionais capazes de identificar padrões, estabelecer relações e realizar previsões a partir das informações analisadas. Conforme apresentado por Gopalakrishnan, Tuptuk e Hailes \cite{gopalakrishnan2023}, esse processo envolve duas etapas principais: o treinamento e o teste do modelo. Na primeira etapa, os dados são submetidos a procedimentos de pré-processamento e extração de características, permitindo que o algoritmo identifique padrões relevantes e desenvolva seu aprendizado. Na etapa seguinte, o modelo treinado é aplicado a dados que não foram utilizados durante o treinamento, possibilitando avaliar sua capacidade de generalização e seu desempenho na realização de classificações \cite{gopalakrishnan2023}.

De acordo com Gopalakrishnan, Tuptuk e Hailes \cite{gopalakrishnan2023}, os algoritmos de aprendizado de máquina são tradicionalmente divididos em três categorias principais: aprendizado supervisionado, aprendizado não supervisionado e aprendizado por reforço. Essa classificação considera a maneira como os dados e as informações de \textit{feedback} são utilizados durante o processo de aprendizagem. Cada abordagem apresenta características específicas e é empregada conforme a natureza do problema a ser solucionado. Nesse contexto, destaca-se o aprendizado supervisionado, utilizado neste trabalho por permitir o treinamento de modelos a partir de dados previamente rotulados, possibilitando a classificação do tráfego de rede entre atividades normais e maliciosas \cite{russell2010,gopalakrishnan2023}.

\paragraph{\textit{Machine Learning} aplicado à detecção de intrusão}

A aplicação de \textit{Machine Learning} (ML) em sistemas de detecção de intrusão tem se consolidado como uma abordagem estratégica para aprimorar a identificação de ameaças em ambientes computacionais, especialmente diante do aumento da complexidade e do volume do tráfego de rede \cite{gopalakrishnan2023}. Diferentemente dos métodos tradicionais baseados em assinaturas, os modelos de aprendizado de máquina permitem analisar grandes volumes de dados e identificar padrões associados a comportamentos normais e anômalos, possibilitando a detecção de ataques previamente desconhecidos e maior capacidade de adaptação a novas ameaças \cite{gopalakrishnan2023}. Nesse contexto, o ML é empregado na extração e análise de características do tráfego de rede, permitindo a classificação automática de eventos e contribuindo para a melhoria da precisão dos sistemas de detecção de intrusão \cite{gopalakrishnan2023}. Entre suas aplicações específicas, destaca-se a detecção de ataques de \textit{Distributed Denial of Service} (DDoS). Dessa forma, a aplicação de ML à detecção de intrusão representa um avanço importante para a segurança cibernética, contribuindo para o desenvolvimento de sistemas mais inteligentes, resilientes e capazes de responder de forma mais eficaz às ameaças emergentes \cite{gopalakrishnan2023,aktar2023}.

\subsubsection{Aprendizado supervisionado}

O aprendizado supervisionado é um dos paradigmas mais tradicionais e amplamente utilizados no campo do aprendizado de máquina, sendo aplicado principalmente em problemas de classificação, regressão e ranqueamento. Nesse modelo, o algoritmo recebe um conjunto de dados previamente rotulados, denominado conjunto de treinamento, e utiliza essas informações para aprender padrões capazes de realizar previsões sobre dados ainda não observados \cite{mohri2018}. Segundo Mohri, Rostamizadeh e Talwalkar \cite{mohri2018}, o aprendizado supervisionado consiste em fornecer ao aprendiz exemplos rotulados para que ele consiga fazer previsões em pontos desconhecidos, sendo esse o cenário mais comum em aplicações de aprendizado de máquina.

De acordo com Luxburg e Schölkopf \cite{luxburg2008}, o problema mais estudado dentro do aprendizado de máquina é a classificação supervisionada. Nesse cenário, existem dois conjuntos fundamentais: o espaço de entrada, representado por $\mathcal{X}$, que contém as instâncias ou objetos, e o espaço de saída, representado por $\mathcal{Y}$, correspondente aos rótulos ou categorias possíveis. O processo de aprendizagem ocorre a partir de pares ordenados compostos por uma entrada e sua respectiva saída correta, $(x_1,y_1),\ldots,(x_n,y_n)$, permitindo ao algoritmo identificar uma relação funcional entre os dados de entrada e os rótulos desejados \cite{luxburg2008}.

O objetivo do aprendizado supervisionado é encontrar uma função de mapeamento $f:\mathcal{X}\rightarrow\mathcal{Y}$. Essa função, denominada classificador, deve atribuir corretamente os rótulos às instâncias analisadas. Para isso, o algoritmo recebe um conjunto de treinamento formado por exemplos rotulados e, a partir deles, produz um modelo capaz de inferir resultados para novos dados \cite{luxburg2008}. Segundo Luxburg e Schölkopf \cite{luxburg2008}, o aprendizado supervisionado recebe esse nome porque o algoritmo possui acesso às respostas corretas durante o treinamento. Dessa forma, o sistema consegue avaliar continuamente se suas previsões estão corretas, caracterizando um processo de aprendizagem guiado ou supervisionado.

\paragraph{Classificação binária}

A classificação binária ocorre quando o espaço de saída ($\mathcal{Y}$) contém exatamente duas categorias possíveis. Por convenção, os rótulos podem ser representados por $\{-1,+1\}$ ou $\{0,1\}$. Nesse cenário, o objetivo do algoritmo é aprender uma função de classificação $f:\mathcal{X}\rightarrow\mathcal{Y}$, capaz de atribuir cada instância do conjunto de entrada à categoria correspondente, buscando minimizar os erros de classificação \cite{vonluxburg2011}.

Segundo Mohri, Rostamizadeh e Talwalkar \cite{mohri2018}, um exemplo comum é a identificação de mensagens de \textit{spam}, em que cada mensagem é classificada como legítima ou indesejada. No contexto da segurança de redes, a classificação binária pode ser utilizada para distinguir tráfego benigno de tráfego malicioso, identificando se uma comunicação representa uma atividade normal ou um possível ataque.

\paragraph{Classificação multiclasse}

A classificação multiclasse é utilizada quando o problema envolve mais de duas categorias possíveis. Nesse caso, cada instância pertence a uma única classe dentro de um conjunto de $k$ categorias, representado por $\mathcal{Y}=\{1,\ldots,k\}$. Diferentemente da classificação binária, o algoritmo precisa identificar qual das diversas classes representa melhor as características da instância analisada \cite{mohri2018}.

Para realizar essa classificação, podem ser utilizadas funções de pontuação (\textit{scoring functions}), que atribuem valores às possíveis classes de uma instância. O modelo seleciona como resultado a categoria que apresenta a maior pontuação, conforme a expressão
\[
f(x)=\arg\max_{y\in\mathcal{Y}} h(x,y),
\]
em que $h(x,y)$ representa a pontuação atribuída à instância $x$ para a classe $y$ \cite{mohri2018}.

Como exemplo, um sistema de reconhecimento de imagens pode classificar objetos em categorias como carro, bicicleta, ônibus ou motocicleta \cite{mohri2018}. No contexto deste trabalho, a classificação multiclasse permite identificar as diferentes categorias de tráfego presentes no CICDDoS2019, distinguindo o tráfego benigno dos diversos tipos de ataques DDoS.

\subsection{Random Forest}

\paragraph{Decision Trees}

As árvores de decisão (\textit{Decision Trees} -- DT) são modelos de aprendizado supervisionado amplamente utilizados em tarefas de classificação e regressão devido à sua simplicidade, interpretabilidade e eficiência computacional. Segundo Breiman et al. \cite{breiman1984}, as árvores de decisão possuem uma estrutura hierárquica composta por nós, ramos e folhas. Os nós representam pontos de decisão nos quais determinados atributos dos dados são avaliados; os ramos representam os possíveis caminhos resultantes dessas decisões; e as folhas correspondem às classes ou aos valores finais previstos pelo modelo.

O funcionamento de uma árvore de decisão ocorre por meio de divisões sucessivas do conjunto de dados. Durante o treinamento, o algoritmo seleciona os atributos mais relevantes para separar os dados em subconjuntos cada vez mais homogêneos. Esse processo é realizado de forma recursiva até que critérios de parada sejam atingidos, como profundidade máxima da árvore ou quantidade mínima de instâncias em um nó. De acordo com Breiman et al. \cite{breiman1984}, o objetivo dessas divisões é reduzir a impureza dos grupos formados, permitindo que o modelo identifique padrões presentes nos dados de treinamento.

Os nós da árvore representam pontos em que uma decisão é tomada com base em uma característica específica do conjunto de dados. Dependendo do valor do atributo analisado, o fluxo segue por diferentes ramos até alcançar uma folha, responsável pela previsão final do modelo. Essa estrutura torna as árvores de decisão facilmente interpretáveis, permitindo compreender de maneira clara como as previsões são realizadas.

Apesar das vantagens relacionadas à simplicidade e à interpretabilidade, árvores de decisão individuais podem apresentar elevada sensibilidade aos dados de treinamento. Isso ocorre porque o modelo pode aprender excessivamente os detalhes e ruídos presentes no conjunto de treinamento, fenômeno conhecido como \textit{overfitting}. O \textit{overfitting} reduz a capacidade de generalização do algoritmo, fazendo com que o modelo apresente bom desempenho nos dados utilizados no treinamento, mas baixa precisão ao lidar com novos exemplos \cite{breiman1984}.

\paragraph{Random Forest}

O algoritmo Random Forest é um dos métodos de aprendizado supervisionado mais utilizados na área de \textit{Machine Learning}, especialmente em problemas de classificação e regressão. Proposto por Leo Breiman em 2001, o modelo pertence à categoria de métodos conhecidos como \textit{ensemble learning}, cuja principal característica é a combinação de múltiplos modelos para produzir resultados mais precisos, robustos e com maior capacidade de generalização \cite{breiman2001}.

O Random Forest utiliza múltiplas árvores de decisão independentes como base para seu funcionamento. Em vez de utilizar apenas uma única árvore para realizar previsões, o algoritmo cria uma ``floresta'' composta por diversas árvores treinadas separadamente. Essa abordagem reduz os problemas de variância e \textit{overfitting} encontrados em árvores individuais, aumentando a robustez e a estabilidade do modelo \cite{breiman2001}.

Durante o treinamento, cada árvore é construída a partir de subconjuntos aleatórios do conjunto de dados original, utilizando a técnica denominada \textit{bootstrap sampling}. Essa técnica consiste em gerar várias amostras aleatórias com reposição a partir do conjunto de treinamento. Dessa forma, algumas instâncias podem ser repetidas, enquanto outras podem não ser selecionadas em determinada amostra. Como consequência, cada árvore é treinada com dados ligeiramente diferentes, aumentando a diversidade entre os modelos da floresta.

Além disso, durante o processo de construção das árvores, o algoritmo considera apenas subconjuntos aleatórios de atributos para realizar as divisões dos nós. Segundo Breiman \cite{breiman2001}, essa aleatoriedade reduz a correlação entre as árvores individuais, contribuindo para melhorar a capacidade de generalização do modelo e reduzir erros de previsão.

No processo de classificação, cada árvore da floresta realiza uma previsão individual para a instância analisada. O resultado final é definido por meio de votação majoritária, na qual a classe mais frequentemente prevista entre as árvores é escolhida como resposta final. Em problemas de regressão, o resultado geralmente corresponde à média dos valores previstos pelas árvores individuais. Essa estratégia reduz significativamente os efeitos do \textit{overfitting} e melhora a estabilidade das previsões realizadas pelo modelo \cite{breiman2001}.

Outra característica importante do Random Forest é sua capacidade de lidar com grandes volumes de dados e conjuntos com alta dimensionalidade. O algoritmo também possibilita medir a importância das variáveis utilizadas durante o treinamento, permitindo identificar quais atributos exercem maior influência nas previsões realizadas. Além disso, apresenta elevada tolerância a ruídos, dados ausentes e variações nos conjuntos de treinamento, fatores que contribuíram para sua ampla utilização em áreas como análise preditiva, reconhecimento de padrões, detecção de fraudes, mineração de dados e sistemas inteligentes \cite{breiman2001}.

Dessa forma, o Random Forest pode ser compreendido como uma evolução das árvores de decisão tradicionais, utilizando a combinação de múltiplas árvores independentes para reduzir erros de generalização e aumentar a precisão das previsões. A utilização conjunta de diferentes árvores permite que o modelo apresente desempenho superior ao de árvores isoladas, consolidando o Random Forest como um dos algoritmos mais eficientes e populares do aprendizado de máquina contemporâneo \cite{breiman2001}.

\paragraph{Hiperparâmetros}

O desempenho do algoritmo Random Forest é influenciado por diferentes hiperparâmetros, responsáveis por controlar a construção das árvores de decisão e a complexidade do modelo \cite{breiman2001}. Entre os aspectos que se destacam estão a quantidade de árvores que compõem a floresta, a profundidade máxima que cada árvore pode alcançar e a quantidade mínima de amostras necessárias para formar uma folha.

A quantidade de árvores determina o número de modelos individuais utilizados na composição da floresta, cujas previsões são combinadas para produzir a classificação final. O aumento desse número pode contribuir para maior estabilidade das previsões e reduzir a influência de variações individuais entre as árvores, embora também eleve o custo computacional e o tempo de processamento. Entretanto, a partir de determinada quantidade, o acréscimo de novas árvores pode proporcionar ganhos cada vez menores no desempenho do classificador \cite{becerra2025}.

A profundidade máxima estabelece o limite de crescimento de cada árvore, determinando quantas divisões podem ser realizadas desde o nó raiz até os nós terminais. Árvores mais profundas conseguem representar relações mais complexas presentes nos dados de treinamento, porém podem aumentar a complexidade do modelo e favorecer o \textit{overfitting}. Em contrapartida, limitar excessivamente a profundidade pode restringir a capacidade de aprendizado do classificador, dificultando a identificação de padrões relevantes \cite{becerra2025}.

Por fim, a quantidade mínima de amostras necessárias em cada folha estabelece o número mínimo de registros que um nó terminal deve conter para que uma divisão seja considerada válida. Valores menores permitem a formação de folhas mais específicas, aumentando a capacidade de adaptação às características dos dados de treinamento. Já valores maiores restringem a criação de folhas muito pequenas, reduzindo a complexidade das árvores e podendo contribuir para evitar o sobreajuste. Entretanto, quando esse limite é excessivamente elevado, o modelo pode perder capacidade de identificar padrões importantes nos dados \cite{becerra2025}.

\subsection{Support Vector Machine (SVM)}

O algoritmo Support Vector Machine (SVM) é um método de aprendizado supervisionado amplamente utilizado em problemas de classificação e regressão, sendo reconhecido por sua elevada capacidade de generalização e eficiência em conjuntos de dados de alta dimensionalidade. O modelo foi formalizado por Corinna Cortes e Vladimir Vapnik em 1995, fundamentando-se nos princípios da Teoria do Aprendizado Estatístico (\textit{Statistical Learning Theory}), desenvolvida por Vapnik \cite{vapnik1998}. O principal objetivo do SVM é encontrar uma fronteira de decisão capaz de separar diferentes classes de dados da forma mais otimizada possível \cite{cortespvapnik1995}.

O funcionamento do SVM baseia-se na construção de um hiperplano de separação entre classes. Em problemas bidimensionais, esse hiperplano pode ser compreendido como uma linha que divide os dados em diferentes categorias; em espaços multidimensionais, corresponde a uma superfície de separação. Segundo Cortes e Vapnik \cite{cortespvapnik1995}, o algoritmo busca identificar o hiperplano que maximize a distância entre as classes, denominada margem (\textit{margin}). Quanto maior essa margem, maior tende a ser a capacidade de generalização do modelo para novos dados.

Os elementos mais importantes para a definição do hiperplano são chamados de vetores de suporte (\textit{support vectors}). Esses vetores correspondem às instâncias de treinamento localizadas mais próximas da fronteira de separação. Dessa forma, são os pontos que efetivamente influenciam a posição e a orientação do hiperplano. De acordo com Vapnik \cite{vapnik1998}, a utilização apenas dos vetores de suporte torna o modelo eficiente e reduz a complexidade do processo de classificação.

O SVM apresenta eficiência no trabalho com conjuntos de dados que possuem muitos atributos e pode ser aplicado a diferentes tarefas de classificação e reconhecimento de padrões. Embora seja tradicionalmente utilizado em classificações binárias, também pode ser adaptado para problemas multiclasse. Sua capacidade de maximizar a margem de separação e lidar com dados não lineares por meio de \textit{kernels} contribui para sua capacidade de generalização e reduz o risco de \textit{overfitting} \cite{cortespvapnik1995,vapnik1998}.

\paragraph{Hiperparâmetros}

Em situações nas quais os dados não podem ser separados linearmente, o SVM utiliza funções conhecidas como \textit{kernels}. Essas funções permitem transformar os dados para espaços de maior dimensionalidade, facilitando a separação entre as classes. Entre os principais tipos de \textit{kernel} estão o linear, o polinomial e o radial (\textit{Radial Basis Function} -- RBF). Conforme descrito por Cortes e Vapnik \cite{cortespvapnik1995}, essa abordagem permite ao SVM trabalhar com problemas mais complexos e dados que não apresentam separação linear.

Além da escolha do \textit{kernel}, o desempenho do SVM também depende da configuração de seus hiperparâmetros, que definem como o modelo realiza a classificação. Entre os principais, destaca-se o parâmetro $C$, responsável por controlar a penalização dos erros de classificação. Valores elevados de $C$ fazem com que o modelo priorize a redução desses erros, enquanto valores menores permitem maior tolerância às classificações incorretas, favorecendo uma margem de separação mais ampla \cite{cortespvapnik1995}.

Outro hiperparâmetro importante é o gama ($\gamma$), utilizado principalmente com o \textit{kernel} RBF. Ele determina o quanto cada ponto de dado influencia a fronteira de decisão do modelo. Valores elevados tornam essa influência mais restrita, enquanto valores menores permitem uma influência mais ampla. A combinação desses parâmetros interfere diretamente no desempenho do SVM e em sua capacidade de identificar padrões sem se ajustar excessivamente aos dados de treinamento \cite{vapnik1998}.

A Teoria do Aprendizado Estatístico, proposta por Vapnik \cite{vapnik1998}, fornece a base teórica para o funcionamento do SVM. Segundo essa teoria, o objetivo de um modelo de aprendizado não é apenas minimizar os erros durante o treinamento, mas também apresentar bons resultados diante de novos dados. Nesse contexto, o SVM busca equilibrar a complexidade do modelo e seu desempenho, reduzindo problemas de \textit{overfitting} \cite{vapnik1998}. O SVM reduz esse problema por meio da maximização da margem entre as classes, estratégia que favorece modelos mais robustos e menos sensíveis às variações dos dados \cite{vapnik1998}.

\subsection{Inteligência Artificial Explicável (XAI)}

A \textbf{Inteligência Artificial Explicável} (\textit{Explainable Artificial Intelligence} -- XAI) compreende métodos e técnicas que buscam tornar mais transparentes e compreensíveis as decisões realizadas por sistemas de Inteligência Artificial. Sua aplicação é especialmente relevante em modelos considerados de \textbf{caixa-preta}, nos quais a complexidade interna dificulta a compreensão dos fatores que levaram a determinada previsão. Nesse contexto, a XAI permite interpretar a influência das características de entrada sobre os resultados produzidos pelo modelo \cite{minh2022}.

As técnicas de XAI podem fornecer explicações sobre o comportamento geral do modelo ou sobre previsões individuais, auxiliando na compreensão dos critérios utilizados pelos algoritmos. Dessa forma, além da análise de métricas de desempenho, torna-se possível investigar quais características apresentam maior influência nas decisões realizadas pelo modelo, aumentando a transparência e a interpretabilidade dos resultados \cite{minh2022}.

\paragraph{SHAP}

Entre os métodos utilizados em XAI, destaca-se o \textbf{SHAP} (\textit{SHapley Additive exPlanations}), apresentado por Scott Lundberg e Su-In Lee em 2017. Esse método atribui a cada característica de entrada um valor que representa sua contribuição para determinada previsão realizada pelo modelo \cite{lundberg2017}.

O SHAP é fundamentado nos \textbf{valores de Shapley}, provenientes da teoria dos jogos cooperativos. A partir dessa abordagem, é possível estimar a contribuição de cada característica para o resultado do modelo e identificar quais variáveis exerceram maior influência sobre determinada classificação. Além de permitir a interpretação de previsões individuais, os valores SHAP podem ser analisados em conjunto para compreender o comportamento geral do modelo \cite{lundberg2017}.

\subsection{Métricas de Avaliação de Modelos}
\subsection{Acurácia}
A acurácia é uma das métricas mais utilizadas para avaliar modelos de classificação em aprendizado de máquina. Essa métrica representa a proporção total de previsões corretas realizadas pelo modelo em relação ao número total de instâncias analisadas. Segundo \citeonline{Rainio2024}, a acurácia expressa o percentual de classificações corretas considerando tanto os exemplos positivos quanto os negativos.
Em problemas de classificação binária, cada previsão realizada pelo modelo pode ser categorizada em quatro tipos principais: verdadeiro positivo (True Positive – TP), verdadeiro negativo (True Negative – TN), falso positivo (False Positive – FP) e falso negativo (False Negative – FN). Um verdadeiro positivo ocorre quando uma instância positiva é corretamente classificada como positiva; um verdadeiro negativo corresponde a uma instância negativa corretamente classificada como negativa; um falso positivo ocorre quando uma instância negativa é incorretamente classificada como positiva; e um falso negativo ocorre quando uma instância positiva é incorretamente classificada como negativa \cite{Rainio2024}.
A fórmula da acurácia é dada por:

\begin{equation}
Accuracy = \frac{TP + TN}{TP + TN + FP + FN}
\end{equation}


Onde:
TP (True Positive): quantidade de verdadeiros positivos;
TN (True Negative): quantidade de verdadeiros negativos;
FP (False Positive): quantidade de falsos positivos;
FN (False Negative): quantidade de falsos negativos.
O valor da acurácia pertence ao intervalo [0,1], no qual valores mais próximos de 1 indicam melhor desempenho do modelo. Entretanto, conforme destacado por \citeonline{Rainio2024}, a acurácia pode não ser suficiente para avaliar modelos em conjuntos de dados desbalanceados, isto é, quando existe grande diferença entre a quantidade de exemplos positivos e negativos. Nesses casos, métricas como recall tornam-se mais adequadas para analisar o comportamento do modelo.

\subsection{Recall}
O recall (revocação) ou taxa de verdadeiros positivos, é uma métrica utilizada para medir a capacidade do modelo de identificar corretamente as instâncias positivas. Segundo \citeonline{Rainio2024}, o recall representa a proporção de exemplos positivos que foram corretamente classificados pelo algoritmo.
Essa métrica é especialmente importante em aplicações nas quais identificar corretamente os casos positivos é essencial, como diagnósticos médicos, detecção de fraudes e sistemas de segurança. Em tais cenários, falsos negativos podem causar impactos significativos, tornando necessário maximizar a taxa de identificação correta dos exemplos positivos.
A fórmula do recall é dada por:

\begin{equation}
Recall = \frac{TP}{TP + FN}
\end{equation}

O valor do recall também pertence ao intervalo [0,1], sendo que valores próximos de 1 indicam que o modelo consegue identificar corretamente a maior parte das instâncias positivas presentes no conjunto de dados. Conforme discutido por \citeonline{Rainio2024}, métricas como recall frequentemente fornecem informações mais relevantes do que a acurácia em problemas com classes desbalanceadas, pois avaliam separadamente o desempenho do modelo em relação às classes positivas e negativas.

\subsection{Precisão}
A precisão é uma métrica que indica a proporção de instâncias classificadas como positivas que foram corretamente identificadas pelo modelo, sendo calculada pela razão entre os verdadeiros positivos (TP) e o total de predições positivas (TP + FP). Embora o par precisão/revocação forneça insights úteis sobre o comportamento de um classificador, \citeonline{Chicco2020} ressaltam que a precisão e o F1-score (sua média harmônica com a revocação) podem apresentar resultados "perigosamente otimistas" e inflados, especialmente em conjuntos de dados desbalanceados. Isso ocorre porque essas métricas são independentes do número de amostras corretamente classificadas como negativas (TN). Por essa razão, os autores argumentam que, para uma avaliação mais informativa e fiel em tarefas de classificação binária, deve-se preferir o Coeficiente de Correlação de Matthews (MCC), que integra as quatro categorias da matriz de confusão e produz um escore alto apenas se o modelo obtiver bons resultados tanto na classe positiva quanto na negativa. 
A fórmula da precisão é dada por:

\begin{equation}
\text{Precisão} = \frac{TP}{TP + FP}
\end{equation}

\subsection{F1-Score}
O F1-score é uma métrica amplamente utilizada para avaliar o desempenho de modelos de classificação, especialmente em cenários com desequilíbrio entre as classes \cite{Hand2021}. Essa métrica corresponde à média harmônica entre a precisão e a revocação, permitindo combinar dois aspectos distintos e complementares do desempenho do modelo \cite{Hand2021}. Diferentemente de medidas que podem ser infladas por grandes volumes de casos negativos, o F1-score foca nos casos positivos de interesse, sendo particularmente relevante em domínios onde os verdadeiros negativos não são a prioridade, como na detecção de fraudes e em problemas de classificação em disciplinas computacionais \cite{Hand2021}. 
Essa métrica corresponde à média harmônica entre a precisão e a revocação, sendo calculada pela fórmula: F1 = 2 × (Precisão × Revocação) / (Precisão + Revocação).
A fórmula do f1-score é dada por:

\begin{equation}
F1 = 2 \cdot \frac{\text{Precisão} \cdot \text{Revocação}}{\text{Precisão} + \text{Revocação}}
\end{equation}

A precisão representa a proporção de instâncias classificadas como positivas que estão corretas, enquanto a revocação indica a capacidade do modelo de identificar corretamente todos os casos positivos presentes no conjunto de dados.

\subsection{Coeficiente de correlação de Matthews (MCC)}
é uma métrica utilizada para avaliar o desempenho de classificadores binários, sendo especialmente adequada para conjuntos de dados desbalanceados. Diferentemente da acurácia, o MCC considera simultaneamente os valores de verdadeiros positivos, verdadeiros negativos, falsos positivos e falsos negativos, fornecendo uma avaliação mais equilibrada do modelo. Seus valores variam entre -1 e 1, em que 1 representa uma classificação perfeita, 0 indica desempenho equivalente ao acaso e valores negativos refletem predições inconsistentes em relação aos valores reais. Por considerar todos os elementos da matriz de confusão, o MCC é frequentemente recomendado como uma métrica robusta para a avaliação de modelos de aprendizado de máquina em problemas de classificação binária \cite{Chicco2023}. 
A fórmula é do MCC dada por:

\begin{equation}
MCC = \frac{(TP \times TN) - (FP \times FN)}
{\sqrt{(TP + FP)(TP + FN)(TN + FP)(TN + FN)}}
\label{eq:mcc}
\end{equation}

% ---------------------------------------------------------------
\section{Trabalhos correlatos}
\label{sec:trabalhos_correlatos}

\subsection{Trabalho 1: An Analysis of Intrusion Detection Classification using Supervised Machine Learning Algorithms on NSL-KDD Dataset}

Rastogi et al. \cite{rastogi2022} apresentam uma análise comparativa de algoritmos supervisionados de aprendizado de máquina aplicados à detecção de intrusões em redes. Utilizando o conjunto de dados NSL-KDD, os autores avaliaram os algoritmos Support Vector Machine (SVM), Random Forest (RF), K-Nearest Neighbors (KNN), Gaussian Naive Bayes (GNB), Decision Tree (DT) e Logistic Regression (LR), considerando a acurácia, o tempo de treinamento e o tempo de teste (inferência).

Nos experimentos de classificação binária, o algoritmo Random Forest obteve a maior acurácia (98,70\%), com tempo de treinamento de 3,51 segundos e tempo de teste de 0,13 segundos, enquanto o SVM alcançou 96,70\% de acurácia, demandando o maior tempo de treinamento do experimento (232,32 segundos) e 11,44 segundos de teste. No cenário de classificação multiclasse, o Random Forest manteve o melhor desempenho preditivo (98,47\% de acurácia), com 3,56 segundos de treinamento e 0,16 segundos de teste, ao passo que o SVM atingiu 95,25\% de acurácia, com 201,62 segundos de treinamento e 22,27 segundos de teste. O estudo evidencia a eficiência computacional do Random Forest em comparação a modelos com elevado tempo de teste, como o KNN (140,06 segundos no binário e 72,06 segundos no multiclasse), ou de menor acurácia, como o GNB (84,33\% no binário) e a DT (89,59\% no multiclasse).

Esse trabalho relaciona-se diretamente à presente pesquisa pela utilização de aprendizado supervisionado e, particularmente, pela avaliação dos algoritmos Random Forest e SVM. Entretanto, enquanto Rastogi et al. \cite{rastogi2022} realizam uma comparação mais ampla entre diferentes classificadores utilizando o NSL-KDD, a presente pesquisa concentra-se especificamente na detecção de ataques DDoS utilizando o CICDDoS2019 e aprofunda a comparação entre RF e SVM em cenários de classificação binária e multiclasse. Além do desempenho preditivo, são consideradas métricas como acurácia, precisão, revocação, F1-score e MCC, bem como o custo computacional, os tempos de treinamento e predição e o comportamento dos modelos antes e após a otimização de hiperparâmetros com validação cruzada. Dessa forma, busca-se avaliar não apenas qual algoritmo apresenta melhores resultados de classificação, mas também as vantagens e limitações de cada modelo diante das diferentes condições experimentais analisadas.

\subsection{Trabalho 2: Detecting DDoS Threats Using Supervised Machine Learning for Traffic Classification in Software Defined Networking}

Hirsi et al. \cite{hirsi2024} investigaram a detecção de ataques DDoS em ambientes de \textit{Software-Defined Networking} (SDN) utilizando algoritmos supervisionados de aprendizado de máquina para classificação binária do tráfego entre benigno e malicioso. Os autores desenvolveram um conjunto de dados próprio, composto por 1.048.575 registros, e avaliaram cinco algoritmos: Logistic Regression, Support Vector Machine (SVM), Random Forest (RF), K-Nearest Neighbors (KNN) e XGBoost. O estudo empregou seleção de características pelo método Chi-square, otimização de hiperparâmetros e validação cruzada com 5 \textit{folds}.

Nos experimentos realizados com o conjunto de dados proposto, o Random Forest apresentou o melhor desempenho, alcançando 98,97\% de acurácia, 98,33\% de precisão, 96,37\% de revocação e 97,34\% de F1-score, além da menor taxa de falsos alarmes (0,023). O SVM, por sua vez, obteve 96,52\% de acurácia, 95,75\% de precisão, 95,75\% de revocação e 95,75\% de F1-score. O desempenho do Random Forest também foi avaliado no conjunto CICDDoS2019, no qual alcançou 97,62\% de acurácia, 98,70\% de precisão, 97,73\% de revocação e 97,69\% de F1-score, reforçando a capacidade do algoritmo na identificação de tráfego DDoS em diferentes conjuntos de dados.

Embora apresente proximidade com a presente pesquisa pela utilização de RF e SVM, validação cruzada e otimização de hiperparâmetros, o estudo de Hirsi et al. \cite{hirsi2024} concentra-se na classificação binária em redes SDN. O presente trabalho amplia essa análise ao comparar diretamente RF e SVM em cenários de classificação binária e multiclasse utilizando o CICDDoS2019, incorporando ainda o MCC e a análise do custo computacional, incluindo os tempos de treinamento e predição.

\subsection{Trabalho 3: Distributed Denial-of-Service (DDoS) Attack Detection Using Supervised Machine Learning Algorithms}

Abiramasundari e Ramaswamy \cite{abiramasundari2025}, no estudo \textit{Distributed Denial-of-Service (DDoS) Attack Detection Using Supervised Machine Learning Algorithms}, propuseram o framework PCA-based Enhanced Distributed DDoS Attack Detection (EDAD) para a identificação de ataques DDoS por meio de aprendizado de máquina supervisionado. A abordagem contempla etapas de pré-processamento, normalização, tratamento do desbalanceamento entre classes e redução de dimensionalidade por \textit{Principal Component Analysis} (PCA), além da aplicação dos algoritmos Random Forest (RF), Support Vector Machine (SVM), K-Nearest Neighbors (KNN), Decision Tree (DT) e Logistic Regression em três conjuntos de dados: CICIDS2017, CICIDS2018 e CICDDoS2019.

No CICDDoS2019, diretamente relacionado ao conjunto utilizado nesta pesquisa, Random Forest e KNN apresentaram os melhores resultados, ambos com acurácia de 98,7\%. Para o Random Forest, foram reportados valores de 0,98 para precisão, revocação e F1-score, enquanto o SVM apresentou precisão de 0,95, revocação de 0,97, F1-score de 0,96 e acurácia de aproximadamente 96\%.

O estudo apresenta forte relação com a presente pesquisa por avaliar RF e SVM na detecção de ataques DDoS utilizando o CICDDoS2019. Entretanto, enquanto Abiramasundari e Ramaswamy \cite{abiramasundari2025} comparam cinco algoritmos e empregam PCA para redução de dimensionalidade, o presente trabalho concentra a análise comparativa em RF e SVM, investigando seu desempenho em cenários de classificação binária e multiclasse.

% ===============================================================
% CAPÍTULO 3 - DESENVOLVIMENTO
% ===============================================================
\chapter{Desenvolvimento / Implementação / Experimentos}
\label{cap:desenvolvimento}

% ---------------------------------------------------------------
% \section{Metodologia da pesquisa}
% \label{sec:metodologia_pesquisa}

% Descreva, de forma rigorosa e reprodutível, como a pesquisa foi conduzida. Informe o tipo de pesquisa (exploratória, empírica, formal etc.) e a hipótese ou o critério de sucesso adotado; detalhe os procedimentos da abordagem proposta, em ordem cronológica (use listas numeradas, diagramas de fluxo ou pseudocódigo para tornar o processo claro); liste as ferramentas, linguagens, bibliotecas, frameworks e ambientes utilizados, justificando cada escolha com base no problema; e explique como os experimentos foram configurados, quais variáveis foram controladas e quais métricas demonstram que a proposta cumpre os objetivos (ex.: acurácia, latência, custo).

% \begin{table}
%   \centering
%   \caption{Plano de experimentação}
%   \label{tab:plano_experimentos}
%   \begin{tabular}{lcc}
%     \toprule
%     \textbf{Experimento} & \textbf{Variável controlada} & \textbf{Métrica principal} \\
%     \midrule
%     Comparação com baseline & algoritmo & acurácia \\
%     Escalabilidade & carga & tempo médio \\
%     \bottomrule
%   \end{tabular}
%   \legend{Fonte: Elaborado pelo autor (2025)}
% \end{table}

% \textit{Propor algo não basta: demonstre como será validado.}

\subsection{Metodologia da pesquisa}
\label{sec:metodologia}

\textbf{Delineamento.} Esta pesquisa é aplicada, quantitativa e empírico-computacional, com caráter experimental e comparativo. É aplicada porque implementa e avalia modelos de aprendizado de máquina para apoiar a detecção de ataques DDoS. É quantitativa porque compara os modelos por métricas numéricas de desempenho e custo computacional. Seu caráter experimental decorre da variação sistemática dos algoritmos, das tarefas de classificação e das configurações de hiperparâmetros, mantendo definidos os conjuntos de dados e os procedimentos de avaliação. A comparação é realizada com dados históricos do CICDDoS2019, e não por meio de um ensaio em uma rede operacional.

\textbf{Procedimentos da pesquisa.} A condução dos experimentos seguiu, em linhas gerais, as etapas abaixo. Os detalhes da implementação são apresentados nas seções subsequentes.

\begin{enumerate}
\def\labelenumi{\arabic{enumi}.}
\item
  Os arquivos do conjunto CICDDoS2019 (18 tabelas, 70.427.637 registros) foram carregados no ClickHouse, preservando a separação temporal adotada: dados do primeiro dia para desenvolvimento e dados do segundo dia para teste.
\item
  Foi realizada a análise exploratória e, em seguida, o pré-processamento dos dados para gerar as entradas destinadas às tarefas de classificação binária e multiclasse.
\item
  Foram definidos os conjuntos de ajuste, validação e teste. O teste do segundo dia foi deduplicado em relação ao treino para reduzir a avaliação em fluxos com atributos já observados durante o aprendizado. No multiclasse, o tipo Portmap, presente apenas no teste, ficou fora das sete classes dos modelos.
\item
  Foram preparados os conjuntos de treino para as duas tarefas e estabelecidas as configurações de referência do Random Forest e do SVM.
\item
  Os hiperparâmetros foram selecionados por busca e validação cruzada. A validação foi usada para avaliar os modelos em dados não utilizados no ajuste e para as decisões previstas no procedimento experimental; a busca e a seleção das configurações não usaram o teste, que foi reservado à avaliação final. Como o teste foi observado mais de uma vez (por exemplo, nas repetições do RF com outras sementes), ele deixa de ser uma avaliação independente a partir da primeira observação, e isso é considerado na discussão dos resultados.
\item
  Os modelos padrão e ajustados foram avaliados nas tarefas binária e multiclasse. Foram calculadas métricas preditivas e medidas de custo computacional, incluindo tempos de treinamento e inferência, e a estabilidade do RF foi examinada com diferentes sementes.
\item
  Foram conduzidas análises complementares de explicabilidade com SHAP.
\item
  Por fim, os modelos salvos foram integrados a uma simulação interativa de IDS. Essa aplicação demonstra seu uso em fluxos selecionados, sem substituir a avaliação quantitativa realizada no teste completo.
\item
  Os resultados foram comparados, de forma qualitativa, com os de estudos da literatura que avaliaram RF e SVM na detecção de intrusões.
\end{enumerate}

\textbf{Diagrama da arquitetura experimental.} O diagrama UML abaixo representa o fluxo geral da pesquisa; a implementação de cada módulo é detalhada nas seções seguintes.

\begin{figure}
\centering
\caption{Diagrama UML da arquitetura experimental}
\label{fig:arquitetura-metodologia}
\includegraphics[width=\textwidth]{assets/arquitetura_metodologia}
\legend{Fonte: Elaborado pelo autor (2026).}
\end{figure}


\textbf{Ferramentas utilizadas.}

\begin{itemize}
\item
  \textbf{Python (3.13.14):} linguagem utilizada na implementação da solução. Sua disponibilidade de bibliotecas para aprendizado de máquina e análise de dados, bem como a integração entre diferentes etapas do desenvolvimento, tornou-a adequada à construção do IDS baseado em aprendizado supervisionado (PYTHON SOFTWARE FOUNDATION, 2026).
\item
  \textbf{Pandas (2.3.3) e NumPy (2.5.3):} empregadas no pré-processamento do CICDDoS2019. O Pandas foi utilizado na leitura, organização e manipulação dos dados; o NumPy, nas operações matemáticas e no processamento numérico (MUALFAH; ARDIANSYAH; GUNAWAN, 2025).
\item
  \textbf{Scikit-learn (1.9.1):} utilizado na implementação e avaliação do Random Forest e do SVM, incluindo validação cruzada, cálculo de métricas e ajuste de hiperparâmetros (SCIKIT-LEARN DEVELOPERS, 2026). Essa é também a versão exigida para carregar os modelos salvos.
\item
  \textbf{Matplotlib (3.11.2):} empregado para gerar gráficos de desempenho e representações visuais das métricas, facilitando a comparação dos modelos (HUNTER, 2007).
\item
  \textbf{Jupyter Notebook:} utilizado como ambiente de desenvolvimento para organizar e executar os procedimentos em Python e examinar os resultados obtidos (PROJECT JUPYTER, 2026).
\item
  \textbf{Joblib (1.6.0):} utilizado para persistir e carregar modelos preditivos e amostras. Essa persistência permitiu reutilizá-los nas etapas posteriores de avaliação, explicabilidade e simulação sem retreinamento (JOBLIB DEVELOPERS, 2026).
\item
  \textbf{SHAP (0.54.1):} utilizado para estimar a contribuição dos atributos às previsões dos modelos binários e multiclasse. A análise possibilitou examinar a importância das variáveis de tráfego nas decisões dos modelos (LUNDBERG; LEE, 2017).
\item
  \textbf{Streamlit e Plotly:} utilizados na simulação do IDS. O Streamlit organiza a interface e os controles do painel, enquanto o Plotly apresenta distribuições e métricas em gráficos interativos (AKASHKUMAR; THAMIZHARASAN, 2025). As versões instaladas não foram registradas; os requisitos da aplicação são Streamlit 1.40 ou superior e Plotly 5.20 ou superior.
\item
  \textbf{Psutil (7.2.2):} utilizado para monitorar recursos computacionais, incluindo uso de memória e informações sobre CPU e processos, apoiando a análise do custo computacional (SAI et al., 2024).
\item
  \textbf{ClickHouse:} empregado como banco de dados analítico, hospedado na Railway, para armazenar e consultar grandes volumes de dados. Sua arquitetura orientada a colunas favorece consultas de filtragem e agregação (SCHULZE et al., 2024). A versão do servidor não foi registrada.
\item
  \textbf{SciPy (1.18.1):} utilizado em análises estatísticas complementares, incluindo comparações de distribuições e avaliação da concordância entre rankings de importância dos atributos (VIRTANEN et al., 2020).
\end{itemize}

\textbf{Ambiente computacional.} Os experimentos locais foram realizados em um computador com sistema operacional de 64 bits, processador Intel Core i5-1235U e 8 GB de memória RAM instalada, dos quais 7,69 GB eram utilizáveis. Portanto, os tempos de treinamento, inferência e processamento apresentados devem ser interpretados considerando esse ambiente. As consultas ao ClickHouse, hospedado na Railway, também podem depender da infraestrutura remota e da comunicação entre o computador local e o serviço.

\textbf{Configuração experimental e critérios de avaliação.} Foram comparados dois algoritmos (Random Forest e SVM) em duas tarefas, classificação binária (BENIGN × ataque) e multiclasse, considerando configurações padrão e ajustadas. Para permitir uma comparação consistente, foram mantidos os conjuntos e critérios de particionamento definidos para cada tarefa, assim como os procedimentos de pré-processamento estabelecidos para os modelos. A busca de hiperparâmetros foi conduzida sem utilizar o conjunto reservado ao teste.

As métricas preditivas incluem acurácia, precisão, revocação, F1-score e MCC. Também são consideradas ROC-AUC, PR-AUC da classe BENIGN, taxa de falso alarme, acurácia balanceada e métricas por classe, conforme a tarefa e a distribuição dos dados. Como a proporção entre benignos e ataques da validação difere da do treino, as métricas que dependem dessa proporção (precisão, F1 e MCC) são calculadas também com ponderação que a recompõe. Os modelos são avaliados no limiar natural de cada algoritmo e, no binário, em pontos de operação definidos na validação para taxas de falso alarme de 1\% e de 0,1\%. O custo computacional é avaliado pelos tempos de treinamento, busca e inferência, além do uso de recursos quando monitorado. Como baselines internos, utilizam-se as configurações padrão de cada algoritmo, comparadas às configurações selecionadas pela busca de hiperparâmetros.

O critério de sucesso não pressupõe que um algoritmo necessariamente supere o outro. Considera-se que os objetivos experimentais são atendidos quando é possível comparar, nos mesmos cenários, o desempenho preditivo e o custo dos modelos padrão e ajustados nas tarefas binária e multiclasse, identificando as condições em que cada configuração apresenta vantagens ou limitações, e situar os resultados em relação à literatura. A hipótese é discutida à luz das diferenças observadas. Como não foi definido um limiar mínimo de melhoria nem foram aplicados testes de significância a todas as comparações, os resultados são interpretados como evidências experimentais descritivas.

\begin{table}
\centering
\caption{Plano de experimentação}
\label{tab:plano-experimentacao}
\footnotesize
\resizebox{\textwidth}{!}{%
\begin{tabular}{p{2.6cm}p{3.2cm}p{3.0cm}p{3.6cm}p{3.0cm}}
\hline
Experimento & Variáveis comparadas & Variáveis controladas & Métricas principais & Baseline \\
\hline
Classificação binária & RF e SVM; padrão e ajustado & Partições e procedimento de avaliação binária & ROC-AUC, acurácia, precisão, revocação, F1, MCC, taxa de falso alarme e custo & Configuração padrão de cada algoritmo \\
Classificação multiclasse & RF e SVM; padrão e ajustado & Partições e procedimento de avaliação multiclasse & ROC-AUC, acurácia balanceada, F1 macro, MCC, métricas por classe e custo & Configuração padrão de cada algoritmo \\
Busca de hiperparâmetros & Configurações candidatas de cada modelo & Conjunto de desenvolvimento e validação cruzada & Métrica de seleção e tempo de busca & Configuração inicial ou padrão \\
Estabilidade entre sementes & Sementes do RF (padrão e ajustado) & Dados, partições e limiares & Média, desvio e intervalo de MCC, recall e F1 macro & Semente original \\
Explicabilidade & Algoritmo, configuração, classe e dia & Modelos finais e amostras explicadas & Importância, direção e estabilidade das contribuições SHAP & Comparações entre RF e SVM, configurações e dias \\
Simulação do IDS & Algoritmo, configuração e cenário (fases de ataque ou balanceado) & Amostra congelada e mesmos fluxos para comparação & Métricas de classificação e latência de inferência & Comparação entre modelos selecionados no painel \\
Comparação com a literatura & Estudos correlatos e este trabalho & — (conjuntos, divisões e métricas diferem) & Acurácia, precisão, revocação, F1 e MCC, quando reportados & — \\
\hline
\end{tabular}}
\par\smallskip
Fonte: Elaborado pelo autor (2026).
\end{table}

A simulação do IDS é demonstrativa e utiliza uma amostra de 3.600 fluxos, menor que o teste completo; seus resultados ilustram o funcionamento do painel e não substituem as métricas dos experimentos principais.

\section{Detalhes da implementação}
\label{sec:implementacao}
\subsection{Ingestão dos dados no ClickHouse}
\label{subsec:ingestao-clickhouse}

Os dados utilizados provêm do conjunto CICDDoS2019, produzido pelo Canadian Institute for Cybersecurity, da Universidade de New Brunswick (SHARAFALDIN et al., 2019). O conjunto contém registros de tráfego benigno e de ataques de negação de serviço distribuída (DDoS), representados por fluxos de rede com 88 atributos e armazenados em arquivos CSV. As capturas foram realizadas em dois dias distintos. Essa separação foi preservada neste trabalho: os 11 arquivos do diretório \texttt{data/}, referentes ao primeiro dia, compõem a base de treinamento; os sete arquivos do diretório \texttt{data2/}, referentes ao segundo dia, compõem a base de teste. A diferença de composição entre os dias é relevante para a avaliação, pois há classes presentes apenas no segundo dia, como Portmap.

O volume total dos arquivos de origem é de aproximadamente 30 GB, o que motivou o uso de um sistema gerenciador de banco de dados analítico. Adotou-se o ClickHouse, hospedado na plataforma Railway. Por ser orientado a colunas, o sistema é adequado a consultas analíticas sobre grandes volumes de dados e permite executar agregações no servidor (SCHULZE et al., 2024).

Cada um dos 18 arquivos CSV foi associado a uma tabela própria no ClickHouse. As tabelas do primeiro dia usam o prefixo \texttt{data\_}, e as do segundo, \texttt{data2\_}. O esquema comum, com 88 colunas, foi derivado do cabeçalho de um dos arquivos e aplicado às tabelas. Os nomes das colunas foram normalizados para padronizar sua representação; por exemplo, a coluna de índice dos arquivos de origem passou a se chamar \texttt{unnamed\_0}.

As colunas foram armazenadas como texto (\texttt{String}), preservando os valores originais durante a ingestão. Essa escolha também permitiu adiar a conversão de tipos e o tratamento de valores como \texttt{Infinity}, \texttt{NaN} e campos vazios para a etapa de pré-processamento (Seção~3.2.4). Para reduzir o espaço ocupado, foi utilizada compressão ZSTD de nível 3 (COLLET; KUCHERAWY, 2021).

\textbf{Fase 1 -- Criação e validação.} Para cada arquivo, o script cria a tabela correspondente, caso ela ainda não exista, e compara o número esperado de registros -- total de linhas menos o cabeçalho -- com a contagem da tabela no ClickHouse. Tabelas completas são ignoradas; cargas incompletas podem ser retomadas; contagens superiores ao esperado interrompem o processo.

\textbf{Fase 2 -- Carga em blocos.} Os registros ainda não carregados são enviados em blocos de 100.000 linhas, em formato CSV comprimido com gzip. Ao final da carga de cada tabela, sua contagem é conferida novamente com o total esperado. Em caso de falha durante o envio, o script verifica a contagem antes de reenviar o bloco, reduzindo o risco de duplicação.

As 18 tabelas somam 70.427.637 registros. O primeiro dia corresponde a 50.063.112 registros (71,1\% do total), e o segundo, a 20.364.525 (28,9\%). A tabela \texttt{data\_tftp} representa 28,5\% do conjunto. Após a compressão, os dados ocupam aproximadamente 2,7 GiB, em comparação com cerca de 30 GB dos arquivos de origem.

\begin{table}
    \centering
    \caption{Registros por tabela bruta no ClickHouse.}
    \label{tab:ingestao}
    \small
    \begin{tabular}{llr}
        \hline
        \textbf{Dia} & \textbf{Tabela} & \textbf{Registros} \\
        \hline
        Primeiro (treino) & \texttt{data\_drdos\_dns}     & 5.074.413 \\
                          & \texttt{data\_drdos\_ldap}    & 2.181.542 \\
                          & \texttt{data\_drdos\_mssql}   & 4.524.498 \\
                          & \texttt{data\_drdos\_netbios} & 4.094.986 \\
                          & \texttt{data\_drdos\_ntp}     & 1.217.007 \\
                          & \texttt{data\_drdos\_snmp}    & 5.161.377 \\
                          & \texttt{data\_drdos\_ssdp}    & 2.611.374 \\
                          & \texttt{data\_drdos\_udp}     & 3.136.802 \\
                          & \texttt{data\_syn}            & 1.582.681 \\
                          & \texttt{data\_tftp}           & 20.107.827 \\
                          & \texttt{data\_udplag}         & 370.605 \\
        \hline
        \multicolumn{2}{l}{\textit{Subtotal do primeiro dia}}
            & \textit{50.063.112} \\
        \hline
        Segundo (teste) & \texttt{data2\_ldap}    & 2.113.234 \\
                        & \texttt{data2\_mssql}   & 5.775.786 \\
                        & \texttt{data2\_netbios} & 3.455.899 \\
                        & \texttt{data2\_portmap} & 191.694 \\
                        & \texttt{data2\_syn}     & 4.320.541 \\
                        & \texttt{data2\_udp}     & 3.782.206 \\
                        & \texttt{data2\_udplag}  & 725.165 \\
        \hline
        \multicolumn{2}{l}{\textit{Subtotal do segundo dia}}
            & \textit{20.364.525} \\
        \hline
        \multicolumn{2}{l}{\textbf{Total}}
            & \textbf{70.427.637} \\
        \hline
    \end{tabular}
    \legend{Fonte: Elaborado pelo autor (2026).}
\end{table}

A validação da ingestão foi feita comparando, para cada tabela, o número de registros carregados com o total esperado do arquivo correspondente. Essa verificação confirma a igualdade das contagens, mas não a igualdade do conteúdo de cada registro. O procedimento também pressupõe que os arquivos tenham a mesma ordem de colunas e não contenham quebras de linha internas nos campos.

\subsubsection{Análise exploratória dos dados}
\label{sec:analise-exploratoria}

A análise exploratória dos dados (EDA, do inglês \textit{Exploratory Data Analysis}) foi realizada após a ingestão no ClickHouse (Seção 3.2.1) e antes das transformações para modelagem. Seu objetivo foi caracterizar a distribuição, a qualidade e as relações entre os atributos, orientando as etapas posteriores de preparação. Em um dos estudos relacionados com o CICDDoS2019, Abiramasundari e Ramaswamy (2025) também avaliaram seleção de atributos e diferentes classificadores, reforçando a relevância de examinar as características do conjunto antes da modelagem. Nesta etapa, os dados de origem não foram alterados nem removidos.

\paragraph{Estratégia computacional.}
Como o conjunto contém 70.427.637 registros distribuídos em 18 tabelas, as contagens, estatísticas descritivas e análises de valores ausentes, inválidos, duplicatas e outliers foram calculadas por consultas SQL no ClickHouse. Correlações e visualizações foram produzidas em Pandas a partir de uma amostra estratificada por rótulo, limitada a 10.000 registros por classe e incluindo integralmente as classes raras. Os resultados foram salvos em arquivos CSV e figuras na pasta \texttt{eda\_outputs}.

\paragraph{Dimensão, atributos e distribuição das classes.}
As tabelas compartilham 88 colunas: sete identificadores, o rótulo \texttt{label} e 80 atributos candidatos. Embora armazenadas como texto, a análise de conteúdo identificou 46 atributos numéricos inteiros, 33 numéricos decimais e um atributo textual ou misto, \texttt{simillarhttp}.

Foram encontrados 19 rótulos distintos. Algumas grafias representam a mesma classe (como \texttt{UDP-lag} e \texttt{UDPLag}) e foram unificadas no pré-processamento (Seção 3.2.4). A distribuição é fortemente desbalanceada: há 113.828 fluxos benignos (0,1616\%) e 70.313.809 fluxos de ataque (99,8384\%). TFTP é a classe mais frequente, com 20.082.580 registros; WebDDoS é a menos frequente, com 439. Esse desequilíbrio limita o uso da acurácia isolada, pois um classificador que previsse sempre ``ataque'' obteria aproximadamente 99,84\% de acurácia sem identificar tráfego benigno. Por esse motivo, a avaliação posterior considera métricas sensíveis às classes (Seção 3.1). A Figura 3 apresenta a distribuição dos rótulos em escala logarítmica (Figura~\ref{fig:distribuicao-rotulos}).

A análise também mostrou que o nome da tabela não determina necessariamente o rótulo de cada registro. Por exemplo, \texttt{data2\_ldap} contém fluxos LDAP, NetBIOS e benignos; \texttt{data2\_udplag} contém registros UDPLag, Syn, UDP e benignos. Assim, a identificação da classe foi feita pela coluna \texttt{label}, e não pelo nome do arquivo ou da tabela.

\begin{figure}
    \centering
    \includegraphics[width=\textwidth]{assets/02_distribuicao_classes_corrigida.png}
    \caption{Distribuição dos rótulos.}
    \label{fig:distribuicao-rotulos}
\end{figure}

\paragraph{Qualidade e variabilidade dos dados.}
Foram identificados 249.152 campos vazios em \texttt{flow\_bytes\_s}, 4.092.348 ocorrências de infinito positivo em \texttt{flow\_bytes\_s} e \texttt{flow\_packets\_s}, e 6.145 valores não numéricos em \texttt{simillarhttp}. No total, 2.170.750 registros apresentam ao menos um desses valores problemáticos. Não foram encontrados \texttt{NaN} ou infinito negativo. Também foram identificados 13 atributos constantes e sete atributos binários. Esses resultados foram usados para orientar o tratamento de valores inválidos e a seleção de atributos no pré-processamento.

A análise das estatísticas descritivas encontrou valores negativos em atributos cujos valores não deveriam ser negativos, como comprimentos de cabeçalho e tamanho mínimo de segmento. Esses valores não foram corrigidos neste trabalho, por motivos expostos posteriormente na Seção 3.2.4, e constituem uma limitação dos dados.

\paragraph{Duplicatas e valores atípicos.}
Na métrica de duplicidade que compara atributos e rótulo, ignorando os identificadores, o notebook registra 59.927.555 ocorrências de duplicidade dentro das tabelas. A investigação de duplicatas entre treino e teste e de conflitos de rótulo foi conduzida nas etapas posteriores de preparação dos dados (Seção 3.2.4).

Os valores atípicos foram identificados pelo critério de Tukey, com limites de 1,5 e 3 vezes o intervalo interquartil, sem remoção de registros. Em alguns atributos de tamanho de pacote, aproximadamente 31\% dos registros foram classificados como outliers. As proporções variam entre tráfego benigno e ataques, de modo que esses valores foram mantidos, e não tratados automaticamente como erros, decisão preservada também no pré-processamento (Seção 3.2.4).

\paragraph{Correlação.}
Desconsiderando os 13 atributos constantes, analisaram-se 67 atributos sobre a amostra. Encontraram-se 47 pares com $|r| > 0{,}95$, envolvendo 25 atributos candidatos à redundância. Os pares mais correlacionados são \texttt{total\_fwd\_packets} $\times$ \texttt{subflow\_fwd\_packets} (1,0000), \texttt{fwd\_header\_length} $\times$ \texttt{fwd\_header\_length\_1}, \texttt{fwd\_packet\_length\_mean} $\times$ \texttt{avg\_fwd\_segment\_size} e \texttt{flow\_packets\_s} $\times$ \texttt{fwd\_packets\_s} (0,9997).

\paragraph{Relação atributo $\times$ classe.}
Mediu-se o poder de separação univariado entre benigno e ataque pela área sob a curva ROC (\textit{Area Under the Curve} --- AUC) de cada atributo, com separação definida como $|\mathrm{AUC} - 0{,}5| \cdot 2$. Os melhores atributos são \texttt{min\_packet\_length} (AUC 0,9257), \texttt{fwd\_packet\_length\_min} (0,9238), \texttt{inbound} (0,9029), \texttt{fwd\_packet\_length\_mean} e \texttt{avg\_fwd\_segment\_size} (0,8933), \texttt{average\_packet\_size} (0,8801), \texttt{packet\_length\_mean} (0,8694) e \texttt{flow\_bytes\_s} (0,8664). Atributos como \texttt{bwd\_header\_length} têm AUC próxima de 0,13, ou seja, discriminam no sentido inverso. As medianas por classe mostram assinaturas claras. Os ataques DrDoS\_DNS, LDAP, NetBIOS, Portmap e MSSQL têm \texttt{flow\_duration} mediano próximo de 1 e \texttt{flow\_packets\_s} próximo de 2.000.000, enquanto o tráfego benigno tem \texttt{flow\_duration} mediano de 20.911,5.

\paragraph{Síntese.}
A EDA mostrou desbalanceamento extremo, alta duplicidade, três atributos com infinitos ou texto, 13 atributos constantes, valores negativos incoerentes, redundância entre atributos e forte capacidade discriminativa de atributos de tamanho de pacote e duração. Esses achados fundamentam as decisões da Seção 3.2.4.

\subsubsection{Pré-processamento}

O pré-processamento transformou as 18 tabelas brutas em bases para modelagem, aplicando as decisões que a EDA identificou (Seção 3.2.3). A etapa foi executada no ClickHouse, em tabelas com prefixo \texttt{prep\_}, sem alterar as tabelas originais. Preservou-se a divisão oficial do conjunto: as tabelas \texttt{data\_*}, referentes ao primeiro dia, compõem o treino; as tabelas \texttt{data2\_*}, do segundo dia, compõem o teste (SHARAFALDIN et al., 2019).

As decisões que dependem dos dados (constantes, colunas duplicadas, correlação, seleção de atributos e os parâmetros de imputação e padronização) foram definidas apenas com o treino e depois aplicadas ao teste, sem que este influenciasse nenhuma delas, para evitar vazamento de dados.

\paragraph{Tipos e identificadores.}
Como as tabelas brutas armazenam os campos como texto, os atributos foram convertidos para tipos numéricos explícitos: \texttt{Int64} ou \texttt{Float64}, ambos anuláveis. Quatro atributos (\texttt{bwd\_iat\_mean}, \texttt{bwd\_packets\_s}, \texttt{active\_mean} e \texttt{idle\_mean}) foram definidos como decimais, pois a conversão para inteiros produzia valores nulos.

Os identificadores (\texttt{unnamed\_0}, \texttt{flow\_id}, endereços IP, portas e timestamp) foram retirados da matriz de atributos e armazenados em uma tabela auxiliar, associados à linha por \texttt{id\_linha}, um hash de 64 bits. Dessa forma, os modelos não usam diretamente endereços, portas ou horários, mas a rastreabilidade dos registros é preservada.

Na tabela \texttt{data2\_syn}, 482.799 linhas são idênticas, inclusive nos identificadores. Para essa origem, o \texttt{id\_linha} passou a incorporar a posição física do registro, e as linhas repetidas foram mantidas no teste.

\paragraph{Valores ausentes, infinitos e texto.}
Campos vazios, texto não numérico, \texttt{NaN} e valores infinitos foram convertidos para \texttt{NULL}, com registro do motivo. Os totais identificados na EDA foram reproduzidos: 4.092.348 ocorrências de infinito positivo, 249.152 campos vazios e 6.145 valores não numéricos.

Os atributos do CICDDoS2019 foram extraídos dos fluxos pelo CICFlowMeter (SHARAFALDIN et al., 2019). Entre eles estão \texttt{flow\_bytes\_s} e \texttt{flow\_packets\_s}, que expressam taxas por segundo, ou seja, o total de bytes ou de pacotes do fluxo dividido pela duração do fluxo. A documentação do conjunto não trata dos valores infinitos. Por isso, a origem deles foi verificada diretamente nos dados: em todas as linhas com taxa infinita, \texttt{flow\_duration} é igual a zero, o que produz divisão por zero. Isso ocorre em 1.363.236 linhas do treino e 807.514 do teste para \texttt{flow\_packets\_s}, e em 1.114.212 do treino e 807.386 do teste para \texttt{flow\_bytes\_s}. Os campos vazios de \texttt{flow\_bytes\_s} (249.024 no treino e 128 no teste) também ocorrem apenas em fluxos de duração zero. Portanto, esses valores não indicam erro de leitura nem de conversão. Eles resultam de fluxos com duração nula, para os quais a taxa não é definida.

Nenhum registro foi removido por conter valores inválidos, pois isso eliminaria justamente os fluxos de duração zero. Também não se substituiu o valor por zero, porque zero indicaria tráfego nulo, o que não é verdade. Os valores infinitos e vazios foram convertidos para ausentes (\texttt{NULL}), com o motivo registrado em uma coluna de status. Também foram criados indicadores de ausência para \texttt{flow\_bytes\_s} e \texttt{flow\_packets\_s}, que preservam a informação de que a taxa não pôde ser calculada. Após a deduplicação, a base de treino contém 24.915 valores nulos (0,302\%) em cada uma dessas taxas; no teste, são 807.514 (3,965\%). Esses valores serão imputados pela mediana calculada exclusivamente no treino, aplicada posteriormente durante o treinamento de cada modelo.

\paragraph{Atributos constantes e redundantes.}
As 13 constantes identificadas pela EDA foram confirmadas no treino, sem a identificação de constantes adicionais, e removidas. Também foram removidos seis atributos idênticos a outros, valor a valor, no treino: \texttt{subflow\_fwd\_packets}, \texttt{subflow\_fwd\_bytes}, \texttt{subflow\_bwd\_packets}, \texttt{subflow\_bwd\_bytes}, \texttt{rst\_flag\_count} e \texttt{fwd\_header\_length\_1}. Em cada par, manteve-se um atributo representativo. Dois pares quase idênticos, relacionados às médias de tamanho de pacote e de segmento, diferiam em algumas linhas e foram mantidos. Após essas remoções, a matriz-base passou a conter 61 atributos dos 80 candidatos.

\paragraph{Rótulos.}
Foram criadas três colunas de rótulo: \texttt{label\_original}, que preserva o rótulo original sem espaços e nunca é alterada; \texttt{label\_binary}, que vale 0 para BENIGN e 1 para qualquer ataque; e \texttt{label\_multi}, com o nome padronizado da classe. Como a EDA mostrou que o nome do arquivo não corresponde ao rótulo, o rótulo vem sempre da coluna \texttt{label}. Os dois dias de captura grafam algumas classes de forma diferente: o treino usa UDP-lag e o teste, UDPLag; além disso, o treino traz o prefixo \texttt{DrDoS\_} (DrDoS\_LDAP, DrDoS\_MSSQL, DrDoS\_NetBIOS, DrDoS\_UDP, além de DNS, NTP, SNMP e SSDP), ausente no teste. Para que a mesma classe tenha o mesmo nome nos dois conjuntos, UDP-lag e UDPLag foram unificadas em UDPLag, e o prefixo \texttt{DrDoS\_} foi removido na leitura dos dados para o treinamento (DrDoS\_LDAP $\rightarrow$ LDAP, e assim por diante). A decisão foi tomada por nome e pela documentação, sem olhar o teste, e se apoia na tabela de ataques por dia do artigo do conjunto de dados (SHARAFALDIN et al., 2019), segundo a qual esses ataques foram executados nos dois dias. Resultam 13 classes no treino e 8 no teste, das quais 7 existem nos dois lados: BENIGN, LDAP, MSSQL, NetBIOS, Syn, UDP e UDPLag. WebDDoS existe só no treino e Portmap só no teste, e portanto ficam fora do multiclasse.

\paragraph{Deduplicação do treino.}
Os 50.063.112 registros do treino foram reduzidos a 8.257.343 registros únicos. Foram considerados repetidos os registros com valores idênticos em todas as features e com o mesmo rótulo, mantendo-se uma única ocorrência de cada. Foram removidos 41.805.769 (83,51\%). Mantém-se a ocorrência de menor \texttt{id\_linha}, e a coluna \texttt{n\_repeticoes} guarda a contagem. A deduplicação é global entre as tabelas de treino, e as contagens por tabela reproduzem as da EDA. A EDA declarou 59.927.555 repetidas, recalculadas no pipeline, as mesmas tabelas somam 59.934.887, uma diferença de 7.332 (0,012\%). A deduplicação evita que cópias idênticas de um fluxo pesem várias vezes no treino e na validação. O teste oficial não é deduplicado, para preservar a avaliação. As Tabelas 4 e 5 resumem o efeito.

\begin{table}
    \centering
    \caption{Registros antes e depois do pré-processamento, por conjunto}
    \label{tab:registros-preprocessamento}
    \begin{tabular}{p{5cm}rrr}
        \hline
        Conjunto & Antes (brutos) & Depois (bases finais) & Variação \\
        \hline
        Treino (\texttt{data\_*}, 11 tabelas) & 50.063.112 & 8.257.343 & $-83{,}51\%$ \\
        Teste (\texttt{data2\_*}, 7 tabelas) & 20.364.525 & 20.364.525 & 0,00\% \\
        Total & 70.427.637 & 28.621.868 & $-59{,}36\%$ \\
        \hline
    \end{tabular}
    \legend{Fonte: Elaborado pelo autor (2026).}
\end{table}

\begin{table}
    \centering
    \caption{Registros por classe antes e depois da deduplicação do treino}
    \label{tab:registros-classe-deduplicacao}
    \begin{tabular}{p{4.5cm}rrr}
        \hline
        Classe (rótulo original) & Antes & Depois & Variação \\
        \hline
        BENIGN & 56.863 & 50.768 & $-10{,}72\%$ \\
        DNS (\texttt{DrDoS\_DNS}) & 5.071.011 & 113.259 & $-97{,}77\%$ \\
        LDAP (\texttt{DrDoS\_LDAP}) & 2.179.930 & 30.099 & $-98{,}62\%$ \\
        MSSQL (\texttt{DrDoS\_MSSQL}) & 4.522.492 & 207.003 & $-95{,}42\%$ \\
        NTP (\texttt{DrDoS\_NTP}) & 1.202.642 & 1.112.936 & $-7{,}46\%$ \\
        NetBIOS (\texttt{DrDoS\_NetBIOS}) & 4.093.279 & 19.490 & $-99{,}52\%$ \\
        SNMP (\texttt{DrDoS\_SNMP}) & 5.159.870 & 115.316 & $-97{,}77\%$ \\
        SSDP (\texttt{DrDoS\_SSDP}) & 2.610.611 & 891.221 & $-65{,}86\%$ \\
        UDP (\texttt{DrDoS\_UDP}) & 3.134.645 & 1.075.337 & $-65{,}70\%$ \\
        Syn & 1.582.289 & 155.515 & $-90{,}17\%$ \\
        TFTP & 20.082.580 & 4.396.919 & $-78{,}11\%$ \\
        UDPLag (UDP-lag) & 366.461 & 89.066 & $-75{,}70\%$ \\
        WebDDoS & 439 & 414 & $-5{,}69\%$ \\
        \hline
        Total & 50.063.112 & 8.257.343 & $-83{,}51\%$ \\
        \hline
    \end{tabular}
    \legend{Fonte: Elaborado pelo autor (2026).}
\end{table}

\paragraph{Conflitos de rótulo.}
A análise encontrou 405.651 linhas únicas de treino (4,91\%) cujos atributos são iguais, mas os rótulos diferem. Apenas 60 desses casos envolvem uma divergência entre BENIGN e ataque; a maioria está associada a rótulos de ataques de reflexão, como DNS, MSSQL, NetBIOS, SSDP, UDP e TFTP.

Como os identificadores (IP, porta e horário) não fazem parte da matriz de atributos, fluxos de ataques diferentes podem apresentar estatísticas de fluxo idênticas e distinguir-se apenas por esses identificadores. Essa é uma interpretação compatível com os dados, mas não uma causa comprovada. Os grupos foram examinados com os identificadores apenas para diagnóstico; não foram corrigidos nem descartados. Foram marcados pelas variáveis \texttt{conflito\_rotulo} e \texttt{conflito\_binario} para análise de sensibilidade.

\paragraph{Sobreposição entre treino e teste.}
Foram identificadas 13.145.336 linhas do teste (64,55\%) com atributos idênticos aos de pelo menos uma linha do treino. A sobreposição varia conforme a origem: de 0,10\% em Syn a 98,55\% em Portmap. Para os registros benignos, os valores ficam entre 0,98\% e 4,16\%.

Cada registro de teste recebeu a marca \texttt{sobrepoe\_treino}. Nas visões usadas pelos modelos, as linhas com sobreposição foram excluídas para reduzir o risco de avaliar como generalização a classificação de um fluxo já visto no treino. A base oficial completa permanece em \texttt{prep\_final\_teste}. A Tabela 6 compara as contagens por classe; após a remoção da sobreposição, UDPLag fica com apenas 557 registros, o que torna suas métricas instáveis.

\begin{table}
\centering
\caption{Registros por classe no teste (sem deduplicação) e nas visões dos modelos}
\label{tab:teste-classes}
\begin{tabular}{lrrr}
\hline
\textbf{Classe} & \textbf{Teste oficial} & \textbf{Teste sem sobreposição} & \textbf{Removidos} \\
 & & \textbf{com o treino} & \\
\hline
BENIGN  & 56.965    & 56.216    & 749        \\
LDAP    & 1.915.122 & 112.491   & 1.802.631  \\
MSSQL   & 5.787.453 & 854.242   & 4.933.211  \\
NetBIOS & 3.657.497 & 91.175    & 3.566.322  \\
Syn     & 4.891.500 & 4.887.060 & 4.440      \\
UDP     & 3.867.155 & 1.214.746 & 2.652.409  \\
UDPLag  & 1.873     & 557       & 1.316      \\
\hline
\textbf{Total (7 classes)} & 20.177.565 & 7.216.487 & 12.961.078 \\
\hline
Portmap\textsuperscript{*} & 186.960 & \multicolumn{2}{c}{fora (sem par no treino)} \\
\hline
\end{tabular}
\par\vspace{2pt}
{\footnotesize \textsuperscript{*} Portmap existe só no teste e não entra no total das 7 classes. Suas 184.258 linhas com features idênticas ao treino estão incluídas nas 13.145.336 do teste oficial, mas não nas 12.961.078 removidas desta tabela.}
\par\vspace{2pt}
{\footnotesize Fonte: Elaborado pelo autor (2026).}
\end{table}

\paragraph{Desbalanceamento, valores atípicos e negativos.}
Após a deduplicação, o treino contém 8.206.575 ataques (99,39\%) e 50.768 registros benignos (0,61\%). No teste oficial, são 20.307.560 ataques (99,72\%) e 56.965 benignos (0,28\%). Como o desbalanceamento permanece, a etapa de balanceamento é aplicada posteriormente, apenas ao treino.

Os valores atípicos e negativos identificados na EDA foram mantidos. As caudas acentuadas podem refletir características reais do tráfego de rede. Os valores negativos concentram-se em poucos atributos, como comprimento de cabeçalho, janela inicial e tamanho mínimo de segmento, e podem decorrer de estouro numérico ou de valores sentinela do extrator. Um valor sentinela é um código usado para indicar uma condição especial, como por exemplo uma informação indisponível, em vez de representar uma medida real. Como não é possível confirmar a origem desses valores apenas pelos dados, eles foram preservados. Removê-los poderia excluir fluxos de ataque inteiros.

Essa escolha tem implicações distintas para os modelos: o Random Forest utiliza a ordem dos valores nas divisões, enquanto o SVM recebe dados padronizados. Valores extremos podem, contudo, afetar as médias e os desvios empregados na padronização, o que constitui uma limitação do trabalho.

\paragraph{Redundância por correlação e capacidade discriminativa.}
Esta etapa teve três partes, todas decididas apenas com o treino.

Na primeira, a correlação foi recalculada no treino deduplicado, sobre uma amostra determinística de 166.089 linhas, e encontrou 43 pares com $|r| > 0{,}95$ (47 na EDA, pois alguns pares já saíram por constância ou duplicidade exata). Disso resultou uma lista reduzida de 40 atributos, mantida apenas como experimento comparativo.

Na segunda, os atributos com AUC muito alto na EDA foram revalidados no teste oficial, como diagnóstico e não como critério de remoção. \texttt{min\_packet\_length} teve AUC de 0,979 no treino e 0,891 no teste, \texttt{fwd\_packet\_length\_min} de 0,977 e 0,890, e \texttt{inbound} de 0,920 e 0,909. A queda indica dependência parcial do período de captura.

Na terceira, aplicou-se um filtro, sem treinar modelos, para indicar quais atributos seriam descartados por serem redundantes ou inúteis. Ele combina o AUC por atributo, a informação mútua e a correlação entre colunas, e usou uma amostra de 104.730 linhas do treino, com todos os 50.768 registros benignos e cerca de 5.000 linhas por tipo de ataque. Os limiares foram 0,02 para a força de AUC ($|\mathrm{AUC} - 0{,}5| \times 2$), 0,02 para a informação mútua relativa e 0,95 para a correlação. Foram consideradas inúteis as colunas abaixo dos dois primeiros limiares, e gêmeas as colunas com $|r| \geq 0{,}95$ em relação a outra de maior nota, da qual só a de maior nota permanece. O filtro foi executado uma vez para o binário e outra para o multiclasse e indicou, em cada caso, o descarte de 16 dos 61 atributos (15 colunas gêmeas e uma inútil, \texttt{syn\_flag\_count}), restando 45. Os conjuntos de 45 atributos não são idênticos nas duas tarefas. Essa indicação não remove colunas das bases finais, que mantêm os 61 atributos: ela é registrada e aplicada na montagem da entrada de cada modelo, descrita na etapa seguinte. O filtro também marcou como suspeitos de atalho os atributos com força de AUC igual ou superior a 0,80, sem descartá-los: três no binário (\texttt{bwd\_packets\_s}, \texttt{inbound} e \texttt{total\_backward\_packets}) e vários no multiclasse, onde a força é medida de uma classe contra as demais.

\paragraph{Escala e imputação.}
A imputação pela mediana utiliza parâmetros calculados exclusivamente no treino e é aplicada na montagem das visões de entrada. O Random Forest recebe os dados imputados; o SVM recebe os dados imputados e, em seguida, padronizados por z-score, com média e desvio-padrão estimados apenas no conjunto de treino de cada ajuste. O procedimento é detalhado nas seções de construção das visões e de treinamento dos modelos.

\paragraph{Bases finais e validação.}
O pré-processamento gerou \texttt{prep\_final\_treino}, com 8.257.343 linhas, e \texttt{prep\_final\_teste}, com 20.364.525 linhas. As duas bases mantêm a matriz de 61 atributos e dois indicadores de ausência (\texttt{flow\_bytes\_s} e \texttt{flow\_packets\_s}), além dos metadados de rastreio (identificador da linha, tabela de origem, rótulos e marcas de repetição, conflito e sobreposição). Os identificadores brutos (IPs, portas, \texttt{flow\_id} e horário) ficam fora da matriz, em tabela auxiliar. A validação final aprovou as 15 checagens do pipeline, entre elas: tabelas originais intactas; treino final sem duplicatas; teste final igual ao teste oficial; soma de \texttt{n\_repeticoes} igual ao treino bruto; coerência entre o rótulo binário e o rótulo original; ausência de infinitos e \texttt{NaN}; e mesmas colunas e tipos nos dois conjuntos.

\subsubsection{Visões de entrada dos modelos binários}
\label{sec:visoes-modelos-binario}

Com base nas tabelas \texttt{prep\_final\_treino} e \texttt{prep\_final\_teste}, foram criadas duas visões de entrada comuns aos modelos binários, o Random Forest e o SVM. Nelas, a variável resposta \texttt{y} assume o valor 0 para a classe BENIGN e 1 para a classe ataque. Na preparação dessas visões, foram aplicadas três decisões: a ablação de atributos suspeitos de atalho, a imputação de valores ausentes e a remoção, do teste, de linhas sobrepostas ao treino.

\paragraph{Atributos de entrada.}
Da seleção de 45 atributos do modelo binário, descrita na Seção 3.2.1.2.9, foram excluídos por ablação \texttt{bwd\_packets\_s}, \texttt{inbound} e \texttt{total\_backward\_packets}, identificados como suspeitos de atalho por apresentarem força de AUC $\geq 0{,}80$. Essa lista foi definida antes da análise dos resultados dos modelos para verificar se o desempenho dependia desses atributos. Restaram 42 atributos. Somado a eles, o indicador de ausência de \texttt{flow\_bytes\_s} totaliza 43 colunas de entrada, listadas no Apêndice A. Não foi criado indicador para \texttt{flow\_packets\_s}, pois essa taxa foi descartada na seleção por ser gêmea de \texttt{fwd\_packets\_s}.

\paragraph{Valores ausentes.}
Os valores ausentes de \texttt{flow\_bytes\_s} foram substituídos pela mediana do conjunto de treino (689,59). O mesmo valor foi aplicado ao treino e ao teste, sem utilizar os dados de teste no cálculo dos parâmetros. O indicador de ausência preserva a informação de que o valor não pôde ser calculado. Como as visões são comuns, o Random Forest e o SVM recebem os valores já imputados.

\paragraph{Padronização dos atributos (SVM).}
O Random Forest utiliza os atributos na escala original. O SVM com kernel RBF, por depender de distâncias entre os registros, recebe os atributos numéricos padronizados por z-score (média 0 e desvio-padrão 1), realizada por um \texttt{StandardScaler} que integra o mesmo \textit{pipeline} do modelo. A média e o desvio-padrão são estimados apenas no conjunto de treino de cada ajuste e aplicados à validação e ao teste, sem utilizar dados de teste. A padronização incide sobre os 42 atributos numéricos; o indicador de ausência permanece em 0 e 1, pois a padronização o distorceria. O procedimento é detalhado na seção do SVM.

\paragraph{Teste sem sobreposição.}
Para evitar vazamento de dados, foram removidas do teste as linhas cujas features são idênticas às de alguma linha do treino. A avaliação de fluxos já vistos pelo modelo poderia inflar as métricas. O conjunto de treino foi mantido completo, com 8.257.343 registros. Dos 20.364.525 registros do teste oficial, 7.219.189 permaneceram e 13.145.336 (64,55\%) foram removidos, como mostra a Tabela~\ref{tab:registros-rf-binario}.

A remoção atingiu uma proporção maior dos ataques (64,7\%) do que dos registros benignos (1,3\%). Como resultado, a proporção de registros benignos aumentou de 0,28\% no teste oficial para 0,78\% no teste sem sobreposição; no treino, essa proporção é de 0,61\%. Após a remoção, restaram 56.216 registros benignos, quantidade considerada suficiente para estimar as métricas dessa classe.

\begin{table}
    \centering
    \caption{Registros por classe no treino e no teste, com e sem sobreposição ao treino (binário)}
    \label{tab:registros-rf-binario}
    \begin{tabular}{lrrrrr}
        \hline
        Classe & Treino & Teste oficial & Teste sem sobreposição & Removidos & \% removido \\
        \hline
        0 (BENIGN) & 50.768 & 56.965 & 56.216 & 749 & 1,3 \\
        1 (ataque) & 8.206.575 & 20.307.560 & 7.162.973 & 13.144.587 & 64,7 \\
        Total & 8.257.343 & 20.364.525 & 7.219.189 & 13.145.336 & 64,6 \\
        \hline
    \end{tabular}
    \legend{Fonte: Elaborado pelo autor (2026).}
\end{table}

Nota: Os totais diferem dos apresentados para o modelo multiclasse (Tabela 6: 7.216.487 registros no teste sem sobreposição e 12.961.078 removidos) porque, no modelo binário, os registros da classe Portmap permanecem classificados como ataques. Assim, suas 184.258 linhas sobrepostas ao treino são incluídas na remoção. No modelo multiclasse, Portmap é excluída por não ter uma classe correspondente no treino. A diferença de 2.702 registros no teste sem sobreposição corresponde às linhas de Portmap que não se sobrepõem ao treino (186.960 $-$ 184.258).

\paragraph{Validação.}
As visões foram aprovadas nas 13 checagens definidas. Foram verificadas as colunas esperadas, a ausência de rótulos ou atributos retirados na ablação, a correspondência entre as contagens de linhas e os totais esperados para os dois conjuntos, a presença das classes 0 e 1 em \texttt{y} e a ausência de valores nulos ou infinitos. Também foi confirmada a equivalência com a imputação realizada no pré-processamento em amostras do treino (28.940 linhas, das quais 1.357 tiveram valores imputados) e do teste (43.004 linhas, das quais 20.299 tiveram valores imputados).

\paragraph{Apêndice A -- Colunas de entrada dos modelos binários (43).}

\texttt{avg\_fwd\_segment\_size}, \texttt{bwd\_header\_length}, \texttt{max\_packet\_length}, \texttt{bwd\_iat\_max}, \texttt{bwd\_iat\_total}, \texttt{bwd\_iat\_mean}, \texttt{fwd\_packet\_length\_max}, \texttt{total\_length\_of\_fwd\_packets}, \texttt{bwd\_iat\_min}, \texttt{down\_up\_ratio}, \texttt{flow\_bytes\_s}, \texttt{flow\_iat\_max}, \texttt{packet\_length\_std}, \texttt{packet\_length\_variance}, \texttt{bwd\_packet\_length\_max}, \texttt{total\_length\_of\_bwd\_packets}, \texttt{avg\_bwd\_segment\_size}, \texttt{flow\_iat\_mean}, \texttt{fwd\_packets\_s}, \texttt{fwd\_header\_length}, \texttt{flow\_duration}, \texttt{init\_win\_bytes\_forward}, \texttt{flow\_iat\_std}, \texttt{init\_win\_bytes\_backward}, \texttt{bwd\_packet\_length\_min}, \texttt{protocol}, \texttt{urg\_flag\_count}, \texttt{total\_fwd\_packets}, \texttt{flow\_iat\_min}, \texttt{min\_seg\_size\_forward}, \texttt{act\_data\_pkt\_fwd}, \texttt{bwd\_iat\_std}, \texttt{cwe\_flag\_count}, \texttt{bwd\_packet\_length\_std}, \texttt{fwd\_packet\_length\_std}, \texttt{fwd\_psh\_flags}, \texttt{idle\_min}, \texttt{active\_min}, \texttt{idle\_std}, \texttt{active\_max}, \texttt{ack\_flag\_count}, \texttt{active\_std} e \texttt{flow\_bytes\_s\_ausente} (indicador de ausência).


\subsection{Balanceamento do treino binário}

\textbf{Conjunto de partida e necessidade de balanceamento.} O balanceamento foi feito após a deduplicação e o pré-processamento dos dados. Nessa etapa, o conjunto de treino continha 8.257.343 registros: 8.206.575 de ataque e 50.768 benignos. Portanto, havia aproximadamente 162 registros de ataque para cada registro benigno, que correspondia a 0,61\% do conjunto. Com essa diferença, um modelo poderia favorecer a classe majoritária durante o aprendizado (HE; GARCIA, 2009). Por exemplo, um classificador que previsse todos os registros como ataques teria 99,39\% de acurácia nesse conjunto, mas não identificaria nenhum registro benigno. Por isso, a acurácia isolada não é adequada para avaliar o classificador binário: é preciso considerar também a capacidade de detectar ataques e de reconhecer o tráfego benigno.

A técnica escolhida para tratar o desbalanceamento foi a subamostragem (\textit{undersampling}) da classe majoritária, formada pelos ataques. Como a quantidade de benignos é limitada, ela foi usada como referência para dimensionar os conjuntos: 80\% dos 50.768 registros benignos (40.614) foram destinados ao ajuste, conjunto com o qual o modelo aprende, e os 20\% restantes (10.154) foram reservados para a validação. Para evitar que o modelo favoreça a classe majoritária durante o aprendizado, o ajuste foi balanceado, com a mesma quantidade de ataques e benignos, totalizando 81.228 registros. Essa divisão aproveita os dados benignos reais, preserva uma parcela deles para a validação e evita ampliar artificialmente a classe minoritária. A sobreamostragem dos benignos não foi adotada porque exigiria ampliar artificialmente uma classe representada por poucos fluxos reais. Além de reduzir a fidelidade dos dados usados no aprendizado, isso elevaria o volume do conjunto para mais de 16 milhões de registros e o custo computacional, sobretudo para o SVM, cujo tempo de treinamento pode crescer pelo menos quadraticamente com o número de amostras (SCIKIT-LEARN, [s.d.]b). A subamostragem dos ataques, portanto, permite formar um ajuste de tamanho viável e equilibrado, respeitando a quantidade de benignos reais disponível.

\textbf{Formação do ajuste e da validação.} O conjunto de ajuste contém 81.228 registros: 40.614 benignos e 40.614 ataques. A validação contém 60.154 registros, dos quais 10.154 são benignos e 50.000 são ataques. A validação não precisa ter a mesma quantidade de registros em cada classe. Para medir recall e taxa de falso alarme, importa haver registros suficientes de cada classe, pois essas métricas são calculadas separadamente dentro de cada uma. A validação tem mais ataques porque os benignos são limitados (apenas os 20\% reservados estão disponíveis), enquanto os ataques são abundantes e a amostra precisa incluir os diferentes tipos.

A validação contém dados que o modelo não viu durante o ajuste. Sua função é medir, e não ensinar: ela permite avaliar o desempenho em dados novos, verificar a estabilidade das métricas em partições da própria validação e definir limiares de decisão. O limiar determina a partir de que valor um fluxo é classificado como ataque; neste trabalho, ele é definido para obter taxas de falso alarme de 1\% ou 0,1\%. A validação não é balanceada em sua contagem física, e isso não impede essas análises. Entretanto, precisão, F1, MCC e PR-AUC dependem da proporção entre as classes. Por esse motivo, essas métricas são calculadas com pesos que recompõem a proporção de benignos do treino deduplicado (aproximadamente 0,615\%): cada ataque recebe peso 1 e cada benigno recebe peso 0,0305. Recall e taxa de falso alarme são calculados dentro de cada classe e não dependem dessa proporção.

A Tabela 8 apresenta as contagens do conjunto deduplicado, do ajuste, da validação e do teste sem sobreposição ao treino. O teste permanece com sua composição original e não é balanceado, para representar o tráfego real.

\begin{table}
\centering
\caption{Registros por classe após a deduplicação e nos conjuntos do balanceamento binário}
\label{tab:registros-balanceamento-binario}
\resizebox{\textwidth}{!}{%
\begin{tabular}{lrrrr}
\hline
Classe & Treino deduplicado & Ajuste & Validação (linhas) & Teste sem sobreposição \\
\hline
0 (BENIGN) & 50.768 & 40.614 & 10.154 & 56.216 \\
1 (ataque) & 8.206.575 & 40.614 & 50.000 & 7.162.973 \\
\textbf{Total} & \textbf{8.257.343} & \textbf{81.228} & \textbf{60.154} & \textbf{7.219.189} \\
\textbf{BENIGN (\% das linhas)} & \textbf{0,61\%} & \textbf{50,0\%} & \textbf{16,9\%} & \textbf{0,78\%} \\
\hline
\end{tabular}%
}
\par\smallskip
Fonte: Elaborado pelo autor (2026).

\begin{minipage}{\textwidth}
\footnotesize
Nota: Os 16,9\% de benignos na validação correspondem à proporção observada nas linhas selecionadas, antes da aplicação dos pesos. Nas métricas ponderadas, os pesos recompõem aproximadamente a proporção de benignos do conjunto deduplicado, de 0,615\% (cada ataque com peso 1 e cada benigno com peso 0,0305).
\end{minipage}
\end{table}

\textbf{Seleção dos ataques e ponderação da validação.} Embora a classificação desta etapa seja binária, os ataques do conjunto contêm diferentes tipos de tráfego. Após a deduplicação, TFTP era o tipo mais numeroso, com 4.396.919 registros, enquanto WebDDoS era o menos numeroso, com 414. Um sorteio sem estratificação poderia concentrar a amostra nos tipos mais frequentes e deixar pouca representação dos demais. Para reduzir esse risco, os ataques foram selecionados separadamente dentro de cada tipo, com quantidades aproximadamente iguais, usando a semente fixa 20261018. Assim, todos os 12 tipos aparecem no ajuste e na validação. WebDDoS é a exceção em volume: como há apenas 414 registros, 331 foram destinados ao ajuste e 83 à validação.

Na validação, a proporção física de benignos (16,9\%) difere da do treino (0,61\%), pois os benignos são os 20\% reservados e os ataques são uma amostra de 50.000 registros. Para recompor a proporção do treino, cada ataque recebe peso 1 e cada benigno recebe o peso $w = \frac{50.768}{8.206.575} \times \frac{50.000}{10.154} \approx 0{,}0305$. Assim, os 10.154 benignos somam cerca de 309 frente a 50.000 ataques, o que corresponde a 0,615\% de benignos. Não se usam pesos proporcionais ao tamanho de cada tipo de ataque: como a amostra é estratificada (cerca de 4.538 ataques por tipo), o peso 1 dá importância semelhante a todos os tipos, enquanto pesos proporcionais fariam o TFTP representar cerca de 54\% dos ataques e dominar as métricas.

A Tabela 9 detalha as quantidades selecionadas por tipo. O tipo de ataque serviu para estratificar e auditar a seleção; não foi incluído como atributo de entrada do classificador.

\begin{table}
\centering
\caption{Distribuição dos registros selecionados por classe e tipo de ataque}
\label{tab:distribuicao-registros-tipo-ataque}
\resizebox{\textwidth}{!}{%
\begin{tabular}{lrrr}
\hline
Classe ou tipo & Total no treino deduplicado & Ajuste & Validação (linhas) \\
\hline
BENIGN & 50.768 & 40.614 & 10.154 \\
TFTP & 4.396.919 & 3.663 & 4.538 \\
NTP & 1.112.936 & 3.662 & 4.538 \\
UDP & 1.075.337 & 3.662 & 4.538 \\
SSDP & 891.221 & 3.662 & 4.538 \\
MSSQL & 207.003 & 3.662 & 4.538 \\
Syn & 155.515 & 3.662 & 4.538 \\
SNMP & 115.316 & 3.662 & 4.538 \\
DNS & 113.259 & 3.662 & 4.538 \\
UDPLag & 89.066 & 3.662 & 4.538 \\
LDAP & 30.099 & 3.662 & 4.538 \\
NetBIOS & 19.490 & 3.662 & 4.537 \\
WebDDoS & 414 & 331 & 83 \\
\textbf{Total} & \textbf{8.257.343} & \textbf{81.228} & \textbf{60.154} \\
\hline
\end{tabular}%
}
\par\smallskip
Fonte: Elaborado pelo autor (2026).
\end{table}

\textbf{Papel do ajuste, da validação e do refit.} O ajuste é o conjunto balanceado usado para o modelo aprender os padrões das duas classes. Como 99,39\% dos registros disponíveis são ataques, treinar diretamente com todos eles faria com que a classe majoritária tivesse muito mais influência no aprendizado. O balanceamento do ajuste reduz esse efeito. A validação fica separada durante essa etapa e é usada para medir o desempenho em dados não vistos, examinar a estabilidade das métricas em suas partições e definir os limiares de decisão associados às taxas de falso alarme de 1\% e 0,1\%. Portanto, as métricas calculadas nessa etapa descrevem o desempenho do modelo avaliado na validação; ela não é usada para ensinar o modelo.

Depois dessas análises, o \textit{refit} aproveita mais registros benignos no ajuste final. Ele mantém os 81.228 registros do ajuste e incorpora 10.154 benignos e 10.154 ataques selecionados da validação, totalizando 101.536 registros, divididos igualmente entre as classes. Essa etapa permite usar todos os 50.768 benignos disponíveis. Como parte da validação é incorporada ao \textit{refit}, ela deixa de ser independente para medir o desempenho do modelo reajustado; sua função de avaliação corresponde à etapa em que os dados ainda estavam separados do ajuste.

\begin{table}
\centering
\caption{Composição do conjunto utilizado no refit}
\label{tab:composicao-conjunto-refit}
\begin{tabular}{lrrr}
\hline
Classe & Ajuste & Registros incorporados da validação & Total no refit \\
\hline
0 (BENIGN) & 40.614 & 10.154 & 50.768 \\
1 (ataque) & 40.614 & 10.154 & 50.768 \\
\textbf{Total} & \textbf{81.228} & \textbf{20.308} & \textbf{101.536} \\
\hline
\end{tabular}
\par\smallskip
Fonte: Elaborado pelo autor (2026).
\end{table}

\textbf{Armazenamento e verificações.} A seleção, os dados selecionados e as visões de ajuste e validação foram registrados no ClickHouse. A tabela de seleção preserva o papel de cada linha, o tipo de ataque e um peso de representatividade por tipo (registros restantes do tipo divididos pelos sorteados). Esse peso serve para auditar a amostragem e não é usado nas métricas de desempenho. Nas visões destinadas ao modelo, o tipo de ataque não aparece como atributo; a visão de validação inclui a coluna de peso (não utilizada no treino).

O notebook aprovou 15 verificações. Entre elas, foram confirmados o balanceamento do ajuste, a ausência de linhas compartilhadas entre ajuste e validação, a presença dos 12 tipos de ataque no ajuste, a soma dos pesos de representatividade dos ataques (8.165.961) igual aos ataques restantes e a proporção de benignos esperada na validação. Também se verificou, em uma amostra de 7.048 linhas, a correspondência dos dados selecionados com a visão de origem. As visões não apresentaram valores nulos nos atributos. Por fim, foram identificadas 18 linhas em conflito: registros com características iguais associados aos rótulos BENIGN e ataque. Esses casos foram contabilizados, mas não removidos.

\subsection{Treino do Random Forest binário}

\textbf{Objetivo e configuração.} Esta etapa treina o Random Forest (RF) binário (BENIGN $\times$ ataque) com uma configuração inicial fixa, mede seu desempenho na validação e registra o custo computacional. A busca de hiperparâmetros é descrita adiante, e o teste é feito uma única vez, em etapa posterior. O RF combina várias árvores de decisão: cada árvore aprende a partir de uma amostra dos dados, e a probabilidade de ataque corresponde à fração de árvores que votam nessa classe.

Foram usadas 100 árvores (\texttt{n\_estimators=100}). Esse parâmetro define quantas árvores compõem o modelo: aumentá-lo pode tornar as previsões mais estáveis, mas também eleva o tempo de treinamento e o tamanho do modelo. A profundidade máxima foi deixada sem limite (\texttt{max\_depth=None}), permitindo que as árvores continuem criando divisões enquanto houver dados suficientes. Para controlar essas divisões, foram definidos \texttt{min\_samples\_split=2}, o número mínimo de registros necessário para dividir um nó, e \texttt{min\_samples\_leaf=1}, o número mínimo de registros permitido em uma folha. Esses valores permitem árvores detalhadas, mas também podem torná-las mais ajustadas às particularidades dos dados de treino.

O parâmetro \texttt{max\_features} foi definido como \texttt{"sqrt"}: em cada divisão, cada árvore considera uma amostra aleatória de \(\sqrt{p}\) atributos, em que \(p\) é o número total de atributos de entrada (Tabela 11). Essa seleção introduz diversidade entre as árvores e reduz a dependência de um único conjunto de atributos. Os demais parâmetros seguem os valores padrão do scikit-learn 1.9.1. A semente fixa (\texttt{random\_state=42}) permite reproduzir a aleatoriedade do treinamento, enquanto \texttt{n\_jobs=-1} permite que o ajuste utilize todos os núcleos disponíveis, reduzindo o tempo de execução.

\begin{table}
\centering
\caption{Configuração inicial do Random Forest}
\label{tab:configuracao-inicial-rf}
\begin{tabular}{ll}
\hline
Hiperparâmetro & Valor \\
\hline
\texttt{n\_estimators} (número de árvores) & 100 \\
\texttt{max\_depth} (profundidade máxima) & sem limite \\
\texttt{min\_samples\_split} & 2 \\
\texttt{min\_samples\_leaf} & 1 \\
\texttt{max\_features} & raiz quadrada do nº de atributos \\
\texttt{random\_state} & 42 \\
\texttt{n\_jobs} & $-1$ (todos os núcleos) \\
\hline
\end{tabular}
\par\smallskip
Fonte: Elaborado pelo autor (2026).
\end{table}

\textbf{Preparação.} O código verifica se as etapas anteriores foram executadas e confere se as contagens dos conjuntos correspondem às do balanceamento: 50.768 benignos e 8.206.575 ataques no treino deduplicado; 40.614 registros de cada classe no ajuste; e 10.154 benignos e 50.000 ataques na validação. O modelo recebe 43 atributos: 42 numéricos e o indicador \texttt{flow\_bytes\_s\_ausente}. A seleção desses atributos é apresentada na Seção 3.2.1.2.9, e a composição das visões de entrada, na Seção 3.2.5. Os valores ausentes já haviam sido substituídos pela mediana do treino nessas visões. Para conferir essa preparação, o notebook reconstruiu os fluxos da validação antes do tratamento: 1.132 dos 60.154 registros tinham algum valor ausente. Depois, comparou os atributos resultantes com os das visões e obteve resultados idênticos. A configuração inicial do Random Forest foi fixada para permitir a reprodução do experimento: 100 árvores, \texttt{max\_features="sqrt"}, profundidade máxima sem limite, mínimo de duas amostras para dividir um nó, mínimo de uma amostra por folha e semente aleatória 42.

\textbf{Métricas de avaliação.} Como 99,39\% do treino e 99,22\% do teste são ataques, a acurácia, a precisão, o F1 e a PR-AUC calculados com os ataques como classe positiva ficam próximos de 1 mesmo para um modelo que erre quase todos os benignos (HE; GARCIA, 2009). Um classificador que previsse tudo como ataque, por exemplo, teria 99,39\% de acurácia no treino. Por isso, essas métricas não bastam para avaliar o modelo. Elas foram mantidas, mas complementadas por métricas que observam diretamente a classe rara (BENIGN) e por métricas que não dependem da proporção entre as classes. As conclusões apoiam-se principalmente em recall, taxa de falso alarme, MCC e nas métricas do BENIGN.

Sejam VP, FP, VN e FN os ataques detectados, os benignos classificados como ataque, os benignos reconhecidos e os ataques perdidos. Não dependem da proporção entre as classes, pois são calculadas dentro de cada classe: o recall, fração dos ataques detectados, \(VP/(VP+FN)\); a taxa de falso alarme, fração dos benignos classificados como ataque, \(FP/(FP+VN)\); a especificidade, fração dos benignos reconhecidos corretamente (1 menos a taxa de falso alarme); a acurácia balanceada, média do recall e da especificidade; e a ROC-AUC, que indica o quanto o modelo ordena os ataques acima dos benignos, de 0,5 (acaso) a 1 (ordenação perfeita), sem depender de um limiar. Dependem da proporção: a acurácia, fração de acertos; a precisão, fração dos fluxos classificados como ataque que são de fato ataques; o F1, média harmônica da precisão e do recall; a PR-AUC, área sob a curva precisão $\times$ recall; e o MCC, coeficiente de correlação de Matthews, que mede a correlação entre a classificação e o rótulo real, de $-1$ a 1, com \(\mathrm{MCC}=\frac{VP\times VN-FP\times FN}{\sqrt{(VP+FP)(VP+FN)(VN+FP)(VN+FN)}}\). Para a classe BENIGN, a precisão é a fração dos fluxos classificados como benignos que são de fato benignos, o F1 é a média harmônica entre essa precisão e o recall do BENIGN (a especificidade), e a PR-AUC é calculada com o BENIGN como classe positiva, tendo como linha de base a proporção de benignos no conjunto. Como as métricas que dependem da proporção mudam com a composição do conjunto, na validação elas são calculadas com pesos que recompõem a proporção do treino, como descrito a seguir; no teste, que já mantém a proporção observada, são calculadas sem pesos.

\textbf{Treino e avaliação na validação.} O Random Forest foi ajustado somente com os 81.228 registros do ajuste e, em seguida, avaliado nos 60.154 registros da validação, que não haviam sido usados para esse aprendizado. As métricas são apresentadas em duas colunas: sem ponderação, considerando as contagens físicas da validação, e ponderadas para representar a proporção do treino deduplicado. Na coluna ponderada, cada ataque recebe peso 1 e cada registro benigno, peso 0,0305. Como a amostra de ataques foi estratificada por tipo, esse procedimento recompõe a proporção entre as classes sem dar maior peso aos tipos de ataque mais numerosos.

O limiar de decisão define a partir de que probabilidade estimada um fluxo é classificado como ataque. Com o limiar padrão de 0,5, o fluxo é classificado como ataque quando pelo menos metade das árvores vota nessa classe. Nessa condição, o modelo identificou corretamente 49.961 dos 50.000 ataques e classificou incorretamente como benignos 39; entre os 10.154 benignos, 8 foram classificados como ataques, como mostra a Tabela 12. Recall e taxa de falso alarme são iguais nas duas colunas porque são calculados separadamente dentro de cada classe. Já as métricas que dependem da proporção entre as classes mudam: o MCC passa de 0,997 sem ponderação para 0,942 com ponderação, e a precisão da classe BENIGN cai de 0,996 para 0,888. Considerando a proporção do treino deduplicado, cerca de 11\% dos fluxos classificados como benignos seriam ataques. Por isso, o MCC e as métricas da classe BENIGN são especialmente relevantes neste problema; a precisão e o F1 dos ataques permanecem próximos de 1 mesmo quando não refletem adequadamente o desempenho sobre a classe minoritária.

\begin{table}
\centering
\caption{Métricas do RF na validação (limiar 0,5)}
\label{tab:metricas-rf-validacao}
\resizebox{\textwidth}{!}{%
\begin{tabular}{lccc}
\hline
Métrica & Sem pesos & Proporção real & Depende da proporção? \\
\hline
ROC-AUC & 0,99994 & 0,99994 & não \\
Recall (ataques) & 0,99922 & 0,99922 & não \\
Taxa de falso alarme & 0,00079 & 0,00079 & não \\
Acurácia balanceada & 0,99922 & 0,99922 & não \\
Acurácia & 0,99922 & 0,99922 & sim \\
Precisão (ataques) & 0,99984 & 1,00000 & sim \\
F1 (ataques) & 0,99953 & 0,99961 & sim \\
MCC & 0,99722 & 0,94157 & sim \\
PR-AUC (ataques) & 0,99998 & 1,00000 & sim \\
Precisão (BENIGN) & 0,99617 & 0,88795 & sim \\
F1 (BENIGN) & 0,99769 & 0,94030 & sim \\
F1 macro & 0,99861 & 0,96996 & sim \\
PR-AUC (BENIGN) & 0,99980 & 0,99612 & sim \\
\hline
\end{tabular}%
}
\par\smallskip
Fonte: Elaborado pelo autor (2026).
\end{table}

\textbf{Pontos de operação.} O limiar de decisão transforma a probabilidade estimada pelo modelo em uma classificação: se ela for igual ou superior ao limiar, o fluxo é classificado como ataque; caso contrário, é classificado como benigno. O limiar padrão é 0,5. Também foram avaliados dois pontos de operação para situações em que se deseja limitar os falsos alarmes: um com taxa máxima de 1\% e outro com taxa máxima de 0,1\%. Em cada caso, escolheu-se na validação o limiar que produz o maior recall sem ultrapassar a taxa de falso alarme definida. Os limiares escolhidos foram registrados para uso posterior (Tabela 13).

Como o modelo tem 100 árvores, a probabilidade de ataque só pode variar em incrementos de 0,01: ela corresponde à fração de árvores que votaram nessa classe. Assim, com limiar 0,05, basta que pelo menos cinco árvores votem em ataque para que o fluxo receba essa classificação. Na validação, esse ponto alcançou recall de 0,99994 e taxa de falso alarme de 0,00837, ou 0,837\%, abaixo do limite de 1\%. Para o limite de 0,1\%, o limiar escolhido foi 0,46: pelo menos 46 árvores precisam votar em ataque. Esse ponto produziu recall de 0,99928 e taxa de falso alarme de 0,00098, ou 0,098\%.

Os pontos mostram o compromisso entre detectar ataques e evitar alarmes falsos: diminuir o limiar tende a aumentar o recall, mas também pode classificar mais tráfego benigno como ataque. As taxas observadas são estimativas de desenvolvimento, pois o limiar foi escolhido e avaliado na mesma validação; por isso, podem ser otimistas. A estimativa de 0,1\% também se baseia em poucos casos: entre 10.154 benignos, essa taxa corresponde a aproximadamente 10 falsos alarmes. Nesse volume, a mudança de apenas um falso alarme altera a taxa em cerca de 0,01 ponto percentual, o que limita a precisão da estimativa.

\begin{table}
\centering
\caption{Recall com taxa de falso alarme fixa (validação)}
\label{tab:recall-falso-alarme-fixo}
\resizebox{\textwidth}{!}{%
\begin{tabular}{lccccc}
\hline
Falso alarme-alvo & Limiar & Recall & Falso alarme obtido & MCC (prop. real) & F1 BENIGN (prop. real) \\
\hline
1\% & 0,05 & 0,99994 & 0,00837 & 0,99092 & 0,99097 \\
0,1\% & 0,46 & 0,99928 & 0,00098 & 0,94558 & 0,94452 \\
\hline
\end{tabular}%
}
\par\smallskip
Fonte: Elaborado pelo autor (2026).
\end{table}

\textbf{Estabilidade do resultado.} Para verificar se as métricas se mantêm semelhantes em diferentes partes da validação, ela foi dividida em cinco partições disjuntas e estratificadas. Cada partição contém 12.031 registros, dos quais 2.031 são benignos, exceto a última, que contém 12.030 registros e 2.030 benignos. O mesmo modelo, sem novo treinamento, foi avaliado em cada partição (Tabela 14). Essa análise mostra a variação causada pela composição dos dados usados na avaliação; ela não mede a variação causada por treinar modelos diferentes. O recall variou de 0,9989 a 0,9996, e o ROC-AUC, de 0,99975 a 1,00000. A taxa de falso alarme variou de 0 a 0,148\%, e o MCC, de 0,921 a 0,968. Essas duas métricas oscilaram mais porque dependem do número de erros entre os benignos, que são a classe minoritária.

Como complemento, foi feita validação cruzada estratificada de cinco dobras (\textit{5-fold cross-validation}) no conjunto de ajuste. Primeiro, os 81.228 registros foram divididos em cinco partes, mantendo a proporção entre benignos e ataques. Em cada uma das cinco rodadas, quatro partes foram usadas para treinar um novo Random Forest e a parte restante para avaliá-lo. Ao final, cada registro foi usado uma vez na avaliação e quatro vezes no treinamento. Cada modelo foi, portanto, treinado com 64.982 ou 64.983 registros, aproximadamente quatro quintos do ajuste. Diferentemente da análise anterior, essa validação cruzada envolve cinco treinamentos; ela mostra como os resultados variam conforme os registros incluídos no treino e na parte reservada para avaliação.

Nas cinco dobras, o recall médio foi 0,99879, com desvio-padrão de 0,00028, e a taxa de falso alarme média foi 0,098\%, com desvio-padrão de 0,047 ponto percentual. O recall médio foi um pouco inferior ao valor de 0,99922 obtido na validação com o modelo treinado em todo o conjunto de ajuste. Essa diferença é compatível com o fato de que, em cada dobra, o modelo foi treinado com menos registros.

\begin{table}
\centering
\caption{Dispersão das métricas (média $\pm$ desvio-padrão; mín--máx)}
\label{tab:dispersao-metricas-rf}
\resizebox{\textwidth}{!}{%
\begin{tabular}{lcc}
\hline
Métrica & 5 partições da validação & 5 dobras da validação cruzada \\
\hline
Recall & 0,99922 $\pm$ 0,00029 (0,9989--0,9996) & 0,99879 $\pm$ 0,00028 (0,99852--0,99914) \\
Falso alarme & 0,00079 $\pm$ 0,00056 (0--0,00148) & 0,00098 $\pm$ 0,00047 (0,00037--0,00160) \\
ROC-AUC & 0,99994 $\pm$ 0,00011 & 0,99996 $\pm$ 0,00006 \\
MCC (prop. real / sem pesos) & 0,94205 $\pm$ 0,01962 & 0,99781 $\pm$ 0,00055 \\
\hline
\end{tabular}%
}
\par\smallskip
Fonte: Elaborado pelo autor (2026).
\end{table}

\begin{figure}
    \centering
    \caption{Validação cruzada no ajuste (5 dobras)}
    \label{fig:validacao-cruzada-ajuste-rf}
    \includegraphics[width=\textwidth]{assets/fig_cv_ajuste_5dobras.png}
    \par\smallskip
    Fonte: Elaborado pelo autor (2026).
\end{figure}

\textbf{Refit (modelo final).} O treino final usa a mesma configuração com 101.536 registros: todos os 50.768 benignos e 50.768 ataques, sendo os 10.154 acrescentados vindos da validação, em quantidade igual por tipo (de 83 a 916). A seleção é gravada em \texttt{prep\_bal\_bin\_refit\_sel} para que o SVM use o mesmo conjunto. O treino foi repetido 3 vezes, e as repetições geraram modelos idênticos (conferido nas previsões), o que confirma que a semente fixa o resultado. Como parte da validação entrou no treino, o modelo final não é medido nela; sua qualidade aparece só no teste.

\textbf{Custo computacional.} O treino foi medido em um computador com Windows 11, processador Intel de 10 núcleos físicos (12 lógicos) e 7,7 GiB de RAM, com Python 3.13.14 e scikit-learn 1.9.1. A medição considera somente o ajuste do modelo (fit); a leitura dos dados no ClickHouse e o pré-processamento ficam de fora. O treino no ajuste levou 5,4 s de relógio e 29,3 s de CPU, o que equivale a cerca de 5,4 dos 12 núcleos lógicos ocupados em média. O refit levou 7,1 s (mediana de 3 repetições; mín. 6,0 s, máx. 8,3 s) e 48,1 s de CPU, cerca de 6,8 núcleos ocupados em média. O tempo do ajuste vem de uma única execução. Como as 5 dobras da validação cruzada, treinadas com menos registros (cerca de 65 mil), levaram de 5,5 s a 8,9 s, a variação entre execuções é do mesmo tamanho da diferença entre os conjuntos, e por isso os tempos não permitem concluir como o custo cresce com o número de registros. O acréscimo de memória durante o treino foi de 60 MiB no ajuste e de até 79 MiB no refit, pequeno diante dos 7,7 GiB disponíveis. O modelo final, com 100 árvores, tem 48.162 nós (cerca de 482 por árvore), profundidade média de 21,1 (máxima de 26) e ocupa 0,9 MiB em arquivo compactado (joblib, compressão de nível 3). O modelo do ajuste tem 42.468 nós, profundidade média de 20,4 (máxima de 27) e 0,8 MiB. A Tabela~\ref{tab:custo-treino-tamanho-modelos}, na seção da configuração ajustada, compara esses valores com os do modelo ajustado.


Na inferência, com o objetivo de medir o tempo necessário para classificar um fluxo individual, contado desde o início do processamento até a obtenção da previsão, a latência de um fluxo com 1 núcleo foi de 11,4 ms (mediana) e 27,4 ms (p99) fim a fim. Com todos os núcleos piorou, para 49,8 e 57,8 ms, pois o custo de dividir o trabalho supera o ganho para um único fluxo. Em lote ocorre o inverso: com todos os núcleos, 257.698 fluxos/s contra 87.512 com 1 núcleo (Tabela~\ref{tab:inferencia-modelo-final}). O pré-processamento acrescenta só 0,09 a 0,16 ms por fluxo. A latência é medida em 1.000 fluxos e a vazão classifica a validação inteira (60.154 fluxos). A leitura do banco não entra, e esses números medem velocidade, não qualidade.

\begin{table}
\centering
\caption{Inferência do modelo final (fim a fim)}
\label{tab:inferencia-modelo-final}
\begin{tabular}{lcc}
\hline
Medida & \texttt{n\_jobs = 1} & \texttt{n\_jobs = $-1$} \\
\hline
Latência p50 / p95 / p99 (ms) & 11,4 / 22,5 / 27,4 & 49,8 / 54,2 / 57,8 \\
Vazão em lote (fluxos/s) & 87.512 & 257.698 \\
\hline
\end{tabular}
\par\smallskip
Fonte: Elaborado pelo autor (2026).
\end{table}

\textbf{Verificação.} As 14 checagens foram aprovadas. Entre elas, o modelo gravado reproduz as previsões, a ordem dos atributos e as medianas gravadas são as usadas, o ajuste e o refit estão balanceados (40.614 + 40.614 e 50.768 + 50.768), o refit contém todo o ajuste, e os pesos preservam a proporção do treino.

\subsubsection{Avaliação do Random Forest binário no teste}
\label{sec:rf-bin-teste}

\textbf{Preparação e composição do teste.} Nesta etapa, o modelo final do Random Forest binário (o treinado no \textit{refit}) foi avaliado no conjunto de teste. O teste foi mantido separado durante todo o processo de treinamento: não foi usado para treinar o modelo, definir sua configuração, escolher os limiares de decisão nem calcular as medianas do pré-processamento, que vêm só do treino. Por isso, o modelo nunca teve contato com ele antes desta avaliação. Ao fim do treinamento, foram gravados o modelo final, a configuração com os hiperparâmetros, o pré-processamento, com os nomes e a ordem dos 43 atributos e as medianas do treino, e os limiares obtidos na validação. Esses limiares são o padrão, de 0,5, e os de 0,05 e 0,46, selecionados para taxas de falso alarme de 1\% e 0,1\%. Eles foram obtidos com o modelo treinado só no ajuste e aplicados ao modelo final, refeito com mais registros benignos. Antes das previsões, o \textit{notebook} de teste carregou esses arquivos e conferiu que os hiperparâmetros do modelo carregado coincidem com os registrados, que a execução avaliada é a do treinamento (mesmo identificador) e que a ordem dos 43 atributos do modelo é a da visão de teste. O teste foi processado uma única vez, e o \textit{notebook} bloqueia, por padrão, nova execução depois que o resultado é gravado. Os limiares definidos a partir dos próprios resultados do teste são apresentados apenas como informação diagnóstica.

O teste foi deduplicado em relação ao treino: foram removidos os registros cujas características eram idênticas às de algum registro usado no treino, mantendo-se o treino completo. Essa decisão evita que o modelo seja avaliado em fluxos com características que já encontrou durante o aprendizado, o que poderia produzir uma estimativa otimista de desempenho. A deduplicação foi feita entre os conjuntos, e não dentro do próprio teste. Assim, as repetições internas do teste foram mantidas; a análise exploratória havia apontado elevada redundância nesse conjunto, o que deve ser considerado ao interpretar os resultados.

Dos 20.364.525 registros do teste completo, foram removidos 13.145.336 por sobreposição com o treino: 13.144.587 ataques (64,7\% dos ataques) e 749 benignos (1,3\% dos benignos). Restaram 7.219.189 registros, dos quais 56.216 eram benignos e 7.162.973 eram ataques. A remoção incidiu principalmente sobre os ataques e alterou a proporção de classes: os benignos passaram de aproximadamente 0,28\% no teste completo para 0,779\% no teste sem sobreposição. Portanto, as métricas apresentadas descrevem esse conjunto filtrado, cuja composição difere da do teste original.

\textbf{Execução e resultado do modelo padrão.} Os registros foram processados em 145 lotes de até 50.000 linhas. A previsão dos 7.219.189 registros levou 20,4~s no total (353.881 fluxos por segundo, com todos os núcleos), sem contar a leitura no ClickHouse nem a extração de atributos do tráfego. Para um fluxo isolado, a latência mediana fim a fim foi de 7,660~ms com um núcleo e 49,970~ms com todos; em lote, a vazão fim a fim foi de 167.330 e 270.218 fluxos por segundo, respectivamente.

Com o limiar de 0,5, o modelo obteve ROC-AUC de 0,98993 e recall de ataques de 0,95111. A taxa de falso alarme foi 2,095\%, com 1.178 benignos classificados como ataques, e 350.177 ataques foram classificados como benignos. O MCC foi 0,35521 e a precisão da classe BENIGN foi 0,13582. Esses resultados mostram por que a avaliação não deve se limitar à precisão ou ao F1 dos ataques: com a predominância dessa classe, essas métricas ficam próximas de 1, embora o modelo ainda cometa muitos erros relevantes sobre os benignos. Dos ataques, 344.675 (4,8\%) receberam escore entre 0,46 e 0,5, o que mostra a sensibilidade do resultado ao limiar.

Os limiares selecionados na validação para taxas de falso alarme de 1\% e 0,1\% (0,05 e 0,46) produziram, no teste, taxas de 26,788\% e 2,654\%. A Tabela~\ref{tab:rf-padrao-teste} resume os três pontos. As taxas estimadas na validação não se mantiveram no teste, apesar do recall elevado.

\begin{table}
\centering
\caption{Random Forest padrão no teste}
\label{tab:rf-padrao-teste}
\small
\begin{tabular}{lrrrrr}
\hline
Limiar & Recall de ataques & Falso alarme & MCC & Falsos negativos & Falsos positivos \\
\hline
0,5 & 0,95111 & 2,095\% & 0,35521 & 350.177 & 1.178 \\
0,05 (meta de 1\% na validação) & 1,00000 & 26,788\% & 0,85450 & 23 & 15.059 \\
0,46 (meta de 0,1\% na validação) & 0,99923 & 2,654\% & 0,94001 & 5.502 & 1.492 \\
\hline
\end{tabular}
\legend{Fonte: Elaborado pelo autor (2026).}
\end{table}

\textbf{Diferença entre os dias de treino e de teste.} A queda de desempenho do treino para o teste precisa ser lida à luz da origem dos dados. O CICDDoS2019 reúne capturas feitas em dois dias distintos: os dados de treino foram coletados em 12/01/2019 e os de teste em 11/03/2019, com conjuntos de ataques parcialmente diferentes \cite{sharafaldin2019}. 
Isso aparece nos dados deste trabalho: o tipo Portmap só existe no teste, e o WebDDoS só existe no treino. A validação, por outro lado, é um sorteio de linhas do próprio dia de treino. Ela mede a capacidade do modelo de reconhecer fluxos do mesmo dia, e não de generalizar para outro dia de captura. A diferença entre os dois resultados é grande. Na validação, o modelo padrão obteve ROC-AUC de 0,99994, recall de 0,99922 e falso alarme de 0,079\% no limiar 0,5. No teste, esses valores foram 0,98993, 0,95111 e 2,095\%, ou seja, um falso alarme mais de vinte vezes maior. Os limiares escolhidos na validação não se transferiram pelo mesmo motivo: a calibração feita num dia não descreve o outro. Essa é a explicação mais plausível para o conjunto dos resultados, mas não foi isolada por um experimento próprio. Efeitos de outras fontes, como a remoção de linhas repetidas do teste (que altera a proporção de classes, mas não a taxa de falso alarme dentro dos benignos), podem contribuir. Por isso, as métricas do teste não podem ser comparadas diretamente com as da validação, e todos os resultados seguintes devem ser lidos com essa ressalva.

% FIGURA (APENDICE): curvas_teste.png do notebook 06 (modelo padrao)
% Pasta: artefatos_modelos/rf_binario/rf_bin_20261021_221311/curvas_teste.png
% Sugestao: citar aqui "(Apêndice X mostra as curvas ROC e de precisão × recall do BENIGN do modelo padrão)".


\subsubsection{Busca de hiperparâmetros}
\label{sec:rf-bin-busca}

\textbf{Sementes e protocolo de busca.} O Random Forest usa o acaso em dois pontos: cada árvore é treinada com uma amostra sorteada das linhas (\textit{bootstrap}) e, em cada divisão, só um subconjunto sorteado dos atributos é considerado. A semente (\texttt{random\_state}) fixa esses sorteios. Com os mesmos dados e os mesmos hiperparâmetros, mas outra semente, a floresta sai ligeiramente diferente, e suas métricas mudam um pouco. Esse efeito costuma ser pequeno quando há muitas árvores, porque a média entre elas reduz o ruído, mas pode ser maior quando poucas árvores decidem casos próximos do limiar de decisão. Até aqui, todos os modelos foram treinados com a semente 42, que garante a reprodutibilidade. Na busca, cada combinação foi treinada com três sementes (42, 43 e 44), e as métricas foram calculadas como a média delas. Assim é possível separar o ruído do sorteio do modelo do efeito real do hiperparâmetro.

A etapa seguinte avaliou 27 combinações, cruzando três quantidades de árvores (100, 300 e 500), três profundidades máximas (sem limite, 15 e 8) e três valores de \texttt{min\_samples\_leaf} (1, 5 e 10). Cada combinação foi examinada com validação cruzada estratificada de cinco dobras no conjunto de ajuste (cerca de 65 mil linhas de treino e 16 mil de avaliação por dobra) e com as três sementes. As mesmas dobras foram usadas em todas as combinações, o que permite compará-las sob partições equivalentes. Cada combinação também foi treinada no ajuste completo e avaliada na validação separada, mas essa validação não determinou a escolha. Como o ajuste e a validação vêm do mesmo dia de captura, a busca mede o desempenho dentro desse dia e não consegue prever o comportamento em outro, como discutido na seção anterior.

A regra de seleção foi definida antes da análise dos resultados e usou somente a validação cruzada. Foram consideradas duas métricas: ROC-AUC e MCC ponderado para a proporção de classes do treino. Pela regra de um erro-padrão, são elegíveis as combinações a até um erro-padrão da melhor nas duas métricas; entre elas vence a de menor custo, medido pelo número de nós. O teste não participou da decisão.

Para reduzir o custo, cada floresta foi treinada com 500 árvores, e as primeiras 100 e 300 árvores avaliaram as configurações menores. O \textit{notebook} verificou, no scikit-learn 1.9.1, que as primeiras $k$ árvores de uma floresta maior são idênticas às de uma floresta treinada com $k$ árvores (diferença máxima de $2{,}9\times10^{-8}$ nas probabilidades).

\textbf{Custo da busca.} Foram treinadas 162 florestas de 500 árvores (9 pares profundidade $\times$ folha, 3 sementes, 5 dobras mais o ajuste completo). A busca levou 1~h 08~min de relógio; a soma dos tempos de treino foi de cerca de 1~h 06~min, ou seja, cerca de 24~s por floresta em média. O cálculo das métricas levou mais 34~s. Cada unidade concluída era gravada em disco, o que permite retomar a busca se for interrompida.

\textbf{Influência de cada hiperparâmetro.} A Figura~\ref{fig:efeitos-principais} mostra o efeito principal de cada hiperparâmetro. A ROC-AUC variou pouco entre as configurações, e as diferenças ficaram dentro do ruído (erro-padrão entre dobras). A profundidade foi o fator mais influente: 83,2\% da variação do MCC entre as configurações vem dela, e 14,1\% de \texttt{min\_samples\_leaf}. Limitar a profundidade a 8 reduziu o MCC médio da validação cruzada para 0,79185, contra 0,89576 sem limite e 0,89350 com profundidade 15. Folhas maiores também reduziram o MCC (0,88708, 0,85482 e 0,83922 para 1, 5 e 10). A amplitude do efeito da profundidade foi 6,3 vezes o ruído, e a de \texttt{min\_samples\_leaf}, 2,9 vezes. A quantidade de árvores ficou dentro do ruído (0,2 vez).

% FIGURA: influencia_efeitos_principais.png
% Pasta: artefatos_modelos/rf_binario_busca/rf_busca_20261024_175105/influencia_efeitos_principais.png
\begin{figure}
\centering
\caption{Efeito principal de cada hiperparâmetro na validação cruzada}
\label{fig:efeitos-principais}
\includegraphics[width=\textwidth]{assets/influencia_efeitos_principais.png}
\legend{Fonte: Elaborado pelo autor (2026).}
\end{figure}

A profundidade e \texttt{min\_samples\_leaf} limitam a complexidade da árvore e atuam em conjunto. A interação entre os dois respondeu por 23,6\% da variação do ROC-AUC, mas apenas 2,7\% da do MCC (Figura~\ref{fig:interacao}).

% FIGURA (opcional): influencia_interacao.png
% Pasta: artefatos_modelos/rf_binario_busca/rf_busca_20261024_175105/influencia_interacao.png
\begin{figure}
\centering
\caption{Interação entre profundidade máxima e \texttt{min\_samples\_leaf}}
\label{fig:interacao}
\includegraphics[width=0.8\textwidth]{assets/influencia_interacao.png}
\legend{Fonte: Elaborado pelo autor (2026).}
\end{figure}

Já o número de árvores quase não alterou o desempenho, mas determinou o custo: de 100 para 500 árvores, os nós totais foram de cerca de 25 mil para 125 mil, o modelo de 0,56 para 2,81~MiB, o tempo de treino de 4,8 para 23,9~s e a previsão na validação de 0,29 para 1,08~s. O número de árvores explica 97,7\% da variação do tempo de treino. Com 300 árvores, o ROC-AUC já é igual ao de 500 (0,99999 na configuração escolhida e 0,99997 sem limite de profundidade, nos dois casos), o que justifica parar em 300 (Figura~\ref{fig:convergencia}).

% FIGURA: convergencia_arvores.png
% Pasta: artefatos_modelos/rf_binario_busca/rf_busca_20261024_175105/convergencia_arvores.png
\begin{figure}
\centering
\caption{Desempenho em função do número de árvores}
\label{fig:convergencia}
\includegraphics[width=0.8\textwidth]{assets/convergencia_arvores.png}
\legend{Fonte: Elaborado pelo autor (2026).}
\end{figure}

\textbf{Configuração selecionada.} Duas combinações atenderam à regra nas duas métricas: 300 árvores, profundidade 15 e folha 1; e 500 árvores, profundidade 15 e folha 1. A primeira foi escolhida por ser a mais barata. O padrão (100 árvores, sem limite, folha 1) não era elegível: seu MCC estava dentro da margem, mas o ROC-AUC (0,99996) ficou além de um erro-padrão da melhor (limite de 0,999986). Como conferência, que não decide nada, na validação separada a escolhida ficou em 4º lugar em ROC-AUC e 5º em MCC. Na validação, o MCC da escolhida (0,94018) praticamente igualou o do padrão (0,94021). O tipo com menor recall na validação foi WebDDoS (0,9639; 83 linhas), que não aparece no teste. A Tabela~\ref{tab:config-busca} compara as configurações.

\begin{table}
\centering
\caption{Configurações do Random Forest comparadas na busca}
\label{tab:config-busca}
\footnotesize
\resizebox{\textwidth}{!}{%
\begin{tabular}{lrrrrrrrr}
\hline
Configuração & Árvores & Prof. máxima & Folha mínima & ROC-AUC (val. cruzada) & MCC ponderado (val. cruzada) & Nós médios & Tamanho do modelo & Tempo de treino* \\
\hline
Padrão         & 100 & Sem limite & 1 & 0,99996 & 0,93007 & 42.328  & 0,81~MiB & 5,29~s \\
Escolhida      & 300 & 15         & 1 & 0,99999 & 0,92818 & 113.643 & 2,23~MiB & 15,14~s \\
Outra elegível & 500 & 15         & 1 & 0,99999 & 0,92973 & 189.531 & 3,72~MiB & 25,23~s \\
\hline
\end{tabular}}
\legend{*Tempo de ajuste completo; para 100 e 300 árvores é estimado proporcionalmente ao da floresta de 500.\\
Fonte: Elaborado pelo autor (2026).}
\end{table}


\subsubsection{Avaliação da configuração ajustada no teste}
\label{sec:rf-bin-teste-ajustado}

A configuração escolhida foi treinada novamente no \textit{refit} (mesmos 101.536 registros do modelo padrão) e avaliada no mesmo teste sem sobreposição, com os limiares definidos na validação: 0,04042 e 0,41599, para metas de 1\% e 0,1\%. Na validação, esses limiares resultaram em falso alarme de 0,906\% e 0,098\%. O modelo padrão foi avaliado nas mesmas linhas. A comparação é descritiva, porque a seleção ocorreu na validação cruzada.

Com limiar 0,5, o modelo ajustado elevou a ROC-AUC de 0,98993 para 0,99136, o recall de 0,95111 para 0,95259 e o MCC de 0,35521 para 0,36144. O falso alarme caiu de 2,095\% para 1,763\%. A melhoria é modesta, e a precisão da classe BENIGN continuou baixa (0,13987): cerca de 14\% dos fluxos classificados como benignos eram de fato benignos. O ajuste de hiperparâmetros foi feito com dados do mesmo dia do treino e, por isso, não podia corrigir a diferença entre os dias: o falso alarme do ajustado foi de 0,079\% na validação e 1,763\% no teste.

A Tabela~\ref{tab:comparacao-teste} compara os modelos. Nos pontos de 1\% e 0,1\%, o falso alarme observado ficou muito acima das metas da validação (26,359\% e 3,017\% no ajustado). A configuração ajustada, portanto, não resolveu a transferência dos limiares da validação para o teste. No ponto de 0,1\%, o ajustado teve mais falsos positivos (1.696 contra 1.492), mas muito menos falsos negativos (2.016 contra 5.502), e o MCC subiu de 0,94001 para 0,96682.

\begin{table}
\centering
\caption{Comparação dos modelos no teste}
\label{tab:comparacao-teste}
\footnotesize
\resizebox{\textwidth}{!}{%
\begin{tabular}{llrrrrrr}
\hline
Limiar de referência & Modelo & Limiar aplicado & Recall de ataques & Falso alarme & MCC & Falsos negativos & Falsos positivos \\
\hline
Padrão (0,5)              & Padrão   & 0,5     & 0,95111 & 2,095\%  & 0,35521 & 350.177 & 1.178 \\
Padrão (0,5)              & Ajustado & 0,5     & 0,95259 & 1,763\%  & 0,36144 & 339.595 & 991 \\
Meta de 1\% na validação   & Padrão   & 0,05    & 1,00000 & 26,788\% & 0,85450 & 23      & 15.059 \\
Meta de 1\% na validação   & Ajustado & 0,04042 & 1,00000 & 26,359\% & 0,85709 & 16      & 14.818 \\
Meta de 0,1\% na validação & Padrão   & 0,46    & 0,99923 & 2,654\%  & 0,94001 & 5.502   & 1.492 \\
Meta de 0,1\% na validação & Ajustado & 0,41599 & 0,99972 & 3,017\%  & 0,96682 & 2.016   & 1.696 \\
\hline
\end{tabular}}
\legend{Fonte: Elaborado pelo autor (2026).}
\end{table}

% FIGURA: curvas_teste.png (notebook 08, modelo ajustado) -- logo apos a analise da Tabela de comparacao
% Pasta: artefatos_modelos/rf_binario_busca/rf_busca_20261024_175105/curvas_teste.png
\begin{figure}
\centering
\caption{Curvas ROC e de precisão $\times$ recall do BENIGN do modelo ajustado no teste}
\label{fig:curvas-teste-ajustado}
\includegraphics[width=\textwidth]{assets/curvas_teste.png}
\legend{Fonte: Elaborado pelo autor (2026).}
\end{figure}

\textbf{Desempenho por tipo de ataque.} A avaliação por tipo mostra o que a métrica global esconde. O tipo Syn, com 4.887.060 registros, teve recall de 0,9307 no limiar 0,5 e concentra cerca de 99,8\% dos 339.595 ataques perdidos (cálculo a partir do recall e do total de Syn). A metade dos fluxos Syn recebeu escore de até 0,62, próximo do limiar. Por isso, baixar o limiar corrige o Syn (recall de 0,9997 no limiar de 0,1\%), ao custo de mais falsos alarmes: a especificidade do BENIGN cai de 0,9824 para 0,9698, e para 0,7364 no limiar de 1\%. Esse comportamento é coerente com a diferença entre os dias: os fluxos Syn do dia de teste ficam mais próximos da fronteira de decisão do que os do dia de treino. O menor recall foi o de UDPLag (0,8995), mas com apenas 557 registros. O tipo Portmap não existe no treino, e mesmo assim teve recall de 0,9304 nos 2.702 registros do teste. É um resultado de uma única execução e não permite generalizar para tipos novos. A Tabela~\ref{tab:recall-tipo} apresenta os valores.

\begin{table}
\centering
\caption{Recall por classe e tipo no teste}
\label{tab:recall-tipo}
\footnotesize
\resizebox{\textwidth}{!}{%
\begin{tabular}{lrcrrr}
\hline
Classe ou tipo & Registros no teste & Presente no treino & Padrão (0,5) & Ajustado (0,5) & Ajustado (limiar de 0,1\%) \\
\hline
BENIGN (especificidade) & 56.216    & Sim & 0,9790 & 0,9824 & 0,9698 \\
LDAP                    & 112.491   & Sim & 0,9992 & 0,9992 & 0,9994 \\
MSSQL                   & 854.242   & Sim & 0,9997 & 0,9998 & 0,9999 \\
NetBIOS                 & 91.175    & Sim & 0,9993 & 0,9993 & 0,9997 \\
Portmap                 & 2.702     & Não & 0,9289 & 0,9304 & 0,9734 \\
Syn                     & 4.887.060 & Sim & 0,9285 & 0,9307 & 0,9997 \\
UDP                     & 1.214.746 & Sim & 0,9999 & 0,9999 & 1,0000 \\
UDPLag                  & 557       & Sim & 0,8977 & 0,8995 & 0,9623 \\
\hline
\end{tabular}}
\legend{Fonte: Elaborado pelo autor (2026).}
\end{table}

% FIGURA: recall_por_tipo_teste.png -- logo apos a tabela de recall por tipo
% Pasta: artefatos_modelos/rf_binario_busca/rf_busca_20261024_175105/recall_por_tipo_teste.png
\begin{figure}
\centering
\caption{Recall por tipo de ataque no teste}
\label{fig:recall-tipo}
\includegraphics[width=\textwidth]{assets/recall_por_tipo_teste.png}
\legend{Fonte: Elaborado pelo autor (2026).}
\end{figure}

\textbf{Custo computacional.} A Tabela~\ref{tab:custo-treino-tamanho-modelos} reúne o custo de treino e o tamanho dos modelos padrão e ajustado. O \textit{refit} da configuração ajustada levou 20,9~s de relógio (mediana de três repetições; 20,4 a 20,9~s) e 152,6~s de CPU, com acréscimo de RAM de até 94,6~MiB. O treino no ajuste levou 16,2~s e 112,9~s de CPU. O modelo final tem 300 árvores, 125.006 nós e 2,5~MiB. Em relação ao padrão (\textit{refit} de 7,1~s, 48,1~s de CPU, 48.162 nós, 0,9~MiB), o custo de treino é cerca de 2,9 vezes maior, e o número de nós, cerca de 2,6 vezes.

\begin{table}
\centering
\caption{Custo de treino e tamanho dos modelos RF (sem os picos de memória do processo)}
\label{tab:custo-treino-tamanho-modelos}
\footnotesize
\resizebox{\textwidth}{!}{%
\begin{tabular}{lcccc}
\hline
Item & Padrão: ajuste (81.228) & Padrão: \textit{refit} (101.536) & Ajustado: ajuste (81.228) & Ajustado: \textit{refit} (101.536) \\
\hline
Tempo de relógio (s) & 5,4 (1 execução) & 7,1 (mediana de 3; 6,0 a 8,3) & 16,2 (1 execução) & 20,9 (mediana de 3; 20,4 a 20,9) \\
Tempo de CPU (s) & 29,3 & 48,1 & 112,9 & 152,6 \\
Acréscimo de RAM no treino (MiB) & 60 & até 79 & 75 & até 94,6 \\
Árvores & 100 & 100 & 300 & 300 \\
Nós (total) & 42.468 & 48.162 & 114.544 & 125.006 \\
Profundidade média (máxima) & 20,4 (27) & 21,1 (26) & 15,0 (15) & 15,0 (15) \\
Arquivo compactado (MiB) & 0,8 & 0,9 & 2,2 & 2,5 \\
\hline
\end{tabular}}
\legend{Fonte: Elaborado pelo autor (2026).}
\end{table}

No teste, a leitura, a previsão pelos dois modelos e a gravação dos 145 lotes levaram 219~s (inclui esperas por falta de memória do servidor). O tempo de previsão foi de 52,3~s para o ajustado (138.091 fluxos/s) e 22,1~s para o padrão (327.026 fluxos/s). Para um fluxo isolado, a latência mediana fim a fim do ajustado foi de 29,693~ms com um núcleo (p95 de 54,632~ms e p99 de 106,515~ms) e 108,118~ms com todos os núcleos (p95 de 131,077~ms e p99 de 195,851~ms). Para um fluxo só, o paralelismo atrapalha, porque o custo de coordenar os núcleos supera o ganho. Em lote ocorre o inverso: a vazão fim a fim passou de 39.649 para 121.911 fluxos/s. Essas medidas incluem o pré-processamento, mas não a leitura do ClickHouse nem a extração de atributos.


\subsubsection{Estudo da variação entre sementes}
\label{sec:rf-bin-sementes}

Por fim, foi analisada a variação causada pela aleatoriedade do Random Forest, isto é, quanto o resultado no teste muda apenas porque a semente mudou. Os modelos padrão e ajustado foram treinados com cinco sementes (42 a 46), com os mesmos dados de treino e teste, de modo que só a semente mudou. A semente 42 corresponde aos modelos já avaliados, e os outros oito modelos foram treinados agora, em 2~min 06~s. A previsão dos 7.219.189 registros pelos dez modelos levou 574~s (cerca de 9,5~min, incluindo esperas por falta de memória do servidor): 19 a 21~s por modelo padrão e 44 a 45~s por modelo ajustado. Os modelos usados no estudo são os mesmos gravados anteriormente, e o ajustado retreinado com a semente 42 saiu idêntico ao do \textit{notebook} 07.

O ROC-AUC médio do ajustado foi 0,99118 (desvio-padrão de 0,00032), e o do padrão foi 0,99014 (desvio de 0,00119). Na comparação pareada, o ajustado teve vantagem média de 0,00104 e foi melhor em quatro das cinco sementes. No entanto, na semente 46 o padrão superou o ajustado, e o desvio das diferenças (0,00126) é maior que a vantagem média. A diferença de ROC-AUC, portanto, não é maior que a variação causada pelo sorteio do modelo.

O principal efeito da configuração ajustada aparece na estabilidade. No limiar 0,5, o MCC do padrão teve média de 0,46266 e desvio de 0,25217, contra 0,35942 e 0,00216 do ajustado. O padrão também apresentou uma oscilação grande do recall do Syn, de 0,9231 a 0,9982, ao passo que no ajustado esse recall ficou entre 0,9289 e 0,9309. A semente 46 foi a que corrigiu o Syn: reduziu os ataques perdidos para 9.652, e por isso o recall médio do padrão (0,95937) supera o do ajustado. Sem essa semente, a média do padrão cai para cerca de 0,9495. No limiar de 0,1\%, que foi fixado a partir do modelo original e não recalculado por semente, o padrão na semente 45 chegou a 343.008 ataques perdidos (MCC de 0,35537), enquanto o ajustado manteve o MCC entre 0,95450 e 0,97212. Os fluxos Syn do dia de teste ficam perto da fronteira de decisão (seção~\ref{sec:rf-bin-teste}), região em que pequenas diferenças entre florestas decidem para que lado eles caem. Isso explica por que a semente influencia tanto o resultado nesse tipo. A Tabela~\ref{tab:sementes} resume a variação.

\begin{table}
\centering
\caption{Variação das métricas entre cinco sementes no teste}
\label{tab:sementes}
\footnotesize
\resizebox{\textwidth}{!}{%
\begin{tabular}{lrr}
\hline
Métrica & RF padrão: média $\pm$ desvio (mín.--máx.) & RF ajustado: média $\pm$ desvio (mín.--máx.) \\
\hline
ROC-AUC                          & 0,99014 $\pm$ 0,00119 (0,98865--0,99147) & 0,99118 $\pm$ 0,00032 (0,99076--0,99145) \\
PR-AUC do BENIGN                 & 0,97345 $\pm$ 0,00616 (0,96551--0,98101) & 0,97613 $\pm$ 0,00275 (0,97221--0,97944) \\
Recall, limiar 0,5               & 0,95937 $\pm$ 0,02201 (0,94741--0,99865) & 0,95202 $\pm$ 0,00063 (0,95143--0,95277) \\
Falso alarme, limiar 0,5         & 2,094\% $\pm$ 0,275 p.p. (1,781--2,433\%) & 1,782\% $\pm$ 0,081 p.p. (1,709--1,902\%) \\
MCC, limiar 0,5                  & 0,46266 $\pm$ 0,25217 (0,34225--0,91365) & 0,35942 $\pm$ 0,00216 (0,35714--0,36186) \\
Recall do Syn, limiar 0,5        & 0,94058 $\pm$ 0,03227 (0,92307--0,99817) & 0,92982 $\pm$ 0,00092 (0,92895--0,93092) \\
MCC, limiar de 0,1\% da validação & 0,82684 $\pm$ 0,26404 (0,35537--0,96612) & 0,96420 $\pm$ 0,00649 (0,95450--0,97212) \\
\hline
\end{tabular}}
\legend{Fonte: Elaborado pelo autor (2026).}
\end{table}

% FIGURA: grafico_sementes.png -- logo apos o paragrafo da estabilidade / tabela de sementes
% Pasta: artefatos_modelos/rf_binario_busca/rf_busca_20261024_175105/estudo_sementes/grafico_sementes.png
\begin{figure}
\centering
\caption{Resultados no teste por semente, modelos padrão e ajustado}
\label{fig:sementes}
\includegraphics[width=\textwidth]{assets/grafico_sementes.png}
\legend{Fonte: Elaborado pelo autor (2026).}
\end{figure}

\textbf{Limitações metodológicas.} Com cinco sementes, os desvios e contagens são descritivos e não constituem teste estatístico. O estudo varia apenas a semente do modelo, não as linhas do ajuste ou do \textit{refit}. O teste vem de um dia de captura diferente do treino, e esse fator influencia fortemente os resultados (seção~\ref{sec:rf-bin-teste}). Como a validação e a validação cruzada usam o mesmo dia do treino, elas não detectam essa diferença, e o desempenho no teste é o único indicador de generalização entre dias. Além disso, o teste não foi usado uma única vez. Foi avaliado com o modelo padrão (seção~\ref{sec:rf-bin-teste}), reavaliado com o padrão e o ajustado (seção~\ref{sec:rf-bin-teste-ajustado}) e reutilizado nesta etapa com dez modelos. Nenhum hiperparâmetro, limiar ou semente foi escolhido com base nesses resultados, mas a decisão de realizar a busca veio depois da observação do desempenho do padrão no teste. O teste, portanto, deixa de ser uma avaliação totalmente independente a partir da seção~\ref{sec:rf-bin-teste}, e os resultados devem ser lidos como descrição, não como estimativa final de desempenho.


\subsection{Treino do SVM binário}
\label{sec:svm-bin-treino}

\textbf{Objetivo e configuração.} Esta etapa treina o SVM binário (BENIGN $\times$ ataque) e registra seu desempenho na validação e seu custo computacional. Os conjuntos de ajuste, validação, \textit{refit} e teste, os 43 atributos de entrada, o tratamento dos valores ausentes (mediana do treino e indicador de ausência), as métricas e a definição dos pontos de operação são os mesmos do Random Forest e já foram descritos nas Seções~\ref{sec:views}, \ref{sec:balanceamento} e \ref{sec:rf-bin-treino}; não são repetidos aqui. O SVM adotado usa o núcleo (\textit{kernel}) RBF, que permite fronteiras de decisão não lineares. Seus principais hiperparâmetros são $C$, que controla a penalização dos erros de classificação, e $\gamma$ (gama), que define o alcance da influência de cada registro sobre a fronteira.

Há três diferenças em relação ao RF. Primeiro, o escore do SVM não é uma probabilidade: é a distância com sinal à fronteira de decisão (\texttt{decision\_function}), e o limiar natural é zero; como os escores podem passar de 1, os limiares dos pontos de operação (falso alarme de 1\% e 0,1\%) também são calculados na validação, mas em outra escala. Segundo, o SVM exige normalização, pois compara fluxos por distância. Terceiro, o treinamento não envolve sorteio: o SVM não amostra linhas nem atributos, e três treinos repetidos do modelo final geraram modelos idênticos (2.661 vetores de suporte em todos). Por isso, não há estudo de sementes para o SVM, ao contrário do RF (Seção~\ref{sec:rf-bin-sementes}).

A configuração inicial usou $C = 1$ e $\gamma =$ \texttt{"scale"}, valor que o scikit-learn calcula a partir dos dados; após a padronização, ele foi de aproximadamente 0,02380. A Tabela~\ref{tab:svm-config-inicial} resume os parâmetros.

\begin{table}
\centering
\caption{Configuração inicial do SVM binário}
\label{tab:svm-config-inicial}
\begin{tabular}{ll}
\hline
Parâmetro & Valor \\
\hline
\textit{Kernel} & RBF \\
$C$ & 1 \\
$\gamma$ & \texttt{"scale"} (referência: 0,02380) \\
Tolerância & 0,001 \\
Normalização & \texttt{StandardScaler} nas 42 colunas numéricas \\
Cache do treino & 200 MiB \\
Processos no treino / na previsão & 1 / 2 \\
Escore e limiar natural & \texttt{decision\_function}; 0 \\
\hline
\end{tabular}
\par\smallskip
Fonte: Elaborado pelo autor (2026).
\end{table}

\textbf{Preparação.} O código conferiu as contagens dos conjuntos (ajuste de 81.228 registros, validação de 60.154, \textit{refit} de 101.536 e teste de 7.219.189) e a disponibilidade dos 43 atributos. A padronização (média 0 e desvio 1) vale só para as 42 colunas numéricas e faz parte de um \textit{Pipeline}: em cada treino, a média e o desvio são calculados somente com os dados daquele treino, de modo que a validação, as dobras de avaliação e o teste não participam desse cálculo. O indicador de ausência permanece binário (0 ou 1), pois padronizá-lo daria cerca de sete desvios a qualquer registro com valor ausente. O notebook conferiu que as colunas numéricas ficam com média 0 e desvio 1 no treino e que o indicador não é alterado. Nas caudas, 9 das 42 colunas numéricas têm $|z|$ máximo acima de 50, e 7 têm mais de 1\% dos registros acima de 5 desvios; como a média e o desvio são sensíveis a valores extremos, isso motivou testar também um normalizador robusto (a seguir).

\textbf{Calibração do custo (notebook 10).} Antes da busca, o notebook 10 mediu o custo do SVM nos dados e no computador, sem escolher hiperparâmetros. Cada treino rodou em um processo separado, com limite de tempo, e as medidas foram feitas na mesma máquina do RF (Windows 11, Intel de 10 núcleos físicos e 12 lógicos, 7,7 GiB de RAM, Python 3.13.14 e scikit-learn 1.9.1). Com $C = 1$ e $\gamma$ de referência, o tempo de treino cresceu de 0,66~s com 10 mil registros para 6,59~s com 40 mil e 341,10~s com os 81.228 registros do ajuste (Tabela~\ref{tab:svm-custo-tamanho}); os vetores de suporte passaram de 552 para 2.152. Ajustando uma lei de potência, o tempo cresce como $n^{2,88}$ e os vetores de suporte como $n^{0,65}$. O ajuste, porém, é só uma média: entre 40 mil e 81 mil registros o tempo aumentou cerca de 52 vezes, bem mais que o previsto pelo expoente. O tempo de treino desse mesmo ponto também varia entre execuções: 341,1~s nesta calibração, 264,2~s e 274,8~s nos testes de \textit{cache} descritos adiante, e 234,6~s na busca. Por isso, esses tempos servem como ordem de grandeza (alguns minutos) e não permitem comparações finas. O custo de prever foi de cerca de 202~ns por par (registro $\times$ vetor de suporte), o que permite estimar o tempo de previsão do teste.

\begin{table}
\centering
\caption{Custo do treino do SVM em função do tamanho da amostra ($C = 1$, $\gamma$ de referência)}
\label{tab:svm-custo-tamanho}
\begin{tabular}{rrrrr}
\hline
Registros & Tempo de treino (s) & CPU (s) & Vetores de suporte & Fração de vetores \\
\hline
10.000 & 0,66 & 0,66 & 552 & 5,5\% \\
20.000 & 1,71 & 1,73 & 808 & 4,0\% \\
40.000 & 6,59 & 6,50 & 1.252 & 3,1\% \\
81.228 & 341,10 & 335,30 & 2.152 & 2,6\% \\
\hline
\end{tabular}
\par\smallskip
Fonte: Elaborado pelo autor (2026).
\end{table}

% FIGURA: custo_por_tamanho.png (notebook 10)
% Pasta: artefatos_modelos/svm_binario_calibracao/svm_cal_20261024_224947/custo_por_tamanho.png
\begin{figure}
\centering
\caption{Tempo de treino e vetores de suporte do SVM conforme o tamanho da amostra}
\label{fig:svm-custo-tamanho}
\includegraphics[width=0.85\textwidth]{assets/svm_custo_por_tamanho.png}
\par\smallskip
Fonte: Elaborado pelo autor (2026).
\end{figure}

Três outras medidas da calibração orientaram o protocolo.

\begin{itemize}
\item \textbf{Normalização.} Em uma amostra de 40 mil registros, com $C = 1$ e $\gamma$ de referência, o \texttt{StandardScaler} concluiu o treino em 6,59~s (ROC-AUC de 0,99895 na validação, apenas informativo). Sem normalizar e com o normalizador robusto (mediana e intervalo interquartil), o treino ultrapassou o limite de 60~s e foi encerrado. A comparação usa uma amostra e um ponto da grade; ela mostra que a normalização é necessária e que o padrão é viável, mas não mede qualidade.
\item \textbf{Sensibilidade a $C$ e $\gamma$.} Na mesma amostra, o tempo de treino variou 2,4 vezes entre as combinações testadas (de 6,6 a 16,1~s); as mais caras foram $C$ baixo (0,1; 16,1~s) e $\gamma$ alto (4 vezes a referência; 13,4~s), e as mais baratas, $C$ alto com $\gamma$ de referência. A grade da busca foi desenhada com essa informação.
\item \textbf{Paralelismo.} O SVM do scikit-learn usa um único núcleo por treino. Treinos simultâneos exigiriam cerca de 1,6 GiB de RAM livre para dois processos, e só havia 1,3 GiB; o teste foi pulado e o treino ficou com um processo por limite de memória, não por ganho medido. Na previsão da validação em vários processos, dois processos deram ganho de 1,84 vez (eficiência de 0,92), e quatro, 1,84 vez (eficiência de 0,46); seis e oito deram 2,17 e 2,03 vezes, com eficiência de 0,36 e 0,25. Adotaram-se dois processos, o maior número com eficiência acima de 0,5, e as previsões em paralelo reproduziram exatamente os escores da previsão sequencial.
\end{itemize}

Um teste complementar verificou se o \textit{cache} de 200 MiB do treino era o gargalo. Com 1.500 e 1.000 MiB, o treino não pôde ser concluído por falta de memória; com 500 MiB levou 274,8~s, e com 200 MiB, 264,2~s. Como o \textit{cache} maior não reduziu o tempo, manteve-se 200 MiB. As medições de calibração levaram, no total, cerca de 13 min (custo $\times$ tamanho, 7 min 05 s; sensibilidade, 2 min 49 s; normalizadores, 2 min 00 s; paralelismo, 1 min 23 s). Com elas, o notebook extrapolou 1 h 02 min para a validação cruzada de cinco dobras e 31 min 36 s para prever o teste com dois processos; a execução real foi mais rápida que essas estimativas (Seções seguintes).

\textbf{Desempenho inicial na validação.} A configuração inicial foi ajustada nos 81.228 registros do ajuste e avaliada nos 60.154 da validação. No limiar natural (zero), obteve ROC-AUC de 0,99930, recall de 0,99406, taxa de falso alarme de 0,423\% e MCC de 0,70985 com a ponderação que recompõe a proporção do treino (a mesma definida para o RF). Os limiares para falso alarme de 1\% e 0,1\% escolhidos na validação foram $-0{,}8443$ e 1,0031; com eles, o recall foi 0,99542 e 0,92490, e o falso alarme obtido, 0,975\% e 0,098\% (Tabela~\ref{tab:svm-pontos-operacao}, adiante). O recall caiu de 0,99542 para 0,92490 ao se passar de 1\% para 0,1\%. Como no RF, os limiares foram escolhidos e avaliados na mesma validação e, por isso, são estimativas de desenvolvimento, possivelmente otimistas; a taxa de 0,1\% corresponde a cerca de dez falsos alarmes entre os 10.154 benignos.

\textbf{Custo do treino inicial.} O ajuste da configuração inicial nos 81.228 registros levou 234,6~s na busca (Tabela~\ref{tab:svm-config-busca}) e o \textit{refit} com 101.536 registros, 257,4~s, com 256,6~s de CPU, ou seja, cerca de um núcleo ocupado. O modelo final inicial tem 2.539 vetores de suporte e ocupa 0,26 MiB.

\textbf{Verificação.} Foram aprovadas as oito checagens do notebook 10: amostras estratificadas (50\% de benignos em cada tamanho), normalização sem vazamento, vetores de suporte sempre menores que o número de linhas, ROC-AUC entre 0 e 1, previsão e treinos em paralelo idênticos aos sequenciais, estimativas finitas e arquivos gravados.


\subsection{Busca de hiperparâmetros do SVM binário}
\label{sec:svm-bin-busca}

\textbf{Protocolo.} O notebook 11 avaliou nove combinações: $C \in \{1, 10, 100\}$ e $\gamma$ igual a 0,25, 1 ou 4 vezes o valor de referência (0,00595, 0,02380 e 0,09521). A configuração inicial ($C = 1$, $\gamma$ de referência) está na grade. Foi usada validação cruzada estratificada de cinco dobras no ajuste, com as mesmas linhas, a mesma ordem e a mesma semente de divisão da busca do RF (Seção~\ref{sec:rf-bin-busca}); as dobras têm 16.245 ou 16.246 registros de avaliação, e o código registra uma impressão digital delas (\texttt{2defef26}) para conferência. A normalização é refeita em cada treino. Como o SVM não sorteia, cada configuração tem cinco treinos de dobra e um treino completo, e não 15 e três, como no RF. A regra de seleção é a do RF: ROC-AUC e MCC ponderado para a proporção do treino (peso do benigno de 0,00619 no ajuste), regra de um erro-padrão e, entre as elegíveis, a configuração de menor custo; aqui o custo é o número de vetores de suporte (que define o tamanho do modelo e o tempo de prever), com desempate pelo menor tempo de treino nas dobras e depois pelo menor $C$.

\textbf{Custo da busca.} Foram realizados 54 ajustes (nove configurações $\times$ cinco dobras, mais um ajuste completo por configuração), em um processo, com limite de 2 h por treino; nenhum atingiu o limite. A busca levou 43 min 05 s de relógio, e a soma dos tempos de treino foi de 34 min 27 s (cerca de 38 s por treino em média). O cálculo das métricas levou mais 5 s. Cada treino era gravado ao terminar, o que permite retomar a busca. O tempo médio de treino nas dobras variou de 14,6~s ($C = 100$, $\gamma = 1\times$) a 81,1~s ($C = 1$, $\gamma = 4\times$); os maiores valores de $\gamma$ foram os mais caros. O número de vetores de suporte importa para a inferência, pois a previsão compara cada registro com todos eles.

\textbf{Influência dos hiperparâmetros.} A ROC-AUC média na validação cruzada ficou entre 0,99844 e 0,99917 nas nove configurações, e o MCC ponderado, entre 0,55497 e 0,70720; o número médio de vetores de suporte nas dobras variou de 1.163,6 a 2.459,6. As amplitudes dos efeitos de $C$ e de $\gamma$ ficaram acima do ruído (erro-padrão entre dobras): no MCC, 0,0864 para $C$ (16,0 vezes o ruído) e 0,0728 para $\gamma$ (13,5 vezes); na ROC-AUC, 0,0003 nos dois (2,6 vezes), uma diferença pequena em valor absoluto. Do MCC, 52,8\% da variação entre as configurações vem de $C$, 38,8\% de $\gamma$ e 8,4\% da interação; na ROC-AUC, a interação responde por 37,7\%. Aumentar $C$ de 1 para 100 elevou o MCC em 0,10045 (16,3 erros-padrão), e reduzir $\gamma$ para 0,25 vez o valor de referência o reduziu em 0,05178. Ou seja, o resultado depende da combinação de $C$ e $\gamma$, não de cada um isoladamente.

% FIGURA: influencia_efeitos_principais.png (notebook 11)
% Pasta: artefatos_modelos/svm_binario_busca/svm_busca_20261025_004828/influencia_efeitos_principais.png
\begin{figure}
\centering
\caption{Efeito dos hiperparâmetros do SVM na validação cruzada}
\label{fig:svm-efeitos}
\includegraphics[width=\textwidth]{assets/svm_influencia_efeitos_principais.png}
\par\smallskip
Fonte: Elaborado pelo autor (2026).
\end{figure}

% FIGURA: influencia_interacao_corrigida.png (notebook 11; existe tambem a influencia_interacao.png original)
% Pasta: artefatos_modelos/svm_binario_busca/svm_busca_20261025_004828/influencia_interacao_corrigida.png
\begin{figure}
\centering
\caption{Interação entre $C$ e $\gamma$ na busca do SVM}
\label{fig:svm-interacao}
\includegraphics[width=0.85\textwidth]{assets/svm_influencia_interacao.png}
\par\smallskip
Fonte: Elaborado pelo autor (2026).
\end{figure}

\textbf{Configuração selecionada.} Só uma combinação foi elegível nas duas métricas: $C = 100$ e $\gamma = 0{,}09521$ (quatro vezes o valor de referência). Ela teve a maior ROC-AUC (0,99917; limite de empate de 0,999031, com três combinações empatadas) e ficou a menos de um erro-padrão da melhor em MCC (0,70093 contra 0,70720 de $C = 100$ e $\gamma$ de referência; limite de 0,70027, com duas combinações empatadas). Na validação cruzada, teve 2.021 vetores de suporte em média nas dobras, e 2.310 no ajuste completo. A combinação de melhor MCC ($C = 100$ e $\gamma = 1\times$) é mais barata (1.314 vetores nas dobras, 1.563 no ajuste completo, 0,15 MiB e 24,0~s de treino completo), mas ficou um pouco abaixo do limite de ROC-AUC (0,99897); portanto, a escolha decorreu da regra e não do menor custo. Como conferência, que não decide nada, na validação separada a escolhida ficou em 1º lugar em ROC-AUC e em 2º em MCC. A Tabela~\ref{tab:svm-config-busca} compara as configurações. O tipo com menor recall na validação, no limiar zero, foi o WebDDoS (0,6747; 83 registros; 0,6224 na validação cruzada), muito abaixo dos demais, todos acima de 0,987; esse tipo não existe no teste.

\begin{table}
\centering
\caption{Configurações do SVM comparadas na busca}
\label{tab:svm-config-busca}
\footnotesize
\resizebox{\textwidth}{!}{%
\begin{tabular}{lrrrrrrr}
\hline
Configuração & $C$ & $\gamma$ & ROC-AUC (val. cruzada) & MCC ponderado (val. cruzada) & Vetores nas dobras (média) & Vetores no ajuste completo & Tempo do ajuste completo \\
\hline
Padrão (inicial)        & 1   & 0,02380 & 0,99889 & 0,60675 & 1.822 & 2.152 & 234,6~s \\
Selecionada             & 100 & 0,09521 & 0,99917 & 0,70093 & 2.021 & 2.310 & 81,8~s \\
Melhor MCC (não elegível) & 100 & 0,02380 & 0,99897 & 0,70720 & 1.314 & 1.563 & 24,0~s \\
\hline
\end{tabular}}
\par\smallskip
Fonte: Elaborado pelo autor (2026).
\end{table}

\textbf{Modelo selecionado na validação.} O modelo selecionado, treinado no ajuste inteiro, foi avaliado na validação (Tabela~\ref{tab:svm-metricas-validacao}). No limiar zero, identificou 49.811 dos 50.000 ataques e classificou como ataque 23 dos 10.154 benignos (189 ataques perdidos e 10.131 benignos reconhecidos). O MCC passa de 0,988 sem ponderação para 0,785 com a proporção do treino, e a precisão da classe BENIGN cai de 0,982 para 0,620, o que reforça a leitura feita para o RF: as métricas dos ataques ficam próximas de 1, e as da classe BENIGN e o MCC são as que revelam os erros.

\begin{table}
\centering
\caption{Métricas do SVM selecionado na validação (limiar 0)}
\label{tab:svm-metricas-validacao}
\resizebox{0.9\textwidth}{!}{%
\begin{tabular}{lccc}
\hline
Métrica & Sem pesos & Proporção real & Depende da proporção? \\
\hline
ROC-AUC & 0,99948 & 0,99948 & não \\
Recall (ataques) & 0,99622 & 0,99622 & não \\
Taxa de falso alarme & 0,00227 & 0,00227 & não \\
Acurácia balanceada & 0,99698 & 0,99698 & não \\
Acurácia & 0,99648 & 0,99623 & sim \\
Precisão (ataques) & 0,99954 & 0,99999 & sim \\
F1 (ataques) & 0,99788 & 0,99810 & sim \\
MCC & 0,98757 & 0,78513 & sim \\
PR-AUC (ataques) & 0,99986 & 1,00000 & sim \\
Precisão (BENIGN) & 0,98169 & 0,62019 & sim \\
F1 (BENIGN) & 0,98965 & 0,76491 & sim \\
F1 macro & 0,99376 & 0,88150 & sim \\
PR-AUC (BENIGN) & 0,99676 & 0,91871 & sim \\
\hline
\end{tabular}}
\par\smallskip
Fonte: Elaborado pelo autor (2026).
\end{table}

\textbf{Pontos de operação.} A Tabela~\ref{tab:svm-pontos-operacao} mostra os limiares escolhidos na validação para a configuração inicial e para a selecionada. Valem as mesmas ressalvas do RF: os limiares foram escolhidos e avaliados na mesma validação, e a taxa de 0,1\% depende de poucos falsos alarmes. Para a meta de 0,1\%, a configuração selecionada manteve recall de 0,98882, contra 0,92490 da inicial.

\begin{table}
\centering
\caption{Pontos de operação do SVM na validação}
\label{tab:svm-pontos-operacao}
\begin{tabular}{llrrrr}
\hline
Configuração & Meta de falso alarme & Limiar & Recall & Falso alarme obtido & MCC ponderado \\
\hline
Inicial & 1\% & $-0{,}8443$ & 0,99542 & 0,975\% & 0,75096 \\
Inicial & 0,1\% & 1,0031 & 0,92490 & 0,098\% & 0,26504 \\
Selecionada & 1\% & $-0{,}7444$ & 0,99724 & 0,936\% & 0,82527 \\
Selecionada & 0,1\% & 0,9911 & 0,98882 & 0,098\% & 0,59301 \\
\hline
\end{tabular}
\par\smallskip
Fonte: Elaborado pelo autor (2026).
\end{table}

\textbf{Estabilidade do resultado.} Como no RF (Seção~\ref{sec:rf-bin-treino}), a validação foi dividida em cinco partições, com o mesmo modelo, e fez-se validação cruzada de cinco dobras no ajuste (Tabela~\ref{tab:svm-dispersao}). Cada modelo das dobras foi treinado com cerca de 64.982 registros. Nas partições, o recall variou de 0,9957 a 0,9968 e a ROC-AUC, de 0,99909 a 0,99979; essa variação descreve subconjuntos de avaliação, não modelos treinados com dados diferentes. Na validação cruzada, a ROC-AUC média foi 0,99917 (desvio de 0,00031) e o MCC ponderado, 0,70093 (desvio de 0,00810). As duas colunas da tabela usam pesos diferentes para recompor a proporção do treino (0,0305 por benigno na validação e 0,00619 no ajuste), portanto seus MCCs ponderados não são comparáveis entre si.

\begin{table}
\centering
\caption{Dispersão das métricas do SVM selecionado (média $\pm$ desvio-padrão; mín--máx)}
\label{tab:svm-dispersao}
\resizebox{\textwidth}{!}{%
\begin{tabular}{lcc}
\hline
Métrica & 5 partições da validação & 5 dobras da validação cruzada \\
\hline
Recall & 0,99622 $\pm$ 0,00054 (0,9957--0,9968) & 0,99372 $\pm$ 0,00029 (0,99323--0,99397) \\
Falso alarme & 0,00227 $\pm$ 0,00056 (0,00148--0,00296) & 0,00271 $\pm$ 0,00068 (0,00172--0,00345) \\
ROC-AUC & 0,99948 $\pm$ 0,00027 (0,99909--0,99979) & 0,99917 $\pm$ 0,00031 (0,99865--0,99944) \\
MCC sem pesos & 0,98757 $\pm$ 0,00129 (0,98653--0,98913) & 0,99102 $\pm$ 0,00075 (0,99016--0,99201) \\
MCC ponderado & 0,78582 $\pm$ 0,02138 (0,76563--0,80920) & 0,70093 $\pm$ 0,00810 (0,68699--0,70755) \\
\hline
\end{tabular}}
\par\smallskip
Fonte: Elaborado pelo autor (2026).
\end{table}

% FIGURA: svm_fig_cv_ajuste_5dobras.png (gerada por graficos_svm.py)
% Pasta: artefatos_modelos/svm_binario_busca/svm_busca_20261025_004828/svm_fig_cv_ajuste_5dobras.png
\begin{figure}
\centering
\caption{Validação cruzada no ajuste do SVM (5 dobras)}
\label{fig:svm-validacao-cruzada-ajuste}
\includegraphics[width=\textwidth]{assets/svm_fig_cv_ajuste_5dobras.png}
\par\smallskip
Fonte: Elaborado pelo autor (2026).
\end{figure}

\textbf{\textit{Refit} (modelo final).} O treino final repete a configuração selecionada no mesmo conjunto de 101.536 registros usado no RF (Seção~\ref{sec:rf-bin-treino}), com a normalização recalculada nesse conjunto (o notebook conferiu que a média e o desvio gravados são os do \textit{refit}). Como parte da validação passou a integrar o treino, o modelo final não é medido nela, e sua qualidade aparece só no teste. Para a comparação no teste, também foi feito um \textit{refit} da configuração inicial. Os três \textit{refits} do ajustado e o do inicial levaram, juntos, 10 min 53 s.

\textbf{Custo computacional.} A Tabela~\ref{tab:svm-custo-treino} reúne o custo de treino e o tamanho dos modelos. O treino do SVM ocupa cerca de um núcleo (o tempo de CPU é igual ao de relógio), diferentemente do RF, que usa vários. O modelo ajustado tem 2.661 vetores de suporte no \textit{refit} (1.824 e 837 nas duas classes) e ocupa 0,37 MiB. O ajuste foi uma única execução; em vista da variação de tempo entre execuções (Seção~\ref{sec:svm-bin-treino}), diferenças pequenas de tempo não são conclusivas.

\begin{table}
\centering
\caption{Custo de treino e tamanho dos modelos SVM (sem os picos de memória do processo)}
\label{tab:svm-custo-treino}
\resizebox{\textwidth}{!}{%
\begin{tabular}{lcccc}
\hline
Item & Padrão: ajuste (81.228) & Padrão: \textit{refit} (101.536) & Ajustado: ajuste (81.228) & Ajustado: \textit{refit} (101.536) \\
\hline
Tempo de relógio (s) & 234,6 (1 execução) & 257,4 (1 execução) & 81,8 (1 execução) & 126,8 (mediana de 3; 123,2 a 129,8) \\
Tempo de CPU (s) & 233,2 & 256,6 & 81,6 & 126,6 \\
Acréscimo de RAM (MiB) & 299 & $\approx$ 311\textsuperscript{*} & 298 & 312 a 329 \\
Vetores de suporte & 2.152 & 2.539 & 2.310 & 2.661 \\
Arquivo compactado (MiB) & 0,23 & 0,26 & 0,33 & 0,37 \\
\hline
\end{tabular}}
\par\smallskip
* Estimado, não medido: o notebook gravou só o pico de RAM do \textit{refit} do padrão (492,4~MiB). O acréscimo foi obtido subtraindo o consumo do processo antes do treino, de cerca de 181~MiB, valor observado nas três repetições do \textit{refit} do modelo ajustado (180,6 a 181,1~MiB). Fonte: Elaborado pelo autor (2026).
\end{table}

Na inferência do modelo final, medida na validação só para velocidade, a latência de um fluxo, em 1.000 fluxos e com um processo, foi de 1,508~ms (mediana), 2,033~ms (p95) e 2,815~ms (p99) de ponta a ponta. A vazão em lote, medida em 20 mil registros (o SVM é lento demais para usar a validação inteira em cada repetição), foi de 2.160 fluxos/s com um processo e 4.558 com dois (Tabela~\ref{tab:svm-inferencia}). O pré-processamento acrescenta cerca de 0,04~ms por fluxo. A leitura do banco não entra, e esses números medem velocidade, não qualidade.

\begin{table}
\centering
\caption{Inferência do modelo final do SVM (fim a fim)}
\label{tab:svm-inferencia}
\begin{tabular}{lcc}
\hline
Medida & 1 processo & 2 processos \\
\hline
Latência p50 / p95 / p99 (ms) & 1,508 / 2,033 / 2,815 & não medida \\
Vazão em lote (fluxos/s) & 2.160 & 4.558 \\
\hline
\end{tabular}
\par\smallskip
Fonte: Elaborado pelo autor (2026).
\end{table}

\textbf{Verificação.} As 23 checagens do notebook 11 foram aprovadas. Entre elas: 54 treinos concluídos, dobras que cobrem cada linha do ajuste uma vez, normalização sem vazamento, modelo gravado que reproduz os escores, três \textit{refits} idênticos, mesmas linhas do \textit{refit} do RF, escolha reproduzível a partir do CSV gravado e pesos que preservam a proporção do treino.


\subsection{Avaliação do SVM binário no teste}
\label{sec:svm-bin-teste}

\textbf{Preparação e composição do teste.} A avaliação usou o mesmo conjunto de teste do RF, com os mesmos 7.219.189 registros (56.216 benignos e 7.162.973 ataques; 0,779\% de benignos), sem balanceamento e sem pesos, processado em 145 lotes de até 50 mil registros. A deduplicação em relação ao treino, a composição e as ressalvas sobre esse conjunto estão na Seção~\ref{sec:rf-bin-teste}. O tipo Portmap, ausente do treino, aparece com 2.702 registros. Como no RF, o teste foi mantido separado do treino e da escolha de hiperparâmetros, e os limiares usados são os gravados na validação (para o SVM, o limiar natural é zero). O notebook foi executado uma única vez, com bloqueio de nova execução depois de gravado o resultado.

\textbf{Resultados dos modelos padrão e ajustado.} O notebook 12 avaliou o modelo padrão (inicial: $C = 1$, $\gamma = 0{,}02380$) e o ajustado ($C = 100$, $\gamma = 0{,}09521$) nos mesmos registros, cada um com seus limiares da validação. A Tabela~\ref{tab:svm-teste} apresenta os três pontos de operação.

\begin{table}
\centering
\caption{Comparação dos modelos SVM no teste}
\label{tab:svm-teste}
\footnotesize
\resizebox{\textwidth}{!}{%
\begin{tabular}{llrrrrrr}
\hline
Ponto de operação & Modelo & Limiar & Recall & Falso alarme & MCC & Falsos negativos & Falsos positivos \\
\hline
Natural & Padrão & 0 & 0,74595 & 0,911\% & 0,14770 & 1.819.788 & 512 \\
Natural & Ajustado & 0 & 0,99642 & 1,160\% & 0,82075 & 25.657 & 652 \\
Meta de 1\% na validação & Padrão & $-0{,}8443$ & 0,99965 & 2,400\% & 0,96575 & 2.517 & 1.349 \\
Meta de 1\% na validação & Ajustado & $-0{,}7444$ & 0,99970 & 3,239\% & 0,96469 & 2.129 & 1.821 \\
Meta de 0,1\% na validação & Padrão & 1,0031 & 0,36210 & 0,021\% & 0,06630 & 4.569.245 & 12 \\
Meta de 0,1\% na validação & Ajustado & 0,9911 & 0,59328 & 0,569\% & 0,10496 & 2.913.331 & 320 \\
\hline
\end{tabular}}
\par\smallskip
Fonte: Elaborado pelo autor (2026).
\end{table}

No limiar natural, o ajustado teve ROC-AUC de 0,99531 (0,99504 no padrão), recall de 0,99642 e MCC de 0,82075. Em relação ao padrão, reduziu os ataques perdidos de 1.819.788 para 25.657, embora o falso alarme tenha subido de 0,911\% para 1,160\%. A busca, portanto, melhorou a posição da fronteira no limiar natural, sem eliminar o compromisso entre detectar ataques e evitar alarmes falsos. Como no RF, as taxas da validação não se mantiveram no teste: para a meta de 1\%, o ajustado teve 3,239\% de falso alarme (recall de 0,99970), e para a meta de 0,1\%, 0,569\% de falso alarme, com recall de apenas 0,59328. Esse resultado é coerente com a diferença entre os dias de captura do treino e do teste discutida na Seção~\ref{sec:rf-bin-teste}: a validação é um sorteio do dia de treino e não anuncia o comportamento em outro dia. No SVM, o efeito foi menor que no RF: o falso alarme no limiar natural foi de 0,227\% na validação e 1,160\% no teste (cerca de cinco vezes), e o recall se manteve (0,99622 na validação e 0,99642 no teste).

% FIGURA: curvas_teste.png (notebook 12)
% Pasta: artefatos_modelos/svm_binario_busca/svm_busca_20261025_004828/curvas_teste.png
\begin{figure}
\centering
\caption{Curvas ROC e de precisão $\times$ recall da classe BENIGN no teste (SVM)}
\label{fig:svm-curvas-teste}
\includegraphics[width=\textwidth]{assets/svm_curvas_teste.png}
\par\smallskip
Fonte: Elaborado pelo autor (2026).
\end{figure}

\textbf{Desempenho por tipo de ataque.} O resultado global é dominado pelo Syn, que tem 4.887.060 registros (cerca de 68\% dos ataques). No limiar natural, o ajustado alcançou recall de 0,9949 no Syn, contra 0,6278 do padrão, o que explica a diferença nos ataques perdidos. No ponto de 0,1\%, o recall do Syn cai para 0,4041 no ajustado e 0,0822 no padrão, e é isso que derruba o recall global nesse ponto. Para o Portmap, ausente do treino, o recall foi 0,9567 no ajustado (resultado de uma única execução, que não permite generalizar para tipos novos). O menor recall do ajustado foi o do UDPLag, 0,8815, mas esse tipo tem só 557 registros e deve ser lido com cautela. A Tabela~\ref{tab:svm-recall-tipo} apresenta os valores.

\begin{table}
\centering
\caption{Recall por classe e tipo no teste para o SVM}
\label{tab:svm-recall-tipo}
\footnotesize
\resizebox{\textwidth}{!}{%
\begin{tabular}{lrcrrr}
\hline
Classe ou tipo & Registros no teste & Presente no treino & Padrão (0) & Ajustado (0) & Ajustado (limiar de 0,1\%) \\
\hline
BENIGN (especificidade) & 56.216 & Sim & 0,9909 & 0,9884 & 0,9943 \\
LDAP & 112.491 & Sim & 0,9992 & 0,9994 & 0,9990 \\
MSSQL & 854.242 & Sim & 0,9997 & 0,9997 & 0,9994 \\
NetBIOS & 91.175 & Sim & 0,9992 & 0,9994 & 0,9969 \\
Portmap & 2.702 & Não & 0,9341 & 0,9567 & 0,9378 \\
Syn & 4.887.060 & Sim & 0,6278 & 0,9949 & 0,4041 \\
UDP & 1.214.746 & Sim & 0,9999 & 1,0000 & 0,9999 \\
UDPLag & 557 & Sim & 0,8115 & 0,8815 & 0,8618 \\
\hline
\end{tabular}}
\par\smallskip
Fonte: Elaborado pelo autor (2026).
\end{table}

% FIGURA: recall_por_tipo_teste.png (notebook 12)
% Pasta: artefatos_modelos/svm_binario_busca/svm_busca_20261025_004828/recall_por_tipo_teste.png
\begin{figure}
\centering
\caption{Recall por tipo de ataque no teste (SVM)}
\label{fig:svm-recall-tipo}
\includegraphics[width=\textwidth]{assets/svm_recall_por_tipo_teste.png}
\par\smallskip
Fonte: Elaborado pelo autor (2026).
\end{figure}

\textbf{Custo computacional.} A previsão dos 7.219.189 registros levou 1.227,6~s para o modelo ajustado e 1.174,5~s para o padrão, ambos com dois processos, o que corresponde a 5.881 e 6.147 fluxos por segundo (2.661 e 2.539 vetores de suporte). O tempo total de relógio do notebook 12 foi de 2.643~s (44 min 03 s), dos quais 2.402~s são de previsão dos dois modelos; o restante (cerca de 4 min) é leitura no ClickHouse, gravação e esperas por falta de memória do servidor. Em 1.000 fluxos individuais, a latência de ponta a ponta do modelo ajustado foi de 1,410~ms (mediana) com um processo, 1,723~ms (p95) e 2,376~ms (p99); a vazão em lote, medida em 20 mil fluxos, foi de 2.520 fluxos/s com um processo e 5.701 com dois. O pré-processamento acrescenta cerca de 0,04~ms. A leitura do banco e a extração de atributos do tráfego não entram. O modelo ajustado ocupa 0,37 MiB, e o padrão, 0,26 MiB. A previsão real foi mais rápida que a estimada na calibração (31 min 36 s por modelo com dois processos; real de 20 min 28 s para o ajustado).

\textbf{Limitações.} Os limiares foram escolhidos na validação e, como mostrado, não se mantêm em outro dia de captura. O teste foi processado uma única vez para o SVM, mas é o mesmo conjunto já utilizado na avaliação do RF (Seções~\ref{sec:rf-bin-teste}, \ref{sec:rf-bin-teste-ajustado} e \ref{sec:rf-bin-sementes}); nenhuma escolha do SVM dependeu desses resultados. Os tempos do SVM variam entre execuções (Seção~\ref{sec:svm-bin-treino}). O SVM não tem variação por semente, mas, como o RF, não foi avaliado com outras divisões de ajuste e \textit{refit}.


\subsection{Comparação entre Random Forest e SVM binários}
\label{sec:comparacao-rf-svm-bin}

\textbf{Desempenho no teste.} A comparação considera os modelos ajustados, avaliados sobre os mesmos 7.219.189 registros. No limiar natural de cada algoritmo (0,5 para o RF e zero para o SVM), o SVM apresentou ROC-AUC superior (0,99531 contra 0,99136), maior recall (0,99642 contra 0,95259), menor taxa de falso alarme (1,160\% contra 1,763\%) e maior MCC (0,82075 contra 0,36144). O RF teve PR-AUC da classe BENIGN ligeiramente superior: 0,97473 contra 0,96823.

A comparação muda conforme o ponto de operação. Com os limiares definidos na validação para a meta de 1\%, os dois modelos tiveram recall próximo de 1, mas o falso alarme no teste foi menor para o SVM: 3,239\%, contra 26,359\% no RF. Na meta de 0,1\%, o RF foi superior em recall e MCC: 0,99972 e 0,96682, contra 0,59328 e 0,10496 no SVM. Em ambos os algoritmos, o falso alarme no teste ficou acima da meta da validação. Em pontos lidos na própria curva do teste (limiar escolhido no teste, apenas informativos), o SVM obteve recall de 0,73909 com falso alarme de 1\% e 0,42572 com 0,1\%, contra 0,56489 e 0,31771 do RF. A Tabela~\ref{tab:comparacao-rf-svm} resume esses resultados.

\begin{table}
\centering
\caption{Comparação dos modelos ajustados no teste}
\label{tab:comparacao-rf-svm}
\begin{tabular}{llrrrr}
\hline
Ponto de operação & Modelo & ROC-AUC & Recall & Falso alarme & MCC \\
\hline
Limiar natural & RF & 0,99136 & 0,95259 & 1,763\% & 0,36144 \\
Limiar natural & SVM & 0,99531 & 0,99642 & 1,160\% & 0,82075 \\
Meta de 1\% na validação & RF & 0,99136 & 1,00000 & 26,359\% & 0,85709 \\
Meta de 1\% na validação & SVM & 0,99531 & 0,99970 & 3,239\% & 0,96469 \\
Meta de 0,1\% na validação & RF & 0,99136 & 0,99972 & 3,017\% & 0,96682 \\
Meta de 0,1\% na validação & SVM & 0,99531 & 0,59328 & 0,569\% & 0,10496 \\
\hline
\end{tabular}
\par\smallskip
Fonte: Elaborado pelo autor (2026).
\end{table}

\textbf{Custo e conclusão.} O ajuste final do SVM levou 126,8~s, contra 20,9~s do \textit{refit} do RF ajustado, e o SVM usa um núcleo, enquanto o RF usa todos. No teste completo, a previsão do SVM levou 1.227,6~s (dois processos) e a do RF, 52,3~s (todos os núcleos): 5.881 contra 138.091 fluxos por segundo. Nas configurações medidas, portanto, o RF teve vazão em lote muito maior. Em contrapartida, a latência mediana de ponta a ponta para um fluxo isolado foi menor no SVM: 1,410~ms, contra 29,693~ms do RF ajustado (e 7,660~ms do RF padrão), todos com um núcleo. O arquivo do SVM também é menor (0,37 MiB contra 2,5 MiB). Uma diferença adicional é a estabilidade: o desempenho do RF varia com a semente, sobretudo no tipo Syn (Seção~\ref{sec:rf-bin-sementes}), enquanto o SVM não tem esse sorteio.

Não há um vencedor único para todos os critérios. No teste, o SVM ajustado apresentou melhor discriminação e melhor resultado no limiar natural, além de menor latência por fluxo isolado. O RF foi mais rápido em grandes lotes e muito superior no ponto definido para falso alarme de 0,1\%. A escolha depende, assim, do requisito operacional: o SVM é favorecido quando importam o desempenho no limiar natural e a resposta por fluxo; o RF, quando a prioridade é a vazão em lote ou o ponto de operação de 0,1\%. Essa conclusão é descritiva: os limiares foram definidos na validação, o teste vem de outro dia de captura, os resultados do teste foram observados em uma única execução, e os modelos usaram estratégias de paralelismo diferentes na previsão.

\subsubsection{Visões de entrada dos modelos multiclasse}
\label{sec:visoes-modelos-multiclasse}

Com base nas tabelas \texttt{prep\_final\_treino} e \texttt{prep\_final\_teste}, foram criadas duas visões de entrada comuns aos modelos multiclasse, o Random Forest e o SVM, no mesmo formato das visões dos modelos binários (Seção~\ref{sec:visoes-modelos-binario}): uma para o treino e outra para o teste sem sobreposição com o treino. A variável resposta \texttt{y} identifica sete classes: BENIGN e seis tipos de ataque presentes tanto no treino quanto no teste. Como no caso binário, foram aplicadas a ablação de atributos, a imputação de valores ausentes e a remoção, do teste, das linhas sobrepostas ao treino. A decisão específica do multiclasse é a definição das classes, e a seleção de atributos e os indicadores de ausência diferem em alguns pontos, descritos a seguir.

\paragraph{Classes e códigos.}
Para que uma classe pudesse ser aprendida e avaliada, ela precisava estar presente nos dois conjuntos. A comparação foi feita pelo nome sem o prefixo \texttt{DrDoS\_}, como no pré-processamento. Foram mantidas BENIGN, LDAP, MSSQL, NetBIOS, Syn, UDP e UDPLag. O rótulo \texttt{y} recebeu um código inteiro: 0 para BENIGN e, em ordem alfabética, de 1 a 6 para os tipos de ataque (LDAP = 1, MSSQL = 2, NetBIOS = 3, Syn = 4, UDP = 5 e UDPLag = 6). A tabela de auditoria \texttt{prep\_classes\_rf\_multi} registra a correspondência entre códigos e nomes.

Foram excluídos os tipos presentes apenas no treino (DNS, NTP, SNMP, SSDP, TFTP e WebDDoS) e o Portmap, presente apenas no teste. Como não há exemplos de Portmap no treino, o modelo não poderia aprender essa classe; as classes exclusivas do treino também não poderiam ser avaliadas no teste. A análise multiclasse considera, portanto, as sete classes compartilhadas. A exclusão dos seis tipos retirou 6.630.065 registros do treino (DNS 113.259, NTP 1.112.936, SNMP 115.316, SSDP 891.221, TFTP 4.396.919 e WebDDoS 414), que passou de 8.257.343 para 1.627.278 registros. Com isso, a proporção de BENIGN no treino subiu de 0,61\% no binário para 3,12\%.

\paragraph{Atributos de entrada.}
Da seleção de 45 atributos do modelo multiclasse, foram retirados por ablação \texttt{bwd\_packets\_s}, \texttt{inbound} e \texttt{total\_backward\_packets}, as mesmas três colunas excluídas no caso binário (Seção~\ref{sec:visoes-modelos-binario}), com a lista definida antes da avaliação dos modelos. Restaram 42 atributos numéricos.

A diferença em relação ao binário está no número de atributos marcados como suspeitos: no pré-processamento, 17 atributos apresentaram força de AUC $\geq 0{,}80$ para separar uma classe das demais, e a ablação retirou apenas três deles, mantendo os outros 14. Esse critério, aplicado a uma classe contra todas as outras, pode destacar características próprias de um tipo de tráfego (por exemplo, a duração do fluxo em um protocolo específico), e por si só não demonstra que a coluna seja um atalho indevido. Retirar todos os 17 deixaria o modelo sem alguns de seus atributos mais informativos e impediria que a comparação com e sem ablação dissesse algo sobre atalhos. Por isso, manteve-se a mesma lista do binário.

As seleções binária e multiclasse têm 42 atributos numéricos cada, mas não são idênticas: diferem em três atributos. Quando dois atributos são quase idênticos ($|r|$ alto), o pré-processamento mantém o de maior nota, e essa nota depende da tarefa. O multiclasse mantém \texttt{flow\_packets\_s}, \texttt{bwd\_packet\_length\_mean} e \texttt{active\_mean}, enquanto o binário mantém \texttt{fwd\_packets\_s}, \texttt{avg\_bwd\_segment\_size} e \texttt{active\_min}. Somados os dois indicadores de ausência descritos a seguir, o modelo recebe 44 colunas de entrada, listadas no Apêndice A.

\paragraph{Valores ausentes.}
Dois atributos têm indicador de ausência: \texttt{flow\_packets\_s} e \texttt{flow\_bytes\_s} (no binário, só \texttt{flow\_bytes\_s}, pois \texttt{flow\_packets\_s} foi descartada na seleção por ser gêmea de \texttt{fwd\_packets\_s}). Os valores ausentes foram substituídos pelas medianas do treino, 1,3374 para \texttt{flow\_packets\_s} e 689,5911 para \texttt{flow\_bytes\_s}, as mesmas para treino e teste, sem uso de dados de teste. Essas medianas vêm do pré-processamento (tabela \texttt{prep\_parametros}) e não foram recalculadas para as sete classes; por isso a de \texttt{flow\_bytes\_s} coincide com a do binário (689,59). Os indicadores \texttt{flow\_packets\_s\_ausente} e \texttt{flow\_bytes\_s\_ausente} preservam a informação de que o valor original estava ausente. Como as visões são comuns, o Random Forest e o SVM recebem os valores já imputados.

\paragraph{Padronização dos atributos (SVM).}
O Random Forest usa os atributos na escala original. O SVM multiclasse recebe os 42 atributos numéricos padronizados por z-score dentro do mesmo \textit{pipeline}, com média e desvio-padrão estimados só no treino de cada ajuste, e os dois indicadores permanecem em 0 e 1, como descrito para o SVM binário (Seção~\ref{sec:visoes-modelos-binario}).

\paragraph{Teste sem sobreposição.}
O critério é o mesmo do binário (Seção~\ref{sec:visoes-modelos-binario}): foram retiradas do teste as linhas cujos atributos coincidem com os de alguma linha do treino, e o treino foi mantido completo, com 1.627.278 registros das sete classes. Dos 20.177.565 registros dessas classes no teste oficial, 7.216.487 permaneceram e 12.961.078 foram removidos, como mostra a Tabela~\ref{tab:registros-rf-multiclasse}.

\begin{table}
    \centering
    \caption{Registros por classe nas visões multiclasse}
    \label{tab:registros-rf-multiclasse}
    \footnotesize
    \resizebox{\textwidth}{!}{%
    \begin{tabular}{lrrrrrr}
        \hline
        Classe & Treino & Teste oficial* & Teste sem sobreposição & Removidos & \% removido & \% do teste sem sobreposição \\
        \hline
        BENIGN & 50.768 & 56.965 & 56.216 & 749 & 1,3 & 0,779 \\
        LDAP & 30.099 & 1.915.122 & 112.491 & 1.802.631 & 94,1 & 1,559 \\
        MSSQL & 207.003 & 5.787.453 & 854.242 & 4.933.211 & 85,2 & 11,837 \\
        NetBIOS & 19.490 & 3.657.497 & 91.175 & 3.566.322 & 97,5 & 1,263 \\
        Syn & 155.515 & 4.891.500 & 4.887.060 & 4.440 & 0,1 & 67,721 \\
        UDP & 1.075.337 & 3.867.155 & 1.214.746 & 2.652.409 & 68,6 & 16,833 \\
        UDPLag & 89.066 & 1.873 & 557 & 1.316 & 70,3 & 0,008 \\
        Total & 1.627.278 & 20.177.565 & 7.216.487 & 12.961.078 & 64,2 & 100,000 \\
        \hline
    \end{tabular}}
    \legend{*Soma apenas das sete classes; o teste oficial completo tem 20.364.525 registros, incluindo 186.960 de Portmap.\\
    Fonte: Elaborado pelo autor (2026).}
\end{table}

Os totais diferem dos do binário (Tabela~\ref{tab:registros-rf-binario}: 7.219.189 registros e 13.145.336 removidos) porque, no binário, o Portmap permanece como ataque. A diferença de 2.702 registros no teste sem sobreposição corresponde às linhas de Portmap que não se sobrepõem ao treino, conforme a nota daquela tabela.

A remoção não afetou todas as classes igualmente: foram removidos 94,1\% dos registros de LDAP e 97,5\% dos de NetBIOS, enquanto apenas 0,1\% dos de Syn. Com isso, a composição do teste mudou muito: o Syn passou de cerca de 24,2\% do teste oficial (sete classes) para 67,721\% do teste sem sobreposição, e o UDPLag ficou com 0,008\%. Essa diferença é relevante para interpretar os resultados por classe, e a baixa quantidade de UDPLag (557 registros) torna as métricas dessa classe menos estáveis.

\paragraph{Validação.}
As visões foram aprovadas nas 16 checagens definidas. Foram verificadas as colunas esperadas (\texttt{id\_linha}, \texttt{y}, os 42 atributos e os dois indicadores, em 46 colunas), a ausência de rótulos e de atributos retirados na ablação, o tipo simples de \texttt{y} (inteiro sem sinal de 8 bits), os códigos e as contagens das sete classes nos dois conjuntos e a inexistência de valores nulos ou infinitos. Também foi confirmada a equivalência com a imputação do pré-processamento em amostras do treino (1.281 registros, dos quais 185 tinham valor ausente imputado) e do teste (8.524 registros, dos quais 4.081 tinham valor ausente imputado).

A tabela \texttt{prep\_views\_rf\_multi} mantém a auditoria dos 45 atributos da seleção multiclasse, identificando os usados e os retirados por ablação, e a \texttt{prep\_classes\_rf\_multi} registra os códigos das classes e suas contagens. Como o treino selecionado tem mais de 1,6 milhão de registros, o \textit{notebook} também prevê sua leitura em lotes, e a estimativa de memória para carregá-lo inteiro em \texttt{float64} foi de cerca de 0,53 GiB.

\paragraph{Apêndice A -- Colunas de entrada dos modelos multiclasse (44).}

\texttt{avg\_fwd\_segment\_size}, \texttt{total\_length\_of\_fwd\_packets}, \texttt{fwd\_packet\_length\_max}, \texttt{flow\_packets\_s}, \texttt{act\_data\_pkt\_fwd}, \texttt{flow\_iat\_mean}, \texttt{max\_packet\_length}, \texttt{flow\_bytes\_s}, \texttt{flow\_duration}, \texttt{total\_fwd\_packets}, \texttt{flow\_iat\_std}, \texttt{flow\_iat\_max}, \texttt{bwd\_header\_length}, \texttt{bwd\_iat\_total}, \texttt{bwd\_iat\_max}, \texttt{bwd\_iat\_mean}, \texttt{ack\_flag\_count}, \texttt{min\_seg\_size\_forward}, \texttt{bwd\_iat\_min}, \texttt{init\_win\_bytes\_forward}, \texttt{fwd\_packet\_length\_std}, \texttt{flow\_iat\_min}, \texttt{fwd\_header\_length}, \texttt{packet\_length\_variance}, \texttt{packet\_length\_std}, \texttt{down\_up\_ratio}, \texttt{idle\_std}, \texttt{protocol}, \texttt{idle\_min}, \texttt{active\_mean}, \texttt{active\_max}, \texttt{total\_length\_of\_bwd\_packets}, \texttt{bwd\_packet\_length\_max}, \texttt{init\_win\_bytes\_backward}, \texttt{bwd\_packet\_length\_mean}, \texttt{active\_std}, \texttt{urg\_flag\_count}, \texttt{bwd\_packet\_length\_min}, \texttt{cwe\_flag\_count}, \texttt{bwd\_iat\_std}, \texttt{bwd\_packet\_length\_std}, \texttt{fwd\_psh\_flags}, \texttt{flow\_packets\_s\_ausente} (indicador de ausência) e \texttt{flow\_bytes\_s\_ausente} (indicador de ausência).

\subsection{Balanceamento do treino multiclasse}
\label{sec:balanceamento-multiclasse}
 
\textbf{Conjunto de partida e necessidade de balanceamento.} Após o pré-processamento e a definição das classes (Seção~\ref{sec:visoes-modelos-multiclasse}), o treino continha 1.627.278 registros das sete classes presentes tanto no treino quanto no teste. A distribuição era desigual: UDP tinha 1.075.337 registros e NetBIOS, 19.490, uma razão de cerca de 55 para 1. O desequilíbrio é bem menor que o do caso binário (162 para 1, Seção~\ref{sec:balanceamento}), porque BENIGN passou a representar 3,12\% do treino, mas ainda faria o aprendizado favorecer as classes mais numerosas. Pelo mesmo motivo apresentado no caso binário, a acurácia isolada não é adequada para avaliar o classificador, e o desempenho será examinado também por classe.
 
\textbf{Limite por classe e ausência de dados sintéticos.} Como no binário, a técnica escolhida foi a subamostragem das classes maiores, sem criar registros artificiais. Foi definido um limite de 50.000 registros reais por classe. Esse valor é próximo à quantidade de BENIGN disponível (50.768), de modo que a classe benigna é aproveitada em 98,5\%. As classes com menos registros que o limite foram mantidas integralmente: LDAP, com 30.099, e NetBIOS, com 19.490. Nas demais, foram selecionados 50.000 registros. Ao todo, 299.589 registros (18,4\% do treino) foram selecionados, e a razão entre a maior e a menor classe no ajuste caiu para aproximadamente 2,57 para 1. Além de reduzir o desequilíbrio, o limite mantém o conjunto em tamanho viável para o SVM, cujo custo de treinamento cresce pelo menos quadraticamente com o número de amostras (Seção~\ref{sec:balanceamento}).
 
A sobreamostragem não foi adotada. Para chegar a 50.000 registros, seria necessário criar artificialmente cerca de 19.901 exemplos de LDAP e 30.510 de NetBIOS, aproximadamente 40\% e 61\% do total dessas classes, respectivamente, uma parcela grande demais para dados sintéticos. Como a diferença após a seleção é moderada, optou-se por compensá-la com pesos de classe durante o treinamento, o que preserva exemplos reais.
 
\textbf{Formação do ajuste e da validação.} De cada classe, foram selecionados até 50.000 registros, dos quais 80\% foram destinados ao ajuste e 20\% à validação. A divisão foi feita separadamente dentro de cada classe, ordenando os registros por um número pseudoaleatório determinístico (função de dispersão aplicada ao identificador da linha e à semente fixa 20261018), de modo que a mesma seleção é obtida a cada execução. Assim, ajuste e validação não compartilham registros e ambos contêm as sete classes. O ajuste ficou com 239.671 registros e a validação, com 59.918. A Tabela~\ref{tab:balanceamento-multiclasse} apresenta a distribuição por classe.
 
\begin{table}
\centering
\caption{Distribuição dos registros por classe no balanceamento multiclasse}
\label{tab:balanceamento-multiclasse}
\resizebox{\textwidth}{!}{%
\begin{tabular}{lrrrrr}
\hline
Classe & Treino & Selecionados & Ajuste & Validação & Não utilizados \\
\hline
BENIGN & 50.768 & 50.000 & 40.000 & 10.000 & 768 \\
LDAP & 30.099 & 30.099 & 24.079 & 6.020 & 0 \\
MSSQL & 207.003 & 50.000 & 40.000 & 10.000 & 157.003 \\
NetBIOS & 19.490 & 19.490 & 15.592 & 3.898 & 0 \\
Syn & 155.515 & 50.000 & 40.000 & 10.000 & 105.515 \\
UDP & 1.075.337 & 50.000 & 40.000 & 10.000 & 1.025.337 \\
UDPLag & 89.066 & 50.000 & 40.000 & 10.000 & 39.066 \\
\textbf{Total} & \textbf{1.627.278} & \textbf{299.589} & \textbf{239.671} & \textbf{59.918} & \textbf{1.327.689} \\
\hline
\end{tabular}%
}
\par\smallskip
Fonte: Elaborado pelo autor (2026).
\end{table}
 
\textbf{Pesos de classe e ponderação da validação.} Como LDAP e NetBIOS têm menos registros que as demais classes no ajuste, os modelos usam \texttt{class\_weight="balanced"}, calculado com as quantidades do ajuste: cada classe recebe o peso $\frac{239.671}{7 \times n_{\text{classe}}}$, em que $n_{\text{classe}}$ é o número de registros dela no ajuste. Cada registro de NetBIOS recebe peso de aproximadamente 2,196, cada registro de LDAP, 1,422, e os das classes com 40.000 registros, 0,856. O notebook apenas exibe esses valores; a aplicação ocorre no treinamento dos modelos, que usam o mesmo parâmetro no Random Forest e no SVM.
 
Esses pesos são diferentes dos atribuídos às linhas da validação, que têm outra finalidade. Na validação, cada registro recebe o peso (registros da classe no treino) $\div$ (registros da classe na validação), de modo que os pesos de cada classe somam sua quantidade no treino. Com eles, precisão, F1, MCC e PR-AUC, que dependem da proporção entre as classes, passam a refletir a composição do treino (por exemplo, UDP com cerca de 66\%), como no binário (Seção~\ref{sec:balanceamento}). A seleção e a comparação dos modelos, porém, usam a validação sem pesos, na qual as classes são mais equilibradas. Recall e taxa de falso alarme por classe não dependem da ponderação. A Tabela~\ref{tab:pesos-multiclasse} reúne os dois conjuntos de pesos.
 
\begin{table}
\centering
\caption{Pesos calculados para o ajuste e para a validação}
\label{tab:pesos-multiclasse}
\begin{tabular}{lrr}
\hline
Classe & Peso por registro na validação & Peso de classe no ajuste \\
\hline
BENIGN & 5,08 & 0,856 \\
LDAP & 5,00 & 1,422 \\
MSSQL & 20,70 & 0,856 \\
NetBIOS & 5,00 & 2,196 \\
Syn & 15,55 & 0,856 \\
UDP & 107,53 & 0,856 \\
UDPLag & 8,91 & 0,856 \\
\hline
\end{tabular}
\par\smallskip
Fonte: Elaborado pelo autor (2026).
\end{table}
 
\textbf{Teste.} Como no binário, o teste não foi balanceado e permanece com a composição do teste sem sobreposição ao treino, apresentada na Tabela~\ref{tab:registros-rf-multiclasse}: 7.216.487 registros, dos quais Syn representa 67,721\% e UDPLag, apenas 557 (0,008\%).
 
\textbf{Papel do ajuste, da validação e do refit.} O ajuste é o conjunto com o qual os modelos aprendem, e a validação fica separada para medir o desempenho em dados não vistos, conforme descrito para o caso binário (Seção~\ref{sec:balanceamento}). Aqui, a validação tem até 10.000 registros por classe (6.020 de LDAP e 3.898 de NetBIOS), quantidade suficiente para estimar o recall de cada classe separadamente. O \textit{refit} também segue o binário: depois das análises, o modelo final é refeito com todos os registros selecionados, ou seja, o ajuste mais a validação (299.589 registros, coluna ``Selecionados'' da Tabela~\ref{tab:balanceamento-multiclasse}), e os pesos de classe são recalculados com esse conjunto, resultando nos mesmos valores da Tabela~\ref{tab:pesos-multiclasse}. O mesmo conjunto de refit é usado pelo Random Forest e pelo SVM, o que torna a comparação justa. Depois do \textit{refit}, a validação deixa de ser independente.
 
\textbf{Armazenamento e verificações.} A seleção foi registrada no ClickHouse na tabela \texttt{prep\_bal\_multi\_selecao}, que preserva o papel de cada registro (ajuste ou validação), a classe, o código \texttt{y} e o peso. Os dados selecionados foram copiados para \texttt{prep\_bal\_multi\_dados}, e dela foram criadas as visões \texttt{prep\_vw\_bal\_multi\_treino} e \texttt{prep\_vw\_bal\_multi\_validacao}. O nome da classe e o papel do registro não aparecem como atributos; a visão de validação inclui a coluna de peso, que não é usada no treino. O ajuste ocupa cerca de 80,5 MiB em memória (\texttt{float64}, 44 colunas de entrada), de modo que é lido por inteiro.
 
O notebook aprovou 15 verificações. Entre elas, foram conferidos os totais por classe e por conjunto, o uso do limite de 50.000 registros (ou de todos os disponíveis, nas classes menores), a divisão 80/20, a ausência de linhas compartilhadas entre ajuste e validação, a correspondência dos códigos com as classes, a soma dos pesos da validação igual ao tamanho de cada classe no treino e a manutenção das 7.216.487 linhas do teste. Também se verificou, em uma amostra de 14.932 linhas, a equivalência dos dados copiados com a visão de origem, além das colunas das visões e a ausência de valores nulos nos atributos.
 
Por fim, 33.322 linhas do ajuste (13,9\%) estão marcadas com conflito de rótulo. Diferentemente do binário, em que o conflito é entre BENIGN e ataque, aqui a marca indica registros cujas características são iguais às de registros com outro rótulo multiclasse. Ela foi calculada no pré-processamento sobre todo o treino, antes da exclusão dos seis tipos, e a lista dos grupos em conflito mostra que muitos envolvem DNS, SNMP e SSDP, tipos que não fazem parte das sete classes; portanto, parte desses registros conflita com classes fora da análise. Dessas 33.322 linhas marcadas pelas 13 classes, 15.896 (47,7\%) são conflitos apenas entre as sete classes do modelo (Seção~\ref{sec:conflitos-multiclasse}). As linhas foram contabilizadas, mas não removidas, para permitir uma análise de sensibilidade posterior. Como registros idênticos com rótulos diferentes não podem ser todos classificados corretamente, os percentuais da Tabela~\ref{tab:conflitos-ajuste-multiclasse}, concentrados em NetBIOS e LDAP, devem ser considerados ao interpretar o desempenho dessas classes.
 
\begin{table}
\centering
\caption{Registros do ajuste marcados com conflito de rótulo}
\label{tab:conflitos-ajuste-multiclasse}
\begin{tabular}{lrrr}
\hline
Classe & Registros no ajuste & Registros em conflito & Percentual em conflito \\
\hline
BENIGN & 40.000 & 17 & 0,0\% \\
LDAP & 24.079 & 10.323 & 42,9\% \\
MSSQL & 40.000 & 3.784 & 9,5\% \\
NetBIOS & 15.592 & 8.149 & 52,3\% \\
Syn & 40.000 & 647 & 1,6\% \\
UDP & 40.000 & 5.559 & 13,9\% \\
UDPLag & 40.000 & 4.843 & 12,1\% \\
\textbf{Total} & \textbf{239.671} & \textbf{33.322} & 13,9\% \\
\hline
\end{tabular}
\par\smallskip
Fonte: Elaborado pelo autor (2026).
\end{table}

\subsection{Treino do Random Forest multiclasse}
\label{sec:rf-multi-treino}
 
\textbf{Objetivo e configuração.} Esta etapa treina o Random Forest (RF) multiclasse, que classifica cada fluxo em uma de sete classes (BENIGN, LDAP, MSSQL, NetBIOS, Syn, UDP e UDPLag), com uma configuração inicial fixa, mede o desempenho na validação e registra o custo computacional. A busca de hiperparâmetros e o teste vêm adiante. Os conceitos do RF, o papel de cada hiperparâmetro e a justificativa da configuração inicial foram apresentados no caso binário (Seção~\ref{sec:rf-bin-treino}) e não são repetidos. A configuração é a mesma, com uma diferença: o modelo recebe \texttt{class\_weight="balanced"}, que multiplica o peso de cada registro pelo inverso da frequência da sua classe no ajuste (Tabela~\ref{tab:pesos-multiclasse}), para que LDAP e NetBIOS, com menos registros, não sejam sacrificados. Outra diferença é a decisão: a classe prevista é a de maior probabilidade estimada (\textit{argmax}), e a probabilidade de cada classe é a média das probabilidades estimadas pelas 100 árvores. Não há limiar de decisão a escolher, e por isso o multiclasse não tem os pontos de operação do binário.
 
\begin{table}
\centering
\caption{Configuração inicial do Random Forest multiclasse}
\label{tab:rf-multi-config-inicial}
\resizebox{\ifdim\width>\textwidth\textwidth\else\width\fi}{!}{%
\begin{tabular}{ll}
\hline
Hiperparâmetro & Valor \\
\hline
\texttt{n\_estimators} (número de árvores) & 100 \\
\texttt{max\_depth} (profundidade máxima) & sem limite \\
\texttt{min\_samples\_split} & 2 \\
\texttt{min\_samples\_leaf} & 1 \\
\texttt{max\_features} & raiz quadrada do nº de atributos \\
\texttt{class\_weight} & \texttt{"balanced"} (0,856; 1,422; 2,196) \\
\texttt{random\_state} & 42 \\
\texttt{n\_jobs} & $-1$ (todos os núcleos) \\
Decisão & classe de maior probabilidade (\textit{argmax}) \\
\hline
\end{tabular}}
\par\smallskip
Fonte: Elaborado pelo autor (2026).
\end{table}
 
\textbf{Preparação.} O código conferiu as contagens do balanceamento (Seção~\ref{sec:balanceamento-multiclasse}): 239.671 registros no ajuste e 59.918 na validação. A validação não é balanceada: tem 10.000 registros por classe, exceto LDAP (6.020) e NetBIOS (3.898), que têm menos fluxos disponíveis. O modelo recebe 44 atributos: 42 numéricos e 2 indicadores de ausência (Seção~\ref{sec:visoes-modelos-multiclasse}). Os valores ausentes já haviam sido substituídos pela mediana do treino. Para conferir essa preparação, o notebook reconstruiu os fluxos da validação antes do tratamento: 1.207 dos 59.918 registros tinham algum valor ausente. Depois, comparou os atributos resultantes com os das visões e obteve resultados idênticos.
 
\textbf{Métricas de avaliação.} As métricas do binário (Seção~\ref{sec:rf-bin-treino}) foram adaptadas às sete classes. O recall e a precisão de cada classe são calculados com essa classe como positiva e as outras seis como negativas. A acurácia balanceada é a média dos recalls das sete classes. O F1 macro é a média dos sete F1, e o F1 ponderado os pondera pelo número de registros de cada classe. A ROC-AUC é calculada de cada classe contra as demais (\textit{one-vs-rest}, OvR) e tirada a média das sete. O MCC multiclasse generaliza o do binário a partir da matriz de confusão. Não dependem da proporção entre as classes: o recall por classe, a acurácia balanceada e a ROC-AUC por classe. Dependem dela: a acurácia, a precisão, o F1 e o MCC. A ROC-AUC macro depende pouco, porque a composição do \textit{resto} muda com os pesos. Na validação, as métricas sensíveis à proporção são calculadas sem pesos e com pesos que recompõem a proporção do treino deduplicado (Tabela~\ref{tab:pesos-multiclasse}); no teste, que já mantém a proporção observada, são calculadas sem pesos. A acurácia isolada é especialmente enganosa aqui, pois uma classe numerosa pode dominar o total de acertos (no teste, o Syn tem 67,7\% dos registros). Por isso as conclusões apoiam-se na acurácia balanceada, no F1 macro, no MCC e no recall e na precisão de cada classe.
 
\textbf{Treino e avaliação na validação.} O RF foi ajustado somente com os 239.671 registros do ajuste e avaliado nos 59.918 da validação. A Tabela~\ref{tab:rf-multi-validacao} apresenta as métricas gerais e a Tabela~\ref{tab:rf-multi-validacao-classe}, as de cada classe. O recall e a acurácia balanceada são iguais nas duas colunas. As métricas que dependem da proporção mudam: o MCC passa de 0,8969 para 0,8520, e a precisão macro, de 0,8992 para 0,8040. A mudança é maior nas classes menores no treino: com a proporção do treino, a precisão do UDPLag cai de 0,871 para 0,453 e a do NetBIOS, de 0,723 para 0,380. Isso ocorre porque os erros vindos de classes grandes no treino, como UDP, passam a pesar mais. Os recalls mais baixos foram os do NetBIOS (0,7630) e do UDPLag (0,8379), e os erros mais frequentes ocorreram entre pares de classes: UDPLag previsto como UDP (1.074 registros, 10,7\% do UDPLag), MSSQL previsto como NetBIOS (999, 10,0\% do MSSQL), NetBIOS previsto como MSSQL (831, 21,3\% do NetBIOS) e UDP previsto como UDPLag (771, 7,7\% do UDP). Esses mesmos pares dominam os conflitos de rótulo entre as classes (Seção~\ref{sec:conflitos-multiclasse}).
 
\begin{table}
\centering
\caption{Métricas do RF multiclasse na validação}
\label{tab:rf-multi-validacao}
\resizebox{\ifdim\width>\textwidth\textwidth\else\width\fi}{!}{%
\begin{tabular}{lccc}
\hline
Métrica & Sem pesos & Proporção do treino & Depende da proporção? \\
\hline
ROC-AUC OvR (macro) & 0,9902 & 0,9888 & pouco \\
Acurácia balanceada & 0,9026 & 0,9026 & não \\
Acurácia & 0,9126 & 0,9139 & sim \\
Precisão macro & 0,8992 & 0,8040 & sim \\
F1 macro & 0,9007 & 0,8352 & sim \\
F1 ponderado & 0,9127 & 0,9249 & sim \\
MCC & 0,8969 & 0,8520 & sim \\
\hline
\end{tabular}}
\par\smallskip
Fonte: Elaborado pelo autor (2026).
\end{table}
 
\begin{table}
\centering
\caption{Desempenho do RF multiclasse padrão por classe na validação}
\label{tab:rf-multi-validacao-classe}
\resizebox{\textwidth}{!}{%
\begin{tabular}{lrrrrr}
\hline
Classe & Registros & Recall & Precisão (sem pesos) & Precisão (proporção do treino) & ROC-AUC OvR \\
\hline
BENIGN  & 10.000 & 0,9990 & 0,9987 & 0,9974 & 1,0000 \\
LDAP    & 6.020  & 0,9646 & 0,9595 & 0,8740 & 0,9983 \\
MSSQL   & 10.000 & 0,8804 & 0,8986 & 0,9628 & 0,9863 \\
NetBIOS & 3.898  & 0,7630 & 0,7234 & 0,3804 & 0,9792 \\
Syn     & 10.000 & 0,9550 & 0,9502 & 0,9704 & 0,9965 \\
UDP     & 10.000 & 0,9180 & 0,8926 & 0,9900 & 0,9914 \\
UDPLag  & 10.000 & 0,8379 & 0,8710 & 0,4533 & 0,9798 \\
\hline
\end{tabular}}
\par\smallskip
Fonte: Elaborado pelo autor (2026).
\end{table}
 
\textbf{Estabilidade do resultado.} Como no binário, a validação foi dividida em cinco partições disjuntas e estratificadas (11.984 ou 11.983 registros cada), e o mesmo modelo, sem novo treinamento, foi avaliado em cada uma. Essa análise mostra a variação causada pela composição dos dados de avaliação, não por treinar modelos diferentes. Como complemento, foi feita validação cruzada estratificada de cinco dobras no ajuste: cada modelo foi treinado com 191.736 ou 191.737 registros (cerca de quatro quintos do ajuste) e avaliado em 47.934 ou 47.935. Essas dobras são as mesmas usadas na busca de hiperparâmetros e no SVM (impressão digital \texttt{7e97c47f}), o que permite comparar os modelos dobra a dobra. A Tabela~\ref{tab:rf-multi-dispersao} e a Figura~\ref{fig:rf-multi-cv-ajuste} mostram que as métricas variam pouco. Os valores da validação cruzada ficaram próximos dos obtidos com o modelo treinado em todo o ajuste, como esperado: cada dobra usa menos registros no treino, mas a diferença é pequena (MCC de 0,8979 contra 0,8969). O tempo de treino por dobra foi, em média, 42,2~s (de 34,1 a 45,8~s).
 
\begin{table}
\centering
\caption{Dispersão das métricas do RF multiclasse (média $\pm$ desvio-padrão; mín--máx)}
\label{tab:rf-multi-dispersao}
\resizebox{\textwidth}{!}{%
\begin{tabular}{lcc}
\hline
Métrica & 5 partições da validação & 5 dobras da validação cruzada \\
\hline
ROC-AUC OvR & 0,99020 $\pm$ 0,00091 (0,98892--0,99118) & 0,99037 $\pm$ 0,00039 (0,98995--0,99080) \\
Acurácia balanceada & 0,90255 $\pm$ 0,00312 (0,89793--0,90582) & 0,90336 $\pm$ 0,00135 (0,90161--0,90482) \\
F1 macro & 0,90068 $\pm$ 0,00264 (0,89705--0,90395) & 0,90148 $\pm$ 0,00138 (0,89960--0,90313) \\
MCC (sem pesos) & 0,89690 $\pm$ 0,00215 (0,89390--0,89981) & 0,89793 $\pm$ 0,00140 (0,89609--0,89937) \\
MCC (proporção do treino) & 0,85203 $\pm$ 0,00358 (0,84612--0,85556) & não calculado* \\
\hline
\end{tabular}}
\par\smallskip
*A métrica com a proporção do treino só foi calculada na validação (Tabela~\ref{tab:pesos-multiclasse}); na validação cruzada, as métricas foram calculadas sem pesos.\\
Fonte: Elaborado pelo autor (2026).
\end{table}
 
\begin{figure}
    \centering
    \caption{Validação cruzada no ajuste do RF multiclasse (5 dobras)}
    \label{fig:rf-multi-cv-ajuste}
    \includegraphics[width=\textwidth]{assets/rf_multi_cv_ajuste_5dobras.png}
    \par\smallskip
    Fonte: Elaborado pelo autor (2026).
\end{figure}
 
\textbf{Refit (modelo final).} O treino final usa a mesma configuração com os 299.589 registros do \textit{refit}, que reúnem o ajuste e a validação (Seção~\ref{sec:balanceamento-multiclasse}), com pesos recalculados nesse conjunto (praticamente iguais aos do ajuste). O treino foi repetido três vezes para medir o tempo (adiante). Como parte da validação entrou no treino, o modelo final não é medido nela; sua qualidade aparece só no teste.
 
\textbf{Custo computacional.} O treino foi medido no mesmo computador do caso binário (Seção~\ref{sec:rf-bin-treino}), considerando somente o ajuste do modelo (\textit{fit}). O treino no ajuste levou 41,3~s de relógio e 258,7~s de CPU, o que equivale a cerca de 6,3 dos 12 núcleos lógicos ocupados em média, com acréscimo de memória de 542~MiB (pico de 1.401~MiB). O \textit{refit} levou mediana de 70,0~s (mín. 55,1~s, máx. 74,2~s), 372,3~s de CPU e acréscimo de até 562~MiB. O tempo do ajuste vem de uma única execução, e a variação entre as três repetições do \textit{refit} (55 a 74~s) é do tamanho da diferença entre os conjuntos; por isso os tempos não permitem concluir como o custo cresce com o número de registros. O modelo final tem 4.288.578 nós (cerca de 42.886 por árvore), profundidade média de 65,1 (máxima de 96) e ocupa 67,1~MiB em arquivo compactado. O modelo do ajuste tem 3.483.364 nós, profundidade média de 61,7 (máxima de 74) e 54,7~MiB. Em relação ao RF binário (48.162 nós e 0,9~MiB no modelo final), o multiclasse tem cerca de 89 vezes mais nós e 75 vezes o tamanho em arquivo (cálculo a partir dos valores das duas seções): as árvores são muito mais profundas e numerosas em nós, o que reflete o maior número de registros (299.589 contra 101.536) e de classes. A Tabela~\ref{tab:rf-multi-custo-treino-tamanho}, na seção da configuração ajustada, compara esses valores com os do modelo ajustado.
 
Na inferência, a latência mediana de um único fluxo, fim a fim, foi de 9,6~ms com 1 núcleo (p99 de 20,4~ms). Com todos os núcleos piorou, para 48,7~ms (p99 de 72,4~ms), pois o custo de dividir o trabalho supera o ganho para um fluxo isolado. Em lote ocorre o inverso: com todos os núcleos, 60.684 fluxos/s contra 33.993 com 1 núcleo (Tabela~\ref{tab:rf-multi-inferencia}). A latência é medida em 1.000 fluxos e a vazão classifica a validação inteira (59.918 fluxos). A leitura do banco não entra, e esses números medem velocidade, não qualidade.
 
\begin{table}
\centering
\caption{Inferência do RF multiclasse final (fim a fim)}
\label{tab:rf-multi-inferencia}
\begin{tabular}{lcc}
\hline
Medida & \texttt{n\_jobs = 1} & \texttt{n\_jobs = $-1$} \\
\hline
Latência p50 / p95 / p99 (ms) & 9,6 / 17,6 / 20,4 & 48,7 / 59,7 / 72,4 \\
Vazão em lote (fluxos/s) & 33.993 & 60.684 \\
\hline
\end{tabular}
\par\smallskip
Fonte: Elaborado pelo autor (2026).
\end{table}
 
\textbf{Verificação.} As 17 checagens do notebook foram aprovadas.
 
 
\subsubsection{Avaliação do Random Forest multiclasse padrão no teste}
\label{sec:rf-multi-teste}
 
\textbf{Preparação e composição do teste.} O modelo final do RF (o treinado no \textit{refit}) foi avaliado no teste, que não foi usado para treinar, definir a configuração nem calcular as medianas do pré-processamento. As garantias e a justificativa da deduplicação em relação ao treino são as do caso binário (Seção~\ref{sec:rf-bin-teste}): o teste foi processado uma única vez, com bloqueio de nova execução, e o notebook conferiu os hiperparâmetros, a execução de origem e a ordem dos 44 atributos. A composição está na Tabela~\ref{tab:registros-rf-multiclasse}: 7.216.487 registros das sete classes, mantida a proporção observada, sem balanceamento. O Syn concentra 67,7\% dos registros e o UDP, 16,8\%; o UDPLag tem apenas 557 (0,008\%). Há também 2.702 registros Portmap, classe que não aparece no treino. Como o classificador só pode escolher entre as sete classes conhecidas, esses registros ficam fora das métricas principais e são analisados à parte (Seção~\ref{sec:comparacao-rf-svm-multi}). Ao contrário do binário, em que o Portmap era um ataque a detectar, aqui ele não tem classe correta possível.
 
\textbf{Execução e resultado do modelo padrão.} A execução completa do teste levou 248,2~s, incluindo leitura dos dados, inferência e gravação. O tempo de previsão foi de 61,2~s (117.989 fluxos por segundo), sem contar a leitura no ClickHouse nem a extração de atributos. No teste, o RF padrão obteve ROC-AUC OvR de 0,9133, acurácia balanceada de 0,6908, acurácia de 0,3917, F1 macro de 0,5626, F1 ponderado de 0,5108 e MCC de 0,4324 (a Tabela~\ref{tab:rf-multi-teste-classe} traz o resultado por classe). A distância entre acurácia (0,39) e acurácia balanceada (0,69) tem uma causa concreta: o Syn, com 67,7\% do teste, teve recall de apenas 0,2445, e 75,4\% dos seus registros (3.684.883) foram classificados como UDPLag. Outros erros relevantes foram MSSQL previsto como NetBIOS (508.850 registros, 59,6\% do MSSQL) e UDP previsto como UDPLag (141.741, 11,7\%). O UDPLag recebeu 3.827.324 previsões, das quais 96,3\% vieram de registros Syn; por isso sua precisão é 0,0001, embora o recall seja 0,4040. A precisão alta de Syn (0,9999) e a baixa de NetBIOS (0,1480) mostram o efeito contrário: o NetBIOS é previsto para muitos registros de outras classes (sobretudo MSSQL). Entre os ataques, apenas 4.078 foram previstos como BENIGN, e agrupando as seis classes de ataque, o recall é 0,99943, o falso alarme (BENIGN previsto como ataque) é 1,76\% e o MCC ataque $\times$ BENIGN é 0,95613. Isto é, o modelo separa bem tráfego benigno de ataque, e seus erros estão principalmente na identificação do \textit{tipo} de ataque.
 
\begin{table}
\centering
\caption{Desempenho do RF multiclasse padrão por classe no teste}
\label{tab:rf-multi-teste-classe}
\begin{tabular}{lrrrrr}
\hline
Classe & Registros & \% do teste & Recall & Precisão & F1 \\
\hline
BENIGN  & 56.216    & 0,78  & 0,9824 & 0,9312 & 0,9561 \\
LDAP    & 112.491   & 1,56  & 0,9632 & 0,7799 & 0,8619 \\
MSSQL   & 854.242   & 11,84 & 0,3681 & 0,9901 & 0,5367 \\
NetBIOS & 91.175    & 1,26  & 0,9990 & 0,1480 & 0,2577 \\
Syn     & 4.887.060 & 67,72 & 0,2445 & 0,9999 & 0,3930 \\
UDP     & 1.214.746 & 16,83 & 0,8744 & 0,9996 & 0,9328 \\
UDPLag  & 557       & 0,01  & 0,4040 & 0,0001 & 0,0001 \\
\hline
\end{tabular}
\par\smallskip
Fonte: Elaborado pelo autor (2026).
\end{table}
 
\begin{figure}
\centering
\caption{Matrizes de confusão e métricas por classe do RF multiclasse padrão no teste}
\label{fig:rf-multi-teste-padrao}
% Arquivo: artefatos_modelos/rf_multiclasse/rfm_20261101_210334/teste_multiclasse.png (notebook 06m)
\includegraphics[width=\textwidth]{assets/rf_multi_teste_padrao.png}
\par\smallskip
Fonte: Elaborado pelo autor (2026).
\end{figure}
 
\textbf{Diferença entre os dias de treino e de teste.} A queda do treino para o teste deve ser lida à luz da origem dos dados, como no binário (Seção~\ref{sec:rf-bin-teste}): o treino vem de 12/01/2019 e o teste, de 11/03/2019, com conjuntos de ataques parcialmente diferentes \cite{sharafaldin2019}. A validação é um sorteio de linhas do próprio dia de treino e mede o reconhecimento de fluxos do mesmo dia, não a generalização para outro dia. A diferença é grande: o F1 macro caiu de 0,9007 (validação) para 0,5626 (teste), o MCC de 0,8969 para 0,4324 e a acurácia balanceada de 0,9026 para 0,6908. O padrão dos erros é coerente com essa explicação: os pares que mais se confundem no teste (MSSQL e NetBIOS, UDP e UDPLag, Syn e UDPLag) são os mesmos pares de classes que já conflitavam no treino (Seção~\ref{sec:conflitos-multiclasse}). Essa é a explicação mais plausível, mas não foi isolada por um experimento próprio; o teste, por construção, não tem linhas idênticas às do treino, e por isso os conflitos medidos no treino não podem ser contados nele. Por isso, as métricas do teste não podem ser comparadas diretamente com as da validação.
 
\textbf{Verificação.} As 11 checagens do notebook de teste foram aprovadas.
 
 
\subsubsection{Busca de hiperparâmetros do Random Forest multiclasse}
\label{sec:rf-multi-busca}
 
\textbf{Sementes e protocolo de busca.} O papel da semente e a razão de treinar cada combinação com três sementes (42, 43 e 44) são os do binário (Seção~\ref{sec:rf-bin-busca}). A busca avaliou 27 combinações: 100, 300 ou 500 árvores; profundidade máxima sem limite, 15 ou 8; e \texttt{min\_samples\_leaf} de 1, 5 ou 10. Cada combinação foi examinada com validação cruzada estratificada de cinco dobras no ajuste (as mesmas dobras do item anterior, com impressão digital \texttt{7e97c47f}) e com as três sementes, e também treinada no ajuste completo e avaliada na validação separada, que não determinou a escolha. A regra de seleção foi definida antes da análise dos resultados e usou somente a validação cruzada: ROC-AUC OvR e MCC, regra de um erro-padrão nas duas métricas e, entre as elegíveis, vence a combinação de menor custo, medido pelo número de nós. O teste não participou da decisão. Como no binário, ajuste, validação e validação cruzada vêm do mesmo dia de captura e não medem o comportamento em outro dia.
 
Para reduzir o custo, cada floresta foi treinada com 500 árvores, e as primeiras 100 e 300 árvores avaliaram as configurações menores. O notebook verificou que as primeiras $k$ árvores de uma floresta maior são idênticas às de uma floresta treinada com $k$ árvores (diferença máxima de $3{,}0\times10^{-8}$ nas probabilidades).
 
\textbf{Custo da busca.} Foram treinadas 162 florestas de 500 árvores (9 pares profundidade $\times$ folha, 3 sementes, 5 dobras mais o ajuste completo). A busca levou 5~h 50~min de relógio (20.993~s), e a soma dos tempos de treino foi de cerca de 5~h 35~min, ou seja, cerca de 2 minutos por floresta em média (cálculo a partir da soma). O cálculo das métricas levou mais 3~min 12~s. O custo é cerca de cinco vezes o da busca binária (1~h 08~min), em parte porque o ajuste tem cerca do triplo de registros (239.671 contra 81.228) e as árvores são muito mais profundas. Cada unidade concluída era gravada em disco, o que permite retomar a busca se for interrompida.
 
\textbf{Influência de cada hiperparâmetro.} A Figura~\ref{fig:rf-multi-efeitos} mostra o efeito principal de cada hiperparâmetro. Diferentemente do binário, aqui todos os efeitos de profundidade e de folha ficaram muito acima do ruído (erro-padrão entre dobras): a amplitude do efeito da profundidade foi 30,7 vezes o ruído na ROC-AUC e 26,1 no MCC, e a de \texttt{min\_samples\_leaf}, 14,2 e 6,7 vezes. A quantidade de árvores teve efeito pequeno: 3,0 vezes o ruído na ROC-AUC (acima do ruído, mas de pequena amplitude absoluta) e 0,6 no MCC (dentro do ruído). A profundidade explica 50,7\% da variação da ROC-AUC entre as configurações e 78,7\% da do MCC; \texttt{min\_samples\_leaf}, 13,5\% e 4,7\%. A interação entre os dois responde por 32,8\% da ROC-AUC e 16,5\% do MCC (Figura~\ref{fig:rf-multi-interacao}), pois os dois limitam a complexidade da árvore e atuam em conjunto.
 
\begin{figure}
\centering
\caption{Efeito principal de cada hiperparâmetro na validação cruzada (RF multiclasse)}
\label{fig:rf-multi-efeitos}
% Arquivo: artefatos_modelos/rf_multiclasse_busca/rfm_busca_20261102_001711/influencia_efeitos_principais.png
\includegraphics[width=\textwidth]{assets/rf_multi_influencia_efeitos_principais.png}
\par\smallskip
Fonte: Elaborado pelo autor (2026).
\end{figure}
 
\begin{figure}
\centering
\caption{Interação entre profundidade máxima e \texttt{min\_samples\_leaf} (RF multiclasse)}
\label{fig:rf-multi-interacao}
% Arquivo: .../rfm_busca_20261102_001711/influencia_interacao.png
\includegraphics[width=0.8\textwidth]{assets/rf_multi_influencia_interacao.png}
\par\smallskip
Fonte: Elaborado pelo autor (2026).
\end{figure}
 
O número de árvores determinou principalmente o custo: de 100 para 500 árvores, os nós totais foram, na média dos efeitos principais, de cerca de 672 mil para 3,36 milhões, o modelo de 13,4 para 66,4~MiB e o tempo de treino de 25,1 para 125,7~s. O número de árvores explica 85,8\% da variação do tempo de treino. A Figura~\ref{fig:rf-multi-convergencia} mostra o desempenho em função do número de árvores.
 
\begin{figure}
\centering
\caption{Desempenho em função do número de árvores (RF multiclasse)}
\label{fig:rf-multi-convergencia}
% Arquivo: .../rfm_busca_20261102_001711/convergencia_arvores.png
\includegraphics[width=0.8\textwidth]{assets/rf_multi_convergencia_arvores.png}
\par\smallskip
Fonte: Elaborado pelo autor (2026).
\end{figure}
 
\textbf{Configuração selecionada.} A melhor combinação em ROC-AUC e MCC foi a de 500 árvores, profundidade sem limite e folha 5. Três das 27 combinações atenderam à regra nas duas métricas (100, 300 e 500 árvores, todas com profundidade sem limite e folha 5), e a mais barata, com 100 árvores, foi escolhida: o número de árvores quase não alterou o desempenho, e o ganho veio de \texttt{min\_samples\_leaf}. O padrão (100 árvores, sem limite, folha 1) não era elegível: seu desempenho na validação cruzada (ROC-AUC de 0,99047 e MCC de 0,89874) ficou abaixo do da escolhida (0,99489 e 0,91411) por mais que um erro-padrão. Ou seja, no multiclasse a escolha melhorou a validação cruzada e, ao mesmo tempo, reduziu o custo: o modelo do ajuste passou de 3.483.364 para 1.045.332 nós e de 54,7 para 23,1~MiB. Como conferência, que não decide nada, na validação separada a escolhida ficou em 4º lugar em ROC-AUC e em 1º em MCC. A Tabela~\ref{tab:rf-multi-config-busca} compara as configurações, e a Tabela~\ref{tab:rf-multi-config-comparacao}, os parâmetros.
 
\begin{table}
\centering
\caption{Configurações do RF multiclasse comparadas na busca}
\label{tab:rf-multi-config-busca}
\footnotesize
\resizebox{\textwidth}{!}{%
\begin{tabular}{lrrrrrrrr}
\hline
Configuração & Árvores & Prof. máxima & Folha mínima & ROC-AUC (val. cruzada) & MCC (val. cruzada) & Nós (ajuste) & Tamanho do modelo & Tempo de treino no ajuste \\
\hline
Padrão    & 100 & Sem limite & 1 & 0,99047 & 0,89874 & 3.483.364 & 54,7~MiB & 41,3~s \\
Escolhida & 100 & Sem limite & 5 & 0,99489 & 0,91411 & 1.045.332 & 23,1~MiB & 39,3~s \\
\hline
\end{tabular}}
\par\smallskip
Fonte: Elaborado pelo autor (2026).
\end{table}
 
\begin{table}
\centering
\caption{Configuração padrão e configuração selecionada do Random Forest multiclasse}
\label{tab:rf-multi-config-comparacao}
\begin{tabular}{lcc}
\hline
Parâmetro & Padrão & Selecionado \\
\hline
Número de árvores & 100 & 100 \\
Profundidade máxima & Sem limite & Sem limite \\
Mínimo de registros por divisão & 2 & 2 \\
Mínimo de registros por folha & 1 & 5 \\
Atributos por divisão & $\sqrt{p}$ & $\sqrt{p}$ \\
Peso das classes & Balanceado & Balanceado \\
\hline
\end{tabular}
\par\smallskip
Fonte: Elaborado pelo autor (2026).
\end{table}
 
 
\subsubsection{Avaliação da configuração ajustada do RF multiclasse no teste}
\label{sec:rf-multi-teste-ajustado}
 
A configuração escolhida foi treinada novamente no \textit{refit} (mesmos 299.589 registros do modelo padrão) e avaliada no mesmo teste sem sobreposição, com as mesmas métricas. Na validação, o modelo do ajuste obteve ROC-AUC de 0,9949, F1 macro de 0,9174 e MCC de 0,9136 sem pesos (com a proporção do treino: ROC-AUC de 0,9936, F1 macro de 0,8549 e MCC de 0,8814), acima do padrão. No teste, obteve ROC-AUC OvR de 0,9248, acurácia balanceada de 0,6771, acurácia de 0,3385, F1 macro de 0,4148, F1 ponderado de 0,4258 e MCC de 0,3688. O padrão foi avaliado nas mesmas linhas, e a comparação é descritiva, porque a seleção ocorreu na validação cruzada.
 
O ajuste elevou a ROC-AUC, mas reduziu o desempenho de classificação. A mudança favoreceu o recall de LDAP (de 0,9632 para 0,9926) e de UDP (de 0,8744 para 0,9081), mas reduziu o de Syn (de 0,2445 para 0,1617) e o de UDPLag (de 0,4040 para 0,3573). O principal efeito é a perda de separação entre benigno e ataque: 40,2\% dos registros Syn (1.965.048) passaram a ser previstos como BENIGN (contra 0,07\% no padrão), a precisão do BENIGN caiu de 0,9312 para 0,0273, o recall agrupado dos ataques caiu de 0,99943 para 0,72547 e o MCC ataque $\times$ BENIGN, de 0,95613 para 0,13825. O total de ataques previstos como BENIGN passou de 4.078 para 1.965.726. A Tabela~\ref{tab:rf-multi-comparacao-classe} compara os dois modelos por classe. Como no binário, o ajuste foi feito com dados do mesmo dia do treino e não podia corrigir a diferença entre os dias; aqui, ele piorou o comportamento justamente na classe que domina o teste. A seção seguinte mostra que o padrão tem uma instabilidade própria, que relativiza essa comparação.
 
\begin{table}
\centering
\caption{Recall, precisão e F1 por classe no teste: RF multiclasse padrão e ajustado}
\label{tab:rf-multi-comparacao-classe}
\footnotesize
\resizebox{\textwidth}{!}{%
\begin{tabular}{lrrrrrr}
\hline
 & \multicolumn{2}{c}{Recall} & \multicolumn{2}{c}{Precisão} & \multicolumn{2}{c}{F1} \\
Classe & Padrão & Ajustado & Padrão & Ajustado & Padrão & Ajustado \\
\hline
BENIGN  & 0,9824 & 0,9806 & 0,9312 & 0,0273 & 0,9561 & 0,0531 \\
LDAP    & 0,9632 & 0,9926 & 0,7799 & 0,7597 & 0,8619 & 0,8607 \\
MSSQL   & 0,3681 & 0,3411 & 0,9901 & 0,9920 & 0,5367 & 0,5076 \\
NetBIOS & 0,9990 & 0,9987 & 0,1480 & 0,1440 & 0,2577 & 0,2517 \\
Syn     & 0,2445 & 0,1617 & 0,9999 & 0,9999 & 0,3930 & 0,2784 \\
UDP     & 0,8744 & 0,9081 & 0,9996 & 0,9999 & 0,9328 & 0,9518 \\
UDPLag  & 0,4040 & 0,3573 & 0,0001 & 0,0001 & 0,0001 & 0,0002 \\
\hline
\end{tabular}}
\par\smallskip
Fonte: Elaborado pelo autor (2026).
\end{table}
 
\begin{figure}
\centering
\caption{Matrizes de confusão e métricas por classe do RF multiclasse ajustado no teste}
\label{fig:rf-multi-teste-ajustado}
% Arquivo: artefatos_modelos/rf_multiclasse_busca/rfm_busca_20261102_001711/teste_multiclasse.png (notebook 08m)
\includegraphics[width=\textwidth]{assets/rf_multi_teste_ajustado.png}
\par\smallskip
Fonte: Elaborado pelo autor (2026).
\end{figure}
 
\textbf{Custo computacional.} A Tabela~\ref{tab:rf-multi-custo-treino-tamanho} reúne o custo de treino e o tamanho dos modelos padrão e ajustado. O \textit{refit} da configuração ajustada levou mediana de 48,5~s (de 47,1 a 54,4~s) e 345,8~s de CPU, com acréscimo de RAM de até 414~MiB. O treino no ajuste levou 39,3~s e 279,2~s de CPU. O modelo final tem 1.267.902 nós, profundidade média de 47,5 (máxima de 76) e 27,8~MiB. Em relação ao padrão (\textit{refit} de 70,0~s; 4.288.578 nós; 67,1~MiB), o custo de treino do \textit{refit} é cerca de 30\% menor e o número de nós, cerca de 70\% menor (cálculo a partir da tabela), o oposto do binário, em que o ajuste encareceu o modelo. No teste, a execução completa levou 294,4~s, com 68,9~s de previsão (104.747 fluxos/s). Para um fluxo isolado, a latência mediana fim a fim do modelo final ajustado foi de 18,8~ms com um núcleo e 104,6~ms com todos; em lote, a vazão foi de 19.260 e 66.172 fluxos/s, respectivamente. A inferência com um núcleo ficou mais lenta que a do padrão (9,6~ms e 33.993 fluxos/s), embora o modelo seja menor; não foi investigada a causa.
 
\begin{table}
\centering
\caption{Custo de treino e tamanho dos modelos RF multiclasse (sem os picos de memória do processo)}
\label{tab:rf-multi-custo-treino-tamanho}
\footnotesize
\resizebox{\textwidth}{!}{%
\begin{tabular}{lcccc}
\hline
Item & Padrão: ajuste (239.671) & Padrão: \textit{refit} (299.589) & Ajustado: ajuste (239.671) & Ajustado: \textit{refit} (299.589) \\
\hline
Tempo de relógio (s) & 41,3 (1 execução) & 70,0 (mediana de 3; 55,1 a 74,2) & 39,3 (1 execução) & 48,5 (mediana; 47,1 a 54,4) \\
Tempo de CPU (s) & 258,7 & 372,3 & 279,2 & 345,8 \\
Acréscimo de RAM no treino (MiB) & 542 & até 562 & 333 & até 414 \\
Árvores & 100 & 100 & 100 & 100 \\
Nós (total) & 3.483.364 & 4.288.578 & 1.045.332 & 1.267.902 \\
Profundidade média (máxima) & 61,7 (74) & 65,1 (96) & 46,4 (78) & 47,5 (76) \\
Arquivo compactado (MiB) & 54,7 & 67,1 & 23,1 & 27,8 \\
\hline
\end{tabular}}
\par\smallskip
Fonte: Elaborado pelo autor (2026).
\end{table}
 
\textbf{Verificação.} As 14 checagens do notebook de teste do modelo ajustado foram aprovadas.
 
 
\subsubsection{Estudo da variação entre sementes do RF multiclasse}
\label{sec:rf-multi-sementes}
 
Por fim, foi analisada a variação causada pela aleatoriedade do RF no teste, como no binário (Seção~\ref{sec:rf-bin-sementes}). Os modelos padrão e ajustado foram treinados com cinco sementes (42 a 46), com os mesmos dados de treino e teste, de modo que só a semente mudou. A semente 42 corresponde aos modelos já avaliados, e os outros oito modelos foram treinados agora, em 7~min 30~s (de 53,5 a 55,7~s por modelo padrão e de 48,1 a 50,0~s por ajustado). O ajustado retreinado com a semente 42 saiu idêntico ao modelo da busca (diferença máxima de $3{,}3\times10^{-16}$ nas probabilidades e mesmos 1.267.902 nós), e a semente 42 de cada configuração reproduziu exatamente a matriz de confusão dos testes anteriores. A leitura do teste para o disco levou 2~min 37~s, e a previsão e as métricas dos dez modelos, 12~min 46~s; o notebook inteiro levou cerca de 24~min.
 
Na média das cinco sementes, o RF padrão obteve F1 macro de 0,4834 e MCC de 0,4022, enquanto o ajustado obteve 0,4116 e 0,3651. O resultado da semente 42, usado nas seções anteriores (F1 macro de 0,5626 e MCC de 0,4324), é o melhor das cinco sementes do padrão: o RF padrão tem, portanto, desempenho no teste bem pior, na média, do que o modelo avaliado antes. A ROC-AUC média aumentou de 0,9136 para 0,9247 com o ajuste, e o ajustado foi melhor em ROC-AUC e em recall de UDP (aumento médio de 0,0322) nas cinco sementes, e pior em F1 macro, acurácia balanceada, MCC, recall de Syn e recall de UDPLag nas cinco. As faixas de F1 macro das duas configurações não se sobrepõem; as de MCC, sim.
 
O principal efeito é a estabilidade. O padrão oscila muito: o recall agrupado dos ataques variou de 0,7210 a 0,9994, e a fração de Syn previsto como BENIGN, de 0,07\% a 40,9\% (média de 17,1\%), conforme a semente; o MCC variou de 0,3666 a 0,4324 e o F1 macro, de 0,4183 a 0,5626 (desvio de 0,0677). No ajustado, essa variação é pequena (F1 macro de 0,4085 a 0,4148), mas o resultado é consistentemente ruim: o Syn previsto como BENIGN ficou entre 40,2\% e 47,6\% (média de 46,0\%), e o recall dos ataques, entre 0,6753 e 0,7255. A semente 42 do padrão, portanto, é uma das três (42, 43 e 46) em que o Syn foi, por sorteio, mandado quase todo para UDPLag em vez de BENIGN. Isso reforça a leitura do binário (Seção~\ref{sec:rf-bin-sementes}): os fluxos Syn do dia de teste ficam perto de uma fronteira de decisão, e pequenas diferenças entre florestas decidem o lado para onde caem. Na validação cruzada da busca, o desvio entre sementes foi muito menor (desvio do MCC de 0,00022 no padrão e 0,00009 no ajustado), isto é, a instabilidade só aparece quando o modelo é aplicado a outro dia. A Tabela~\ref{tab:rf-multi-sementes} resume a variação.
 
\begin{table}
\centering
\caption{Variação das métricas do RF multiclasse entre cinco sementes no teste}
\label{tab:rf-multi-sementes}
\footnotesize
\resizebox{\textwidth}{!}{%
\begin{tabular}{lrr}
\hline
Métrica & RF padrão: média (desvio; mín.--máx.) & RF ajustado: média (desvio; mín.--máx.) \\
\hline
ROC-AUC OvR & 0,9136 (0,0045; 0,9072--0,9196) & 0,9247 (0,0014; 0,9224--0,9260) \\
Acurácia balanceada & 0,6843 (0,0057; 0,6775--0,6908) & 0,6731 (0,0027; 0,6697--0,6771) \\
F1 macro & 0,4834 (0,0677; 0,4183--0,5626) & 0,4116 (0,0027; 0,4085--0,4148) \\
F1 macro sem UDPLag & 0,5639 (0,0790; 0,4880--0,6564) & 0,4802 (0,0031; 0,4766--0,4839) \\
MCC & 0,4022 (0,0313; 0,3666--0,4324) & 0,3651 (0,0031; 0,3616--0,3688) \\
Recall do Syn & 0,2029 (0,0362; 0,1668--0,2445) & 0,1505 (0,0108; 0,1366--0,1617) \\
Recall dos ataques (6 classes) & 0,8833 (0,1479; 0,7210--0,9994) & 0,6858 (0,0222; 0,6753--0,7255) \\
Syn previsto como BENIGN (\%) & 17,09 (21,67; 0,07--40,86) & 46,02 (3,25; 40,21--47,57) \\
\hline
\end{tabular}}
\par\smallskip
Fonte: Elaborado pelo autor (2026).
\end{table}
 
\begin{figure}
\centering
\caption{Resultados do RF multiclasse no teste por semente, modelos padrão e ajustado}
\label{fig:rf-multi-sementes}
% Arquivo: .../rfm_busca_20261102_001711/estudo_sementes/grafico_sementes.png
\includegraphics[width=\textwidth]{assets/rf_multi_grafico_sementes.png}
\par\smallskip
Fonte: Elaborado pelo autor (2026).
\end{figure}
 
\textbf{Limitações metodológicas.} Valem as do binário (Seção~\ref{sec:rf-bin-sementes}). Com cinco sementes, os desvios e contagens são descritivos e não constituem teste estatístico. O estudo varia apenas a semente do modelo, não as linhas do ajuste ou do \textit{refit}. O teste vem de um dia de captura diferente do treino, e esse fator influencia fortemente os resultados. Além disso, o teste não foi usado uma única vez: foi avaliado com o RF padrão, depois com o ajustado, reutilizado neste estudo com dez modelos e, adiante, com os SVMs. Nenhum hiperparâmetro ou semente foi escolhido com base nesses resultados, mas a decisão de realizar este estudo veio depois de observar o desempenho dos modelos no teste. O teste deixa de ser, portanto, uma avaliação totalmente independente a partir da Seção~\ref{sec:rf-multi-teste}, e os resultados devem ser lidos como descrição, não como estimativa final de desempenho. Os intervalos de confiança de 95\% calculados nos notebooks são otimistas, pois assumem registros independentes, e o teste repete vetores: 74,1\% dos registros repetem o vetor de atributos de outro registro (medido na avaliação do SVM).
 
 
\subsection{Treino do SVM multiclasse}
\label{sec:svm-multi-treino}
 
\textbf{Objetivo e configuração.} Esta etapa treina o SVM com núcleo RBF nas mesmas sete classes e registra seu desempenho na validação e seu custo. Os conjuntos, os 44 atributos de entrada, o tratamento dos valores ausentes, as métricas e a razão de usar o mesmo modelo de avaliação já foram descritos nas Seções~\ref{sec:visoes-modelos-multiclasse}, \ref{sec:balanceamento-multiclasse}, \ref{sec:rf-multi-treino} e \ref{sec:svm-bin-treino}, e não são repetidos. As diferenças em relação ao SVM binário são três. Primeiro, o scikit-learn resolve o problema de sete classes treinando um SVM para cada par de classes, 21 no total, cada um vendo apenas os registros das suas duas classes; por isso o custo do binário não serve de régua para o multiclasse. Segundo, o escore do SVM não é uma probabilidade: é a decisão um-contra-resto (\texttt{decision\_function}, com sete colunas) derivada dos duelos entre pares, e a classe prevista é a de maior escore; o escore de cada classe é aproximadamente o número de duelos vencidos (de 0 a 6), acrescido de um termo de confiança. Terceiro, assim como no binário, o treino não envolve sorteio, e não há estudo de sementes. O mesmo modelo recebe \texttt{class\_weight="balanced"}, como o RF multiclasse. A normalização é a do binário: \texttt{StandardScaler} nas 42 colunas numéricas, ajustado dentro de cada treino, e os indicadores de ausência permanecem em 0 ou 1; o notebook só conferiu que ela funciona, sem repetir a comparação entre normalizadores, já feita no binário.
 
A configuração inicial usou $C = 1$ e $\gamma =$ \texttt{"scale"}, valor calculado a partir do ajuste inteiro e que, após a padronização, foi de aproximadamente 0,02379 (0,02380 no binário). A Tabela~\ref{tab:svm-multi-config-inicial} resume os parâmetros.
 
\begin{table}
\centering
\caption{Configuração inicial do SVM multiclasse}
\label{tab:svm-multi-config-inicial}
\resizebox{\ifdim\width>\textwidth\textwidth\else\width\fi}{!}{%
\begin{tabular}{ll}
\hline
Parâmetro & Valor \\
\hline
\textit{Kernel} & RBF \\
$C$ & 1 \\
$\gamma$ & \texttt{"scale"} (referência: 0,02379) \\
\texttt{class\_weight} & \texttt{"balanced"} (0,856; 1,422; 2,196) \\
Normalização & \texttt{StandardScaler} nas 42 colunas numéricas \\
Cache do treino & 200 MB por processo \\
Classificadores internos & 21 (um por par de classes) \\
Escore e decisão & \texttt{decision\_function} um-contra-resto; classe de maior escore \\
\hline
\end{tabular}}
\par\smallskip
Fonte: Elaborado pelo autor (2026).
\end{table}
 
\textbf{Calibração do custo (notebook 10m).} Antes da busca, o notebook mediu o custo do SVM nos dados e no computador, sem escolher hiperparâmetros. Cada treino rodou em um processo separado, com limite de tempo. Com $C = 1$ e $\gamma$ de referência, o tempo de treino cresceu de 0,96~s com 7.000 registros (1.000 por classe) para 443,39~s com 111.592 registros (16.000 por classe, exceto NetBIOS, que tem 15.592) (Tabela~\ref{tab:svm-multi-custo-tamanho}); os vetores de suporte passaram de 2.993 para 34.923. Ajustando uma lei de potência, o tempo cresce como $n^{2{,}27}$ e os vetores de suporte como $n^{0{,}88}$, e o custo de prever foi de cerca de 185~ns por par (registro $\times$ vetor de suporte). Como no binário, o ajuste é só uma média: entre 28 mil e 56 mil registros o tempo aumentou 6,8 vezes (de 15,6 para 105,5~s), e a medida serve como ordem de grandeza. Dois pontos distinguem o multiclasse do binário: a fração de registros que viram vetores de suporte é muito maior (31\% a 43\%, contra 2,6\% a 5,5\%), e o tempo do treino completo no ajuste (239.671 registros) foi estimado por extrapolação, pois o maior tamanho medido foi 111.592.
 
\begin{table}
\centering
\caption{Custo do treino do SVM multiclasse em função do tamanho da amostra ($C = 1$, $\gamma$ de referência)}
\label{tab:svm-multi-custo-tamanho}
\resizebox{\ifdim\width>\textwidth\textwidth\else\width\fi}{!}{%
\begin{tabular}{rrrrr}
\hline
Registros & Tempo de treino (s) & CPU (s) & Vetores de suporte & Fração de vetores \\
\hline
7.000   & 0,96   & 0,97   & 2.993  & 42,8\% \\
14.000  & 3,24   & 3,28   & 5.500  & 39,3\% \\
28.000  & 15,61  & 15,58  & 9.913  & 35,4\% \\
56.000  & 105,52 & 104,66 & 18.218 & 32,5\% \\
111.592 & 443,39 & 442,28 & 34.923 & 31,3\% \\
\hline
\end{tabular}}
\par\smallskip
Fonte: Elaborado pelo autor (2026).
\end{table}
 
\begin{figure}
\centering
\caption{Tempo de treino e vetores de suporte do SVM multiclasse conforme o tamanho da amostra}
\label{fig:svm-multi-custo-tamanho}
% Arquivo: artefatos_modelos/svm_multiclasse_calibracao/svm_cal_20261106_233729/custo_por_tamanho.png (notebook 10m)
\includegraphics[width=0.85\textwidth]{assets/svm_multi_custo_por_tamanho.png}
\par\smallskip
Fonte: Elaborado pelo autor (2026).
\end{figure}
 
Três outras medidas da calibração orientaram o protocolo.
 
\begin{itemize}
\item \textbf{Sensibilidade a $C$ e $\gamma$.} Com 56.000 registros (8.000 por classe), as nove configurações da grade levaram, em média, 1,07 vez o tempo da configuração base (no máximo, 1,41 vez), e o canto $C = 100$, $\gamma = 4\times$ custou 1,31 vez o tempo e 0,91 vez os vetores de suporte da base. Isto é, a grade custa pouco a mais que a configuração inicial; o custo vem do tamanho do ajuste, não dos hiperparâmetros.
\item \textbf{Estimativas.} Com a lei de potência, o notebook estimou a busca com validação cruzada de cinco dobras em 1,9~h com 56.000 registros, 7,1~h com 111.592 e 34,6~h com o ajuste inteiro (extrapolado), o \textit{refit} em 1,1 a 1,5~h e a previsão do teste, registro a registro, em 15 a 16,5~h. O notebook recomendou, para a busca, uma amostra de 8.000 registros por classe, mas a busca real usou o ajuste inteiro, como o RF, para que os dois algoritmos fossem comparados com a mesma quantidade de dados.
\item \textbf{Cache e paralelismo.} O \textit{cache} do kernel e o ganho de paralelismo não foram medidos de novo: valem os do binário (Seção~\ref{sec:svm-bin-treino}), isto é, cache de 200 MB por processo, treino com um processo e previsão com dois processos (ganho de 1,84 vez). Na prática, a memória livre do computador limitou a busca e a previsão a um processo (adiante).
\end{itemize}
 
As medições de calibração levaram, no total, cerca de 1~h 03~min (custo $\times$ tamanho, 23~min 24~s; sensibilidade, 40~min 06~s).
 
\textbf{Desempenho inicial na validação.} A configuração inicial ($C = 1$, $\gamma = 1\times$ a referência) foi ajustada nos 239.671 registros do ajuste e avaliada nos 59.918 da validação. Obteve ROC-AUC OvR de 0,97572, acurácia balanceada de 0,86296, F1 macro de 0,86621 e MCC de 0,87212 sem pesos; com a proporção do treino, o F1 macro foi 0,83907 e o MCC, 0,90571. O recall por classe foi baixo em NetBIOS (0,5380) e UDPLag (0,7033), como no RF, e mais alto em BENIGN (0,9959), Syn (0,9729) e UDP (0,9816). Na validação, o RF padrão (F1 macro de 0,9007 e MCC de 0,8969) foi melhor que o SVM padrão. No teste, essa ordem se inverte (Seção~\ref{sec:comparacao-rf-svm-multi}).
 
\textbf{Custo do treino inicial.} O \textit{refit} da configuração inicial, com 299.589 registros, levou 4.257,1~s (1~h 11~min), 4.247,1~s de CPU e pico de 797~MiB de memória. O modelo final inicial tem 84.556 vetores de suporte e ocupa 4,99~MiB.
 
\textbf{Verificação.} Foram aprovadas as oito checagens do notebook 10m: amostras com o mínimo entre o tamanho pedido e as linhas de cada classe, normalização sem vazamento, vetores de suporte presentes em todas as classes e sempre menores que o número de linhas, ROC-AUC entre 0 e 1, sensibilidade cobrindo a grade inteira, estimativas finitas e arquivos gravados.
 
 
\subsection{Busca de hiperparâmetros do SVM multiclasse}
\label{sec:svm-multi-busca}
 
\textbf{Protocolo.} O notebook 11m avaliou nove combinações: $C \in \{1, 10, 100\}$ e $\gamma$ igual a 0,25, 1 ou 4 vezes o valor de referência (0,00595, 0,02379 e 0,09515), a grade do binário (Seção~\ref{sec:svm-bin-busca}). A configuração inicial está na grade. Foi usada a validação cruzada estratificada de cinco dobras no ajuste inteiro (239.671 registros), com as mesmas dobras do RF (impressão digital \texttt{7e97c47f}, conferida por código antes do início: se fosse diferente, a busca parava). Como o SVM não sorteia, cada configuração tem cinco treinos de dobra e um treino completo, num total de 54 SVMs, cada um com os 21 pares de classes. A regra de seleção é a do RF multiclasse: ROC-AUC OvR e MCC, regra de um erro-padrão nas duas métricas e, entre as elegíveis, a de menor custo, medido pelo número de vetores de suporte (que define o tamanho do modelo e o tempo de prever). Se nenhuma configuração for elegível, vale a melhor posição média nas duas métricas (plano B), com o mesmo desempate pelo menor número de vetores de suporte.
 
\textbf{Custo da busca.} Os 54 ajustes foram feitos em um processo (a memória livre, de 0,2~GiB no início, só comportava um), com limite de 4~h por treino; nenhum atingiu o limite e nenhum foi interrompido por falta de memória. A busca levou 28~h 13~min de relógio (em uma sessão), e a soma dos tempos de treino foi de 21~h 16~min (cerca de 24~min por treino em média); a diferença corresponde, provavelmente, à previsão de cada dobra e da validação, feita em cada processo. O cálculo das métricas levou mais 9~s. Os 54 ajustes não foram paralelizados, e o ganho medido (soma dos tempos de treino $\div$ relógio) foi de 0,75, abaixo de 1 pelo motivo apontado. O tempo médio de treino nas dobras variou de 962~s (16~min; $C = 10$, $\gamma = 1\times$) a 1.532~s (25,5~min; $C = 100$, $\gamma = 4\times$), e o de treino completo, de 1.812~s (30~min; $C = 100$, $\gamma = 1\times$) a 3.650~s (61~min; $C = 1$, $\gamma = 4\times$). A estimativa da calibração para a busca com o ajuste inteiro era de 35~h (faixa de 33 a 40~h); a execução real foi mais rápida. Cada treino era gravado ao terminar.
 
\textbf{Influência dos hiperparâmetros.} A ROC-AUC média na validação cruzada ficou entre 0,96921 e 0,97545 nas nove configurações, e o MCC, entre 0,86490 e 0,88105; o número de vetores de suporte, entre 62.797 e 75.238 (nos treinos completos). As diferenças absolutas são pequenas, mas as amplitudes dos efeitos de $C$ e de $\gamma$ ficaram bem acima do ruído (erro-padrão entre dobras): na ROC-AUC, 16,6 vezes para $C$ e 20,5 para $\gamma$; no MCC, 13,0 e 17,0 vezes. Da ROC-AUC, 33,4\% da variação entre as configurações vem de $C$, 48,4\% de $\gamma$ e 18,2\% da interação; do MCC, 33,6\%, 58,6\% e 7,7\%. A partir do padrão, aumentar $C$ de 1 para 100 reduziu a ROC-AUC em 0,00554 (34,0 erros-padrão) e elevou o MCC em 0,00830 (16,2 erros-padrão), e aumentar $\gamma$ de 1 para 4 vezes a referência elevou o MCC em 0,00669 (13,0 erros-padrão), com ROC-AUC praticamente igual (diferença de 0,00023). Reduzir $\gamma$ para 0,25 vez a referência piorou as duas métricas. A ROC-AUC e o MCC, portanto, não apontam para a mesma direção em $C$, o que explica a ausência de configuração elegível (a seguir). Em relação ao custo, os maiores valores de $\gamma$ foram os mais caros no treino (tempo médio de treino completo de 3.281~s com $\gamma = 4\times$, contra 2.460~s com $\gamma = 1\times$), mas produziram os menores números de vetores de suporte.
 
\begin{figure}
\centering
\caption{Efeito dos hiperparâmetros do SVM multiclasse na validação cruzada}
\label{fig:svm-multi-efeitos}
% Arquivo: artefatos_modelos/svm_multiclasse_busca/svm_busca_20261107_013148/influencia_efeitos_principais.png
\includegraphics[width=\textwidth]{assets/svm_multi_influencia_efeitos_principais.png}
\par\smallskip
Fonte: Elaborado pelo autor (2026).
\end{figure}
 
\begin{figure}
\centering
\caption{Interação entre $C$ e $\gamma$ na busca do SVM multiclasse}
\label{fig:svm-multi-interacao}
% Arquivo: .../svm_busca_20261107_013148/influencia_interacao.png
\includegraphics[width=0.8\textwidth]{assets/svm_multi_influencia_interacao.png}
\par\smallskip
Fonte: Elaborado pelo autor (2026).
\end{figure}
 
\textbf{Configuração selecionada.} A melhor configuração em ROC-AUC foi $C = 1$ e $\gamma = 4\times$ (0,97545), e em MCC foi $C = 100$ e $\gamma = 4\times$ (0,88105), e nenhuma configuração ficou dentro de um erro-padrão da melhor nas duas métricas ao mesmo tempo. Foi aplicado, portanto, o plano B. Três configurações ($C = 1$, 10 e 100, com $\gamma = 4\times$) ficaram com a mesma posição média (2,5, cálculo a partir das nove configurações), e o desempate pelo menor número de vetores de suporte escolheu $C = 100$ e $\gamma = 4\times$ a referência (aproximadamente 0,09515), com 62.797 vetores de suporte contra 68.351 do padrão. A escolhida teve o maior MCC na validação cruzada (0,88105, contra 0,87135 do padrão), mas ROC-AUC menor (0,97334, contra 0,97523). Como conferência, que não decide nada, na validação separada a escolhida ficou em 4º lugar em ROC-AUC e em 1º em MCC. Entre as classes, o maior ganho de recall na validação, em relação ao padrão, foi do MSSQL (+0,0195), e a única queda foi a do UDPLag ($-0{,}0019$). A Tabela~\ref{tab:svm-multi-config-busca} compara as configurações, e a Tabela~\ref{tab:svm-multi-config-comparacao}, os parâmetros.
 
\begin{table}
\centering
\caption{Configurações do SVM multiclasse comparadas na busca}
\label{tab:svm-multi-config-busca}
\footnotesize
\resizebox{\textwidth}{!}{%
\begin{tabular}{lrrrrrrr}
\hline
Configuração & $C$ & $\gamma$ ($\times$ ref.) & ROC-AUC (val. cruzada) & MCC (val. cruzada) & Vetores de suporte & Tamanho do modelo & Tempo de treino no ajuste* \\
\hline
Padrão        & 1   & 1 & 0,97523 & 0,87135 & 68.351 & 4,06~MiB & 3.033~s \\
Escolhida     & 100 & 4 & 0,97334 & 0,88105 & 62.797 & 4,00~MiB & 2.730~s \\
Empatada (posição média) & 10  & 4 & 0,97399 & 0,88016 & 64.149 & 4,10~MiB & 3.461~s \\
Empatada (melhor ROC-AUC) & 1   & 4 & 0,97545 & 0,87805 & 65.940 & 4,24~MiB & 3.650~s \\
\hline
\end{tabular}}
\par\smallskip
*Tempo de relógio do treino completo no ajuste (239.671 registros).\\
Fonte: Elaborado pelo autor (2026).
\end{table}
 
\begin{table}
\centering
\caption{Configuração padrão e configuração selecionada do SVM multiclasse}
\label{tab:svm-multi-config-comparacao}
\begin{tabular}{lcc}
\hline
Parâmetro & Padrão & Selecionado \\
\hline
\textit{Kernel} & RBF & RBF \\
$C$ & 1 & 100 \\
$\gamma$ & Referência (\texttt{scale}) & 4 $\times$ referência \\
Peso das classes & Balanceado & Balanceado \\
\hline
\end{tabular}
\par\smallskip
Fonte: Elaborado pelo autor (2026).
\end{table}
 
\textbf{Modelo selecionado na validação.} O modelo do ajuste da configuração escolhida obteve, na validação, ROC-AUC OvR de 0,97355, acurácia balanceada de 0,87104, F1 macro de 0,87566 e MCC de 0,88245 sem pesos; com a proporção do treino, o F1 macro foi 0,85978 e o MCC, 0,92193 (Tabela~\ref{tab:svm-multi-validacao}). O recall foi alto em BENIGN (0,9983), Syn (0,9912) e UDP (0,9884), e baixo em NetBIOS (0,5441) e UDPLag (0,7014) (Tabela~\ref{tab:svm-multi-validacao-classe}). Os erros mais frequentes foram NetBIOS previsto como MSSQL (1.614 registros, 41,4\% do NetBIOS), UDPLag previsto como UDP (2.179, 21,8\%), UDPLag previsto como Syn (687, 6,9\%) e MSSQL previsto como LDAP (614, 6,1\%). São os pares que mais conflitam no treino (Seção~\ref{sec:conflitos-multiclasse}).
 
\begin{table}
\centering
\caption{Métricas do SVM multiclasse selecionado na validação}
\label{tab:svm-multi-validacao}
\begin{tabular}{lcc}
\hline
Métrica & Sem pesos & Proporção do treino \\
\hline
ROC-AUC OvR (macro) & 0,97355 & 0,96889 \\
Acurácia balanceada & 0,87104 & 0,87104 \\
Acurácia & 0,89869 & 0,95838 \\
Precisão macro & 0,90105 & 0,86153 \\
F1 macro & 0,87566 & 0,85978 \\
F1 ponderado & 0,89385 & 0,95783 \\
MCC & 0,88245 & 0,92193 \\
\hline
\end{tabular}
\par\smallskip
Fonte: Elaborado pelo autor (2026).
\end{table}
 
\begin{table}
\centering
\caption{Desempenho do SVM multiclasse selecionado por classe na validação}
\label{tab:svm-multi-validacao-classe}
\resizebox{\textwidth}{!}{%
\begin{tabular}{lrrrrr}
\hline
Classe & Registros & Recall & Precisão (sem pesos) & Precisão (proporção do treino) & ROC-AUC OvR \\
\hline
BENIGN  & 10.000 & 0,9983 & 0,9897 & 0,9766 & 0,9995 \\
LDAP    & 6.020  & 0,9560 & 0,8878 & 0,6801 & 0,9873 \\
MSSQL   & 10.000 & 0,9179 & 0,8263 & 0,9201 & 0,9748 \\
NetBIOS & 3.898  & 0,5441 & 0,8636 & 0,5381 & 0,9213 \\
Syn     & 10.000 & 0,9912 & 0,9344 & 0,9616 & 0,9978 \\
UDP     & 10.000 & 0,9884 & 0,8192 & 0,9820 & 0,9813 \\
UDPLag  & 10.000 & 0,7014 & 0,9864 & 0,9722 & 0,9530 \\
\hline
\end{tabular}}
\par\smallskip
Fonte: Elaborado pelo autor (2026).
\end{table}
 
\textbf{Estabilidade do resultado.} Como no RF, a validação foi dividida em cinco partições (Tabela~\ref{tab:svm-multi-dispersao}). Na validação cruzada da configuração escolhida (as mesmas cinco dobras, Figura~\ref{fig:svm-multi-cv-ajuste}), cada SVM foi treinado com cerca de 191.737 registros, levou em média 1.531,9~s (de 1.503,7 a 1.554,8~s) e teve em média 50.682 vetores de suporte. As métricas variaram pouco, e os valores das dobras (MCC de 0,88105) ficaram próximos dos da validação.
 
\begin{table}
\centering
\caption{Dispersão das métricas do SVM multiclasse selecionado (média $\pm$ desvio-padrão; mín--máx)}
\label{tab:svm-multi-dispersao}
\resizebox{\textwidth}{!}{%
\begin{tabular}{lcc}
\hline
Métrica & 5 partições da validação & 5 dobras da validação cruzada \\
\hline
ROC-AUC OvR & 0,97355 $\pm$ 0,00079 (0,97236--0,97425) & 0,97334 $\pm$ 0,00032 (0,97288--0,97367) \\
Acurácia balanceada & 0,87104 $\pm$ 0,00334 (0,86597--0,87493) & 0,86908 $\pm$ 0,00160 (0,86745--0,87157) \\
F1 macro & 0,87563 $\pm$ 0,00353 (0,87013--0,87929) & 0,87373 $\pm$ 0,00162 (0,87211--0,87628) \\
MCC (sem pesos) & 0,88246 $\pm$ 0,00299 (0,87759--0,88472) & 0,88105 $\pm$ 0,00106 (0,88007--0,88281) \\
MCC (proporção do treino) & 0,92194 $\pm$ 0,00412 (0,91547--0,92601) & não calculado* \\
\hline
\end{tabular}}
\par\smallskip
*A métrica com a proporção do treino só foi calculada na validação (Tabela~\ref{tab:pesos-multiclasse}); na validação cruzada, as métricas foram calculadas sem pesos.\\
Fonte: Elaborado pelo autor (2026).
\end{table}

\begin{figure}
    \centering
    \caption{Validação cruzada no ajuste do SVM multiclasse (5 dobras)}
    \label{fig:svm-multi-cv-ajuste}
    \includegraphics[width=\textwidth]{assets/svm_multi_cv_ajuste_5dobras.png}
    \par\smallskip
    Fonte: Elaborado pelo autor (2026).
\end{figure}
 
\textbf{\textit{Refit} (modelo final).} O treino final usa a configuração escolhida com os 299.589 registros do \textit{refit} (os mesmos do RF). Como o SVM é determinístico, foi feita uma única repetição, que levou 5.308,4~s (1~h 28~min) e 5.300,5~s de CPU, com pico de 743~MiB. O modelo final tem 77.699 vetores de suporte (3.331 de BENIGN, 4.022 de LDAP, 27.393 de MSSQL, 10.217 de NetBIOS, 4.368 de Syn, 12.225 de UDP e 16.143 de UDPLag) e ocupa 4,88~MiB. O MSSQL concentra 35\% dos vetores de suporte, possivelmente por sua sobreposição com o NetBIOS e o LDAP (Seção~\ref{sec:conflitos-multiclasse}).
 
\textbf{Custo computacional.} O treino no ajuste levou 2.730,1~s (45~min 30~s) de relógio e 2.723,6~s de CPU, cerca de um núcleo ocupado, com acréscimo de 429~MiB de memória (pico de 757~MiB). O modelo do ajuste tem 62.797 vetores de suporte e 4,00~MiB. A Tabela~\ref{tab:svm-multi-custo-treino} reúne o custo dos modelos. Em relação ao RF padrão (\textit{refit} de 70~s), o \textit{refit} do SVM escolhido (5.308~s) é cerca de 76 vezes mais lento (cálculo a partir dos valores), embora o arquivo do modelo seja muito menor (4,9 contra 67,1~MiB). Na inferência, com um processo (a memória livre, de 1,2~GiB, não comportava dois), a latência mediana de um fluxo do modelo final, fim a fim, foi de 14,6~ms (p95 de 18,8~ms e p99 de 25,0~ms), da mesma ordem que a do RF com um núcleo (9,6~ms). Em lote, porém, a vazão foi de apenas 77,9 fluxos/s (lote de 5.000 fluxos), contra 33.993 do RF com um núcleo, cerca de 436 vezes menos (cálculo a partir dos valores); como no binário, é a comparação de cada fluxo com todos os vetores de suporte que torna a previsão cara (Tabela~\ref{tab:svm-multi-inferencia}).
 
\begin{table}
\centering
\caption{Custo de treino e tamanho dos modelos SVM multiclasse}
\label{tab:svm-multi-custo-treino}
\footnotesize
\resizebox{\textwidth}{!}{%
\begin{tabular}{lccc}
\hline
Item & Selecionado: ajuste (239.671) & Selecionado: \textit{refit} (299.589) & Padrão: \textit{refit} (299.589) \\
\hline
Tempo de relógio (s) & 2.730,1 (1 execução) & 5.308,4 (1 execução) & 4.257,1 (1 execução) \\
Tempo de CPU (s) & 2.723,6 & 5.300,5 & 4.247,1 \\
Pico de RAM (MiB) & 757 & 743 & 797 \\
Acréscimo de RAM no treino (MiB) & 429 & 493 & $\approx$ 547* \\
Vetores de suporte & 62.797 & 77.699 & 84.556 \\
Arquivo compactado (MiB) & 4,00 & 4,88 & 4,99 \\
\hline
\end{tabular}}
\par\smallskip
*Estimado, não medido: o notebook gravou só o pico de RAM do \textit{refit} do padrão (797~MiB). O acréscimo foi obtido subtraindo o consumo do processo antes do treino, de cerca de 250~MiB, valor obtido no \textit{refit} do modelo selecionado (pico de 742,5~MiB menos acréscimo de 493,0~MiB).\\
Fonte: Elaborado pelo autor (2026).
\end{table}
 
\begin{table}
\centering
\caption{Inferência do SVM multiclasse (fim a fim, 1 processo)}
\label{tab:svm-multi-inferencia}
\begin{tabular}{lcc}
\hline
Medida & Modelo do ajuste & Modelo final (\textit{refit}) \\
\hline
Latência p50 / p95 / p99 (ms) & 12,7 / 20,2 / 22,9 & 14,6 / 18,8 / 25,0 \\
Vazão em lote (fluxos/s) & 96,5 & 77,9 \\
\hline
\end{tabular}
\par\smallskip
Fonte: Elaborado pelo autor (2026).
\end{table}
 
\textbf{Verificação.} As 20 checagens do notebook 11m foram aprovadas. Entre elas, a busca completa (54 treinos e 9 configurações), as dobras iguais às do RF (impressão digital \texttt{7e97c47f}), a reprodução dos escores pelos modelos gravados, os pesos de classe conferidos com o scikit-learn, o \textit{refit} com as mesmas linhas do RF, a reprodutibilidade da escolha a partir dos resultados gravados e a ausência de treinos interrompidos por falta de memória.
 
 
\subsection{Avaliação do SVM multiclasse no teste}
\label{sec:svm-multi-teste}
 
\textbf{Preparação e composição do teste.} Os dois modelos finais do SVM (o padrão, com $C = 1$ e $\gamma = 1\times$ a referência, e o selecionado, com $C = 100$ e $\gamma = 4\times$) foram avaliados no mesmo teste do RF (7.216.487 registros, mais os 2.702 do Portmap à parte), com as garantias do RF: nenhum modelo foi treinado, nenhum arquivo existente foi sobrescrito (o notebook conferiu os 672 arquivos preexistentes antes e depois), o servidor só foi lido e o teste foi usado uma única vez para cada modelo, com bloqueio de nova execução.
 
\textbf{Previsão dos vetores únicos.} Prever as 7,2 milhões de linhas com cerca de 78 mil vetores de suporte levaria cerca de 25 a 27~h por modelo (estimativa do notebook). O teste tem muitos registros com exatamente os mesmos 44 atributos, e o SVM é uma função determinística deles: o mesmo vetor sempre dá a mesma previsão. Por isso o notebook leu o teste uma vez, identificou os vetores distintos (comparação exata, por uma chave de 128 bits), previu cada um uma única vez e espalhou o resultado para todas as linhas, cada uma contada com o seu rótulo real. Foram encontrados 1.865.985 vetores distintos (25,9\% das linhas; 74,1\% são repetidas), o que levou 2~min 10~s. Para provar a equivalência, 3.000 linhas do teste foram previstas diretamente pelo modelo: as classes foram iguais em 100\% das linhas, com diferença máxima de $2{,}4\times10^{-7}$ nos escores (diferença de precisão entre \textit{float32} e \textit{float64}). Os tempos de previsão abaixo, portanto, descrevem a previsão dos vetores únicos e a expansão para as linhas do teste, não uma execução independente em cada registro.
 
\textbf{Resultados dos modelos padrão e ajustado.} A previsão dos 1.865.985 vetores levou 6~h 17~min (22.647~s) para o SVM selecionado e 6~h 37~min (23.837~s) para o padrão, ambos com um processo (ganho de paralelismo de 0,98 e 0,99) e memória por processo de cerca de 0,6~GiB. Isso equivale a 82,4 e 78,3 vetores por segundo, ou 318,6 e 302,7 registros por segundo considerando as linhas do teste. A Tabela~\ref{tab:svm-multi-teste} apresenta os resultados. O SVM padrão obteve ROC-AUC OvR de 0,8855, acurácia balanceada de 0,6914, acurácia de 0,5996, F1 macro de 0,6119, F1 ponderado de 0,7213 e MCC de 0,5347. O ajustado elevou a ROC-AUC e a acurácia balanceada (de 0,6914 para 0,7335), mas ficou abaixo do padrão em acurácia, F1 macro e MCC: a configuração ajustada mudou o equilíbrio entre as classes sem melhorar todas as métricas.
 
\begin{table}
\centering
\caption{Comparação dos modelos SVM multiclasse no teste}
\label{tab:svm-multi-teste}
\footnotesize
\resizebox{\textwidth}{!}{%
\begin{tabular}{lrrrrrrr}
\hline
Modelo & ROC-AUC OvR & Acurácia balanceada & Acurácia & F1 macro & F1 macro sem UDPLag & F1 ponderado & MCC \\
\hline
SVM padrão    & 0,8855 & 0,6914 & 0,5996 & 0,6119 & 0,7139 & 0,7213 & 0,5347 \\
SVM ajustado  & 0,8997 & 0,7335 & 0,5646 & 0,6022 & 0,7026 & 0,6895 & 0,5172 \\
\hline
\end{tabular}}
\par\smallskip
Fonte: Elaborado pelo autor (2026).
\end{table}
 
\begin{figure}
\centering
\caption{Matrizes de confusão e métricas por classe do SVM multiclasse padrão no teste}
\label{fig:svm-multi-teste-padrao}
% Arquivo: artefatos_modelos/svm_multiclasse_busca/svm_busca_20261107_013148/teste_padrao/teste_multiclasse.png (notebook 12m)
\includegraphics[width=\textwidth]{assets/svm_multi_teste_padrao.png}
\par\smallskip
Fonte: Elaborado pelo autor (2026).
\end{figure}
 
\begin{figure}
\centering
\caption{Matrizes de confusão e métricas por classe do SVM multiclasse ajustado no teste}
\label{fig:svm-multi-teste-ajustado}
% Arquivo: .../svm_busca_20261107_013148/teste_ajustado/teste_multiclasse.png (notebook 12m)
\includegraphics[width=\textwidth]{assets/svm_multi_teste_ajustado.png}
\par\smallskip
Fonte: Elaborado pelo autor (2026).
\end{figure}
 
\textbf{Desempenho por classe.} O padrão dos erros é parecido com o do RF (Tabela~\ref{tab:svm-multi-teste-classe}). O Syn (67,7\% do teste) teve recall de 0,5273 no padrão e 0,4736 no ajustado, e a maior parte do resto foi classificada como UDPLag (47,1\% e 52,4\% do Syn: 2.304.231 e 2.560.291 registros). O MSSQL foi confundido com o NetBIOS (57,2\% e 58,9\% do MSSQL), e por isso o recall do MSSQL ficou em 0,38 e 0,37, e a precisão do NetBIOS, em 0,15. Diferentemente do RF ajustado, o SVM ajustado quase não manda Syn para BENIGN (0,18\%, 8.553 registros, contra 40,2\% do RF ajustado), e sua precisão do BENIGN foi 0,8604. Entre os ataques, apenas 3.968 (SVM padrão) e 8.986 (SVM ajustado) foram previstos como BENIGN, e o MCC ataque $\times$ BENIGN foi 0,96017 e 0,91994, com falso alarme de 1,16\% e 1,50\%. O UDPLag, com 557 registros, é o caso extremo: o SVM padrão não acertou nenhum (recall de 0,0000; a maior parte foi classificada como MSSQL, 61,6\%), e o ajustado acertou 193 (recall de 0,3465; intervalo de 95\% otimista de 0,308 a 0,387). Os dois modelos receberam mais de 2,3 milhões de previsões de UDPLag, das quais 99,98\% (ajustado) e 99,999\% (padrão) vieram de registros Syn; a precisão do UDPLag é praticamente zero e sua ROC-AUC OvR (0,2180 no padrão e 0,3270 no ajustado) fica abaixo de 0,5. Por isso, o F1 macro sem UDPLag (secundário, definido antes de ver os resultados) é apresentado ao lado do F1 macro.
 
\begin{table}
\centering
\caption{Desempenho dos SVMs multiclasse por classe no teste}
\label{tab:svm-multi-teste-classe}
\footnotesize
\resizebox{\textwidth}{!}{%
\begin{tabular}{lrrrrrr}
\hline
 & \multicolumn{2}{c}{Recall} & \multicolumn{2}{c}{Precisão} & \multicolumn{2}{c}{F1} \\
Classe & Padrão & Ajustado & Padrão & Ajustado & Padrão & Ajustado \\
\hline
BENIGN  & 0,9884 & 0,9850 & 0,9333 & 0,8604 & 0,9601 & 0,9185 \\
LDAP    & 0,9847 & 0,9823 & 0,7476 & 0,7830 & 0,8500 & 0,8714 \\
MSSQL   & 0,3835 & 0,3748 & 0,9616 & 0,9842 & 0,5483 & 0,5429 \\
NetBIOS & 0,9964 & 0,9987 & 0,1463 & 0,1451 & 0,2552 & 0,2533 \\
Syn     & 0,5273 & 0,4736 & 0,9999 & 1,0000 & 0,6904 & 0,6427 \\
UDP     & 0,9596 & 0,9736 & 0,9998 & 0,9999 & 0,9793 & 0,9866 \\
UDPLag  & 0,0000 & 0,3465 & 0,0000 & 0,0001 & 0,0000 & 0,0002 \\
\hline
\end{tabular}}
\par\smallskip
Fonte: Elaborado pelo autor (2026).
\end{table}
 
\textbf{Custo computacional.} A previsão dos vetores únicos levou 6~h 17~min (ajustado) e 6~h 37~min (padrão), mais de 12~h somadas, e o notebook inteiro levou cerca de 13~h. Mesmo com o atalho dos vetores únicos, a vazão equivalente (303 a 319 registros por segundo) é cerca de 370 a 390 vezes menor que a do RF padrão (117.989 por segundo; cálculo a partir dos valores).
 
\textbf{Limitações.} Valem as do RF (Seção~\ref{sec:rf-multi-sementes}): o teste vem de outro dia de captura, é descritivo e foi reutilizado. O SVM não tem semente, então seu resultado é uma nota só, sem faixa de variação, ao contrário do RF. O intervalo de confiança de 95\% é otimista. O resultado do UDPLag se apoia em 557 registros.
 
\textbf{Verificação.} As 9 checagens do notebook 12m foram aprovadas: matrizes que fecham, equivalência com a previsão direta, mapa do teste coerente (1.865.985 vetores distintos em 7.216.487 linhas), todos os vetores previstos, todas as linhas Portmap lidas e previstas, nenhum modelo treinado e nenhum arquivo preexistente alterado.
 
 
\subsection{Análise dos conflitos de rótulo entre as sete classes}
\label{sec:conflitos-multiclasse}
 
\textbf{Motivação e definição.} Registros com exatamente os mesmos atributos e rótulos diferentes não podem ser todos classificados corretamente. O notebook de balanceamento marcou, para o conjunto das 13 classes do treino, as linhas em conflito (\texttt{conflito\_rotulo}), mas essa marca exagera o problema do multiclasse: o LDAP, por exemplo, teve 42,9\% das linhas marcadas, quase sempre em conflito com DNS e SNMP, que ficaram fora do modelo (Seção~\ref{sec:balanceamento-multiclasse}). Este notebook recalculou o conflito considerando somente as sete classes (\texttt{conflito7}): uma linha está em conflito quando seus atributos aparecem em registros do treino com dois ou mais rótulos entre as sete. Os grupos de atributos iguais são os já calculados no pré-processamento, e nada é recalculado do zero. O teste não é tocado.
 
\textbf{Resultados.} Das 299.589 linhas do ajuste e da validação, 41.668 (13,9\%) estão marcadas com conflito pelas 13 classes, e 19.864 (6,6\%) estão em conflito entre as sete; isto é, 47,7\% da marca original é um conflito real do problema de sete classes, e as outras 21.804 só tinham como parceiras classes que saíram do modelo. O ajuste tem 15.896 linhas em conflito, e a validação, 3.968 (6,6\% de cada conjunto). A frequência varia muito por classe (Tabela~\ref{tab:conflito7-classe}): o NetBIOS tem 44,8\% do ajuste em conflito, o UDPLag, 10,8\%, e o LDAP e o MSSQL, 5,8\% e 5,6\%, enquanto BENIGN, UDP e Syn têm 1,4\% ou menos. Os percentuais da validação são quase iguais aos do ajuste, como esperado de uma amostra do mesmo treino.
 
\begin{table}
\centering
\caption{Registros em conflito de rótulo por classe: 13 classes e 7 classes}
\label{tab:conflito7-classe}
\footnotesize
\resizebox{\textwidth}{!}{%
\begin{tabular}{lrrrrrrr}
\hline
 & \multicolumn{3}{c}{Ajuste} & \multicolumn{3}{c}{Validação} \\
Classe & Linhas & Marca das 13 classes & Entre as 7 (\%) & Linhas & Marca das 13 classes & Entre as 7 (\%) \\
\hline
BENIGN  & 40.000 & 17     & 3 (0,0\%)      & 10.000 & 4     & 2 (0,0\%) \\
LDAP    & 24.079 & 10.323 & 1.399 (5,8\%)  & 6.020  & 2.589 & 358 (5,9\%) \\
MSSQL   & 40.000 & 3.784  & 2.254 (5,6\%)  & 10.000 & 988   & 591 (5,9\%) \\
NetBIOS & 15.592 & 8.149  & 6.986 (44,8\%) & 3.898  & 1.999 & 1.712 (43,9\%) \\
Syn     & 40.000 & 647    & 561 (1,4\%)    & 10.000 & 178   & 138 (1,4\%) \\
UDP     & 40.000 & 5.559  & 391 (1,0\%)    & 10.000 & 1.373 & 94 (0,9\%) \\
UDPLag  & 40.000 & 4.843  & 4.302 (10,8\%) & 10.000 & 1.215 & 1.073 (10,7\%) \\
\hline
\end{tabular}}
\par\smallskip
Fonte: Elaborado pelo autor (2026).
\end{table}
 
Os conflitos se concentram em poucos pares de classes. As combinações de rótulos mais frequentes entre as sete (ajuste e validação juntos) foram MSSQL e NetBIOS (9.009 linhas), UDP e UDPLag (4.354), Syn e UDPLag (1.923) e LDAP e MSSQL (1.835). Esses pares coincidem com os erros mais frequentes dos modelos, na validação (MSSQL e NetBIOS, UDP e UDPLag; Seções~\ref{sec:rf-multi-treino} e \ref{sec:svm-multi-busca}) e no teste (MSSQL para NetBIOS e Syn e UDP para UDPLag; Seções~\ref{sec:rf-multi-teste} e \ref{sec:svm-multi-teste}). A coincidência é sugestiva, mas não demonstra causa: os erros no teste não podem ser explicados por conflitos com o treino, pois o teste foi deduplicado em relação a ele.
 
\textbf{Versão do ajuste sem conflitos.} O notebook também preparou, como análise de sensibilidade, o ajuste sem as 15.896 linhas em conflito. O ajuste passa de 239.671 para 223.775 linhas e a composição muda, sobretudo no NetBIOS, que fica com 8.606 linhas (55,2\% do original) e é a menor classe, com peso de classe de 3,715 (contra 2,196). A razão entre a maior e a menor classe passa de 2,57 para 1 a 4,65 para 1. A validação não muda (as mesmas 59.918 linhas e pesos), para que um modelo sem conflitos e o normal sejam medidos exatamente nas mesmas linhas. As 15 checagens do notebook foram aprovadas, entre elas, que todo conflito entre as sete classes é também conflito entre as 13, que a conta feita aqui com 13 classes coincide linha a linha com a marca do pré-processamento (zero divergências) e que a seleção sem conflitos é a anterior menos as linhas conflitantes do ajuste. Nenhum modelo foi treinado com essa versão, e portanto seus resultados não devem ser apresentados como desempenho de um modelo sem conflitos. Essa comparação, medida na mesma validação, não foi executada nesta etapa.
 
\textbf{Análise dos erros do RF por conflito.} O modelo do ajuste do RF padrão foi usado para separar os erros da validação pelas linhas com e sem conflito entre as sete classes (o notebook recalculou a acurácia, 0,91265, e confirmou que coincide com a gravada). Dos 59.918 registros, 3.968 eram conflitantes. O modelo acertou 47,3\% deles, contra 94,4\% dos sem conflito. Dos 5.234 erros, 2.092 (40,0\%) ocorreram nos registros conflitantes, que são 6,6\% da validação, e entre esses erros, 2.015 (96,3\%) foram previsões da classe parceira (a do rótulo gêmeo). A taxa de erro foi 52,7\% nos conflitantes e 5,6\% nos demais (Tabela~\ref{tab:erros-conflito}). Isto é, nos registros conflitantes o modelo quase sempre previu a outra classe associada ao mesmo vetor, o que é compatível com a ambiguidade dos rótulos.
 
\begin{table}
\centering
\caption{Desempenho do RF padrão na validação segundo a presença de conflito}
\label{tab:erros-conflito}
\begin{tabular}{lrrrr}
\hline
Grupo & Registros & Acurácia & Erros & Taxa de erro \\
\hline
Sem conflito & 55.950 & 0,9438 & 3.142 & 5,6\% \\
Com conflito & 3.968  & 0,4728 & 2.092 & 52,7\% \\
\hline
\end{tabular}
\par\smallskip
Fonte: Elaborado pelo autor (2026).
\end{table}
 
O recall por classe nas duas situações (Tabela~\ref{tab:recall-conflito}) mostra que o efeito é desigual: nos registros conflitantes o recall foi especialmente baixo em MSSQL (0,0964), Syn (0,2174) e UDP (0,1383), e mais alto em LDAP (0,6704), NetBIOS (0,5613) e UDPLag (0,5359). Nos registros sem conflito, o recall do NetBIOS (0,9209) é bem maior que o recall geral da classe (0,7630): 1.712 dos 3.898 registros do NetBIOS (43,9\%) estão em conflito, e eles concentram cerca de 81\% dos erros da classe (cálculo a partir do recall e das contagens). O BENIGN tem só 2 registros conflitantes, e o resultado dessa linha não é interpretável. Este resultado ajuda a interpretar os limites do modelo, mas não substitui uma avaliação causal: os conflitos são uma associação entre vetores e rótulos e não demonstram, isoladamente, a origem do problema; e a análise foi feita com um único modelo (RF padrão), só na validação.
 
\begin{table}
\centering
\caption{Recall do RF padrão na validação, nas linhas com e sem conflito entre as sete classes}
\label{tab:recall-conflito}
\resizebox{\ifdim\width>\textwidth\textwidth\else\width\fi}{!}{%
\begin{tabular}{lrrrr}
\hline
Classe & Linhas sem conflito & Recall sem conflito & Linhas em conflito & Recall em conflito \\
\hline
BENIGN  & 9.998 & 0,9992 & 2     & 0,0000 \\
LDAP    & 5.662 & 0,9832 & 358   & 0,6704 \\
MSSQL   & 9.409 & 0,9296 & 591   & 0,0964 \\
NetBIOS & 2.186 & 0,9209 & 1.712 & 0,5613 \\
Syn     & 9.862 & 0,9653 & 138   & 0,2174 \\
UDP     & 9.906 & 0,9254 & 94    & 0,1383 \\
UDPLag  & 8.927 & 0,8742 & 1.073 & 0,5359 \\
\hline
\end{tabular}}
\par\smallskip
Fonte: Elaborado pelo autor (2026).
\end{table}
 
 
\subsection{Comparação entre Random Forest e SVM multiclasse}
\label{sec:comparacao-rf-svm-multi}
 
\textbf{Desempenho no teste.} A comparação foi feita nos mesmos 7.216.487 registros, com os mesmos atributos e rótulos, e é descritiva (o teste tem uma partição e muitos registros repetidos). A Tabela~\ref{tab:comparacao-multi-teste} reúne os quatro modelos. O RF ajustado teve a maior ROC-AUC OvR (0,9248), seguido do RF padrão (0,9133); o SVM ajustado teve a maior acurácia balanceada (0,7335); e o SVM padrão teve o melhor F1 macro (0,6119), F1 ponderado (0,7213), acurácia (0,5996) e MCC (0,5347), com acurácia balanceada praticamente igual à do RF padrão (0,6914 e 0,6908). Como o RF padrão oscila muito entre sementes (média de F1 macro de 0,4834 e MCC de 0,4022; Seção~\ref{sec:rf-multi-sementes}), a vantagem do SVM padrão sobre ele é maior do que a Tabela~\ref{tab:comparacao-multi-teste} sugere, pois o resultado da semente 42, ali usado, é o melhor das cinco. O SVM não tem essa variação. A escolha depende da métrica priorizada: o RF ordena melhor os escores segundo a ROC-AUC, mas os SVMs apresentaram melhor desempenho agregado nas classificações finais.
 
\begin{table}
\centering
\caption{Desempenho dos modelos multiclasse no teste}
\label{tab:comparacao-multi-teste}
\footnotesize
\resizebox{\textwidth}{!}{%
\begin{tabular}{lrrrrrrr}
\hline
Modelo & ROC-AUC OvR & Acurácia balanceada & Acurácia & F1 macro & F1 macro sem UDPLag & F1 ponderado & MCC \\
\hline
RF padrão    & 0,9133 & 0,6908 & 0,3917 & 0,5626 & 0,6564 & 0,5108 & 0,4324 \\
RF ajustado  & 0,9248 & 0,6771 & 0,3385 & 0,4148 & 0,4839 & 0,4258 & 0,3688 \\
SVM padrão   & 0,8855 & 0,6914 & 0,5996 & 0,6119 & 0,7139 & 0,7213 & 0,5347 \\
SVM ajustado & 0,8997 & 0,7335 & 0,5646 & 0,6022 & 0,7026 & 0,6895 & 0,5172 \\
\hline
\end{tabular}}
\par\smallskip
Fonte: Elaborado pelo autor (2026).
\end{table}
 
\textbf{Por classe.} O SVM padrão melhora o recall de BENIGN, LDAP, MSSQL, Syn e UDP em relação ao RF padrão, enquanto o RF tem recall ligeiramente superior para NetBIOS e é o único dos modelos padrão a identificar parte dos UDPLag (Tabela~\ref{tab:comparacao-multi-recall}). A maior diferença entre os modelos está no Syn, que domina o teste: o recall passa de 0,2445 (RF padrão) e 0,1617 (RF ajustado) para 0,5273 (SVM padrão) e 0,4736 (SVM ajustado). A classe UDPLag continua sendo um ponto crítico, pois todos os modelos têm precisão praticamente zero nela, e o SVM padrão não acertou nenhum dos 557 registros. O SVM ajustado recupera 0,3465 de recall, comparável ao RF.
 
\begin{table}
\centering
\caption{Recall por classe nos modelos multiclasse avaliados no teste}
\label{tab:comparacao-multi-recall}
\begin{tabular}{lrrrr}
\hline
Classe & RF padrão & RF ajustado & SVM padrão & SVM ajustado \\
\hline
BENIGN  & 0,9824 & 0,9806 & 0,9884 & 0,9850 \\
LDAP    & 0,9632 & 0,9926 & 0,9847 & 0,9823 \\
MSSQL   & 0,3681 & 0,3411 & 0,3835 & 0,3748 \\
NetBIOS & 0,9990 & 0,9987 & 0,9964 & 0,9987 \\
Syn     & 0,2445 & 0,1617 & 0,5273 & 0,4736 \\
UDP     & 0,8744 & 0,9081 & 0,9596 & 0,9736 \\
UDPLag  & 0,4040 & 0,3573 & 0,0000 & 0,3465 \\
\hline
\end{tabular}
\par\smallskip
Fonte: Elaborado pelo autor (2026).
\end{table}
 
\textbf{Ataque contra benigno.} Como o objetivo do sistema de detecção de intrusão é, antes de tudo, separar tráfego benigno de ataque, os seis tipos de ataque foram agrupados em uma só classe (cálculo posterior, descritivo). Nesse recorte, três dos quatro modelos têm desempenho muito próximo (MCC de 0,92 a 0,96), e o RF ajustado destoa (Tabela~\ref{tab:comparacao-multi-ataque}). O RF ajustado classificou 1.965.726 ataques como BENIGN, quase todos Syn. O SVM padrão e o RF padrão (semente 42) são equivalentes, mas o RF padrão só tem esse resultado em parte das sementes: com outras sementes, o recall de ataques cai a 0,72 (Seção~\ref{sec:rf-multi-sementes}). O SVM é determinístico.
 
\begin{table}
\centering
\caption{Ataque contra BENIGN no teste (as seis classes de ataque agrupadas)}
\label{tab:comparacao-multi-ataque}
\footnotesize
\resizebox{\textwidth}{!}{%
\begin{tabular}{lrrrr}
\hline
Medida & RF padrão & RF ajustado & SVM padrão & SVM ajustado \\
\hline
Recall dos ataques & 0,99943 & 0,72547 & 0,99945 & 0,99875 \\
Falso alarme (BENIGN previsto como ataque) & 1,76\% & 1,94\% & 1,16\% & 1,50\% \\
MCC ataque $\times$ BENIGN & 0,95613 & 0,13825 & 0,96017 & 0,91994 \\
Ataques previstos como BENIGN & 4.078 & 1.965.726 & 3.968 & 8.986 \\
Syn previsto como BENIGN & 0,07\% & 40,21\% & 0,07\% & 0,18\% \\
Syn previsto como UDPLag & 75,40\% & 43,56\% & 47,15\% & 52,39\% \\
\hline
\end{tabular}}
\par\smallskip
Fonte: Elaborado pelo autor (2026).
\end{table}
 
\textbf{Linha a linha.} Comparando as previsões dos modelos do mesmo tipo (padrão com padrão, ajustado com ajustado) nas mesmas linhas, há muito mais linhas que só o SVM acerta do que linhas que só o RF acerta, e a diferença vem sobretudo do Syn (Tabela~\ref{tab:comparacao-multi-linha}): no ajustado, o SVM acerta 1.706.130 linhas que o RF erra, e o RF acerta 74.492 que o SVM erra (saldo de 1.631.638 a favor do SVM; Syn, 1.524.139). Os dois erram juntos em 3.067.778 linhas. Nos outros tipos o saldo é bem menor (UDP, 79.583; MSSQL, 28.825) ou negativo (LDAP, $-1.151$).
 
\begin{table}
\centering
\caption{Comparação linha a linha entre SVM e RF do mesmo tipo no teste}
\label{tab:comparacao-multi-linha}
\footnotesize
\resizebox{\textwidth}{!}{%
\begin{tabular}{lrrrrr}
\hline
Comparação & Certas nos dois & Só o SVM acerta & Só o RF acerta & Erram nos dois & Saldo (SVM $-$ RF) \\
\hline
Padrão  & 2.472.104 & 1.855.094 & 354.420 & 2.534.869 & 1.500.674 \\
Ajustado & 2.368.087 & 1.706.130 & 74.492  & 3.067.778 & 1.631.638 \\
\hline
\end{tabular}}
\par\smallskip
Fonte: Elaborado pelo autor (2026).
\end{table}
 
\textbf{Validação e teste.} O ranking dos algoritmos se inverte entre a validação e o teste (Tabela~\ref{tab:comparacao-multi-validacao-teste}). Na validação, que vem do mesmo dia do treino, o RF foi melhor que o SVM em F1 macro e MCC (0,9007 e 0,8969 contra 0,8662 e 0,8721 nos padrões; 0,9174 e 0,9136 contra 0,8757 e 0,8825 nos ajustados). No teste, de outro dia, o SVM foi melhor. Os quatro modelos perderam muito (o F1 macro caiu de 0,87--0,92 para 0,41--0,61), e o RF, mais. Isso sugere que o SVM generaliza melhor para o dia de teste, mas a diferença é descritiva e o motivo não foi isolado; a validação, sozinha, indicaria a escolha oposta.
 
\begin{table}
\centering
\caption{F1 macro e MCC dos modelos multiclasse na validação e no teste}
\label{tab:comparacao-multi-validacao-teste}
\footnotesize
\begin{tabular}{lrrrr}
\hline
 & \multicolumn{2}{c}{F1 macro} & \multicolumn{2}{c}{MCC} \\
Modelo & Validação & Teste & Validação & Teste \\
\hline
RF padrão    & 0,9007 & 0,5626 & 0,8969 & 0,4324 \\
RF ajustado  & 0,9174 & 0,4148 & 0,9136 & 0,3688 \\
SVM padrão   & 0,8662 & 0,6119 & 0,8721 & 0,5347 \\
SVM ajustado & 0,8757 & 0,6022 & 0,8825 & 0,5172 \\
\hline
\end{tabular}
\par\smallskip
Fonte: Elaborado pelo autor (2026).
\end{table}
 
\textbf{Registros Portmap.} O Portmap só existe no teste (2.702 registros) e não tem classe correta possível. Os modelos o empurraram, na maioria, para o NetBIOS (Tabela~\ref{tab:comparacao-multi-portmap}), com confiança alta: a mediana da probabilidade do RF padrão nos registros previstos como NetBIOS foi 1,00, e o escore do SVM para a classe escolhida (cerca de 6,3, equivalente a ganhar os seis duelos de pares) foi praticamente igual ao dos acertos de classes conhecidas (cerca de 6,3 nos dois grupos). Isto é, nenhum dos modelos distingue um tipo desconhecido de um conhecido pela confiança, e o desempenho em classes novas não deve ser extrapolado a partir deste resultado.
 
\begin{table}
\centering
\caption{Destino dos 2.702 registros Portmap no teste}
\label{tab:comparacao-multi-portmap}
\begin{tabular}{lrr}
\hline
Modelo & Previsto como NetBIOS & Previsto como BENIGN \\
\hline
RF padrão    & 88,7\% & 4,6\% \\
RF ajustado  & 88,3\% & 8,6\% \\
SVM padrão   & 85,7\% & 4,6\% \\
SVM ajustado & 90,1\% & 3,3\% \\
\hline
\end{tabular}
\par\smallskip
Fonte: Elaborado pelo autor (2026).
\end{table}
 
\textbf{Custo e conclusão.} O ajuste de hiperparâmetros não melhorou todos os modelos. No RF, aumentou a ROC-AUC, mas reduziu F1 macro, acurácia balanceada e MCC, e piorou muito a separação entre benigno e ataque. No SVM, a configuração ajustada obteve a maior acurácia balanceada e recuperou parte do UDPLag, mas ficou abaixo do SVM padrão em F1 macro e MCC. Assim, não há evidência de que a configuração ajustada seja superior sob todos os critérios. O custo diferencia fortemente as abordagens: o RF padrão processou o teste em 61~s de previsão, enquanto o SVM padrão levou cerca de 6~h 37~min, ainda que tenha reutilizado previsões para vetores repetidos; o \textit{refit} do RF levou cerca de 70~s e o do SVM, 4.257 a 5.308~s. A Tabela~\ref{tab:comparacao-multi-sintese} resume a comparação.
 
\begin{table}
\centering
\caption{Síntese da comparação entre os modelos multiclasse}
\label{tab:comparacao-multi-sintese}
\footnotesize
\resizebox{\textwidth}{!}{%
\begin{tabular}{lll}
\hline
Critério & Melhor resultado & Observação \\
\hline
ROC-AUC OvR & RF ajustado (0,9248) & Não acompanha melhora nas métricas de classificação final \\
Acurácia balanceada & SVM ajustado (0,7335) & Maior equilíbrio médio entre recalls \\
F1 macro & SVM padrão (0,6119) & RF padrão: média de 0,4834 em cinco sementes \\
MCC & SVM padrão (0,5347) & RF padrão: média de 0,4022 em cinco sementes \\
Ataque $\times$ BENIGN (MCC) & SVM padrão (0,96017) & RF padrão equivalente (0,95613), mas instável entre sementes \\
Estabilidade entre sementes & SVM (determinístico) & RF padrão: recall de ataques de 0,72 a 0,99 \\
UDPLag (recall) & RF padrão (0,4040) & SVM padrão teve recall zero; SVM ajustado, 0,3465 \\
Tempo de previsão no teste & RF padrão (61,2~s) & SVM padrão levou cerca de 6~h 37~min \\
Validação (F1 macro, MCC) & RF ajustado (0,9174; 0,9136) & Ordem se inverte no teste \\
\hline
\end{tabular}}
\par\smallskip
Fonte: Elaborado pelo autor (2026).
\end{table}
 
\textbf{Conclusão.} Entre as configurações padrão, o SVM apresentou melhor desempenho agregado no teste, sobretudo em F1 macro e MCC, enquanto o RF foi muito mais rápido. Entre as ajustadas, o SVM superou o RF em todas as métricas agregadas, exceto a ROC-AUC. O RF ajustado não superou o RF padrão nas principais métricas agregadas. Portanto, considerando a qualidade geral das previsões e a estabilidade, o SVM padrão foi o melhor dos modelos avaliados; considerando custo e rapidez, o RF padrão é a alternativa mais eficiente, mas com resultado menos estável entre sementes. A decisão final depende de como o sistema tratará os erros por classe, em especial os de Syn (confundido com UDPLag) e UDPLag, e do custo de inferência aceitável. Essa conclusão deve ser qualificada pela diferença entre os dias de treino e de teste, pelo desempenho nulo do SVM padrão em UDPLag, pela baixa quantidade de exemplos dessa classe e pela presença de conflitos entre rótulos (Seção~\ref{sec:conflitos-multiclasse}), que deixam parte dos erros fora do alcance de qualquer modelo treinado com esses atributos.


% ---------------------------------------------------------------
% EXEMPLO DE INSERÇÃO DE FIGURA (substitua pela sua imagem)
% A imagem deve ficar na pasta assets/ (ex.: assets/figures/arquitetura.png).
% ORDEM ABNT: o título (\caption) vem ACIMA da figura e a fonte
% (\legend) vem ABAIXO. Por isso \caption e \label aparecem ANTES do
% \includegraphics. Use \ref{...} para citar a figura no texto.
% ---------------------------------------------------------------
% \begin{figure}
%     \centering
%     \caption{Exemplo de inserção de figura --- substitua pelo diagrama da arquitetura da sua solução.}
%     \label{fig:arquitetura}
%     \includegraphics[width=0.35\textwidth]{assets/logo-senac.png}
%     \legend{Fonte: Elaborado pelo autor (2025)}
% \end{figure}

% Toda figura deve ser numerada, ter legenda (\verb|\caption|) e indicação de fonte, além de ser citada no texto --- por exemplo, a Figura~\ref{fig:arquitetura}. As figuras com legenda aparecem automaticamente na \textit{Lista de ilustrações}, nos elementos pré-textuais.

% \section{Experimentos e cenário de testes (opcional)}
% \label{sec:experimentos}

% Esta seção é \textbf{opcional} e deve ser incluída quando o trabalho prevê avaliação experimental. Descreva os cenários planejados, os parâmetros testados e por que cada experimento está alinhado aos objetivos. Liste os conjuntos de dados, as variações e o número de repetições.

% \section{Resultados e análises}
% \label{sec:resultados}

% Apresente os resultados com tabelas, gráficos e figuras e, sobretudo, \textbf{analise-os} à luz da revisão bibliográfica e da hipótese --- resultados devem ser analisados, não apenas exibidos. Diferencie:
% \begin{itemize}
%   \item \textbf{Figura:} imagens, diagramas conceituais ou arquiteturais que ilustram estrutura ou fluxo.
%   \item \textbf{Gráfico:} plot de dados (linhas, barras, boxplots) que compara métricas obtidas.
% \end{itemize}

% Exemplo de tabela:
% \begin{table}
%   \centering
%   \caption{Comparação de métricas reportadas}
%   \label{tab:metricas}
%   \begin{tabular}{lcc}
%     \toprule
%     Método & Acurácia & Tempo médio (ms) \\
%     \midrule
%     Proposta & 92\% & 180 \\
%     Baseline & 88\% & 220 \\
%     \bottomrule
%   \end{tabular}
%   \legend{Fonte: Elaborado pelo autor (2025)}
% \end{table}

% Ao inserir um gráfico (por exemplo \texttt{assets/figures/desempenho.png}), use --- lembrando que, pela ABNT, o \texttt{\textbackslash caption} (título) vem \textbf{acima} da imagem e a fonte (\texttt{\textbackslash legend}) vem \textbf{abaixo}:
% \begin{verbatim}
% \begin{figure}
%   \centering
%   \caption{Gráfico de comparação das métricas de acurácia.}
%   \label{fig:grafico-desempenho}
%   \includegraphics[width=0.8\textwidth]{assets/figures/desempenho.png}
%   \legend{Fonte: Elaborado pelo autor (2025)}
% \end{figure}
% \end{verbatim}

% Na sequência, faça a \textbf{análise crítica e comparação}, discutindo as descobertas em relação aos trabalhos correlatos e evidenciando ganhos e perdas; por fim, na \textbf{discussão}, relacione os resultados aos objetivos e à hipótese, explicitando se o critério de sucesso foi atingido.

% \textit{O desenvolvimento só é pesquisa se estiver claramente ligado ao problema e à metodologia definida.}

\subsection{Experimentos e cenário de testes}
\label{sec:experimentos}

Além do treinamento, da busca de hiperparâmetros e da avaliação no teste (Seções~\ref{sec:rf-bin-treino} a~\ref{sec:comparacao-rf-svm-multi}), foram realizados dois experimentos complementares. O primeiro é a análise de explicabilidade com SHAP, que identifica os atributos que sustentam as previsões, se eles se mantêm entre os dias de treino e de teste e se RF e SVM se apoiam nos mesmos atributos. O segundo é uma simulação de sistema de detecção de intrusão, que aplica os modelos já treinados a uma amostra fixa de fluxos para demonstrar seu uso.

\begin{table}
\centering
\caption{Cenários dos experimentos de explicabilidade (SHAP)}
\label{tab:cenarios-experimentos}
\footnotesize
\resizebox{\textwidth}{!}{%
\begin{tabular}{p{3.2cm}p{5.8cm}p{5.8cm}}
\hline
 & Binário (notebooks 14 e 15) & Multiclasse (notebooks 14m e 15m) \\
\hline
Modelos & RF e SVM, padrão e ajustado & RF e SVM, padrão e ajustado \\
Amostra do SHAP (até) & 300 fluxos por grupo no RF e 60 no SVM; 6.300 fluxos por RF (3.900 do treino e 2.400 do teste) e 1.260 por SVM (780 e 480) & 150 fluxos por classe no RF e 30 no SVM; 2.100 por RF (1.050 por dia) e 420 por SVM (210 por dia) \\
Explicadores & \texttt{TreeExplainer} no RF e \texttt{PermutationExplainer} no SVM (referência de 24 fluxos do treino); no notebook 15, \texttt{PermutationExplainer} também no RF & O mesmo (referência de 7 fluxos do treino, um por classe) \\
Atributos de entrada & 43 & 44 \\
\hline
\end{tabular}}
\par\smallskip
Fonte: Elaborado pelo autor (2026).
\end{table}

\textbf{Explicabilidade com SHAP.} O SHAP distribui a saída do modelo entre os atributos de cada fluxo, em relação a uma referência. Os fluxos foram sorteados de forma estratificada por classe (ou tipo de ataque) e por dia, e a importância global de cada atributo foi calculada com peso igual para cada grupo, para que os grupos mais numerosos não dominassem o resultado. Como RF e SVM produzem saídas em escalas diferentes, as comparações usam a participação percentual de cada atributo, e não os valores absolutos. A análise foi feita nos dois dias, para verificar se a importância se mantém. A estabilidade foi avaliada por reamostragem (metades da amostra sorteadas 50 vezes, frações da amostra e \textit{bootstrap}), com limites fixados antes: correlação de postos entre metades de pelo menos 0,9, sobreposição de pelo menos 7 dos 10 principais atributos e convergência de pelo menos 0,95 com metade da amostra. O notebook 15 acrescentou quatro análises: o \texttt{PermutationExplainer} também no RF, para medir a concordância com o \texttt{TreeExplainer}; a comparação entre RF e SVM com o mesmo explicador; a contagem de atributos com valores negativos (no multiclasse, relacionada à taxa de erro no teste); e a verificação de que as previsões do SVM nos fluxos com valores negativos não dependem só do intercepto. Em todos os casos, a aditividade (referência mais soma das contribuições igual à saída do modelo) foi conferida.

\textbf{Simulação do sistema de detecção de intrusão.} Foi construída uma aplicação interativa, em Python com Streamlit, para demonstrar o funcionamento dos modelos nas tarefas binária e multiclasse. A aplicação carrega os oito modelos já treinados (RF e SVM, em configuração padrão e ajustada, nas duas tarefas), sem retreiná-los ou alterá-los.

\textit{Amostra.} Para tornar a demonstração reproduzível, foi preparada uma amostra fixa de 3.600 fluxos do conjunto de teste: 1.200 benignos e 400 de cada um dos seis tipos de ataque da tarefa multiclasse (LDAP, MSSQL, NetBIOS, Syn, UDP e UDPLag); o Portmap, que existe só no teste, não foi incluído. Os fluxos foram lidos do ClickHouse, com consultas somente de leitura, a partir das visões de teste sem sobreposição com o treino, em ordem estável definida por \textit{hash} do identificador. Só entraram fluxos presentes nas duas visões (binária e multiclasse), de modo que todos os modelos classificam os mesmos fluxos, e cada fluxo aparece no máximo uma vez. A amostra foi gravada em arquivo local, e o rótulo real é usado apenas no cálculo das métricas, nunca como atributo. Por conter 66,7\% de ataques (2.400 de 3.600), sua proporção entre as classes difere da do teste completo.

\textit{Cenários e configuração.} A interface oferece dois cenários (Tabela~\ref{tab:cenarios-simulacao}). Como as visões não contêm horários, a ordem dos fluxos é definida pelo cenário para a demonstração e não representa a cronologia da captura. Também é possível escolher a família do modelo, a configuração (padrão ou ajustada), o tamanho do lote (de 100 a 500 fluxos) e, opcionalmente, comparar RF e SVM sobre os mesmos fluxos. Os modelos operam nos limiares naturais, sem ajuste no teste: probabilidade de pelo menos 0,5 no RF binário, função de decisão maior que zero no SVM binário e, no multiclasse, a classe de maior probabilidade (RF) ou de maior pontuação um-contra-resto (SVM).

\begin{table}
\centering
\caption{Composição dos cenários da simulação}
\label{tab:cenarios-simulacao}
\footnotesize
\begin{tabular}{p{3.0cm}p{2.0cm}p{8.8cm}}
\hline
Cenário & Fluxos & Composição \\
\hline
Fases de ataque & 3.600 & 300 benignos; seis fases na ordem Syn, UDP, LDAP, MSSQL, NetBIOS e UDPLag, cada uma com 400 fluxos do ataque e 100 benignos misturados; 300 benignos ao final \\
Balanceado & 2.800 & 400 fluxos de cada uma das sete classes, embaralhados com semente fixa \\
\hline
\end{tabular}
\par\smallskip
Fonte: Elaborado pelo autor (2026).
\end{table}

\textit{Painel.} Durante a execução, o painel mostra os fluxos processados, a fase atual, a distribuição das classes previstas, a linha do tempo das detecções por lote, as classes reais e previstas, os últimos fluxos classificados e a matriz de confusão. No binário, calcula precisão, recall, F1, MCC e acurácia da classe ataque, além de falsos positivos e falsos negativos. No multiclasse, calcula acurácia, F1 macro (sobre as classes já apresentadas), MCC e precisão, recall e F1 de cada classe; nele, falso positivo é um benigno classificado como ataque, e falso negativo, um ataque classificado como benigno. O painel também mostra o tempo médio de inferência por fluxo e a velocidade em fluxos por segundo, obtidos dividindo o tempo gasto apenas na chamada de previsão de cada lote pelo número de fluxos do lote. A aplicação usa Streamlit na interface, Plotly nos gráficos, pandas e NumPy na manipulação dos dados, scikit-learn (versão 1.9.1, a mesma que salvou os modelos) para executar os modelos e joblib para armazenar os modelos e a amostra.

\textit{Alcance.} A simulação demonstra como os modelos classificam fluxos, quais erros cometem e qual é o tempo de inferência nas condições da aplicação. Seus resultados não substituem nem devem ser comparados aos do teste completo, porque usam uma amostra menor, com proporção de classes diferente e cenários de ordem construída. A latência exibida é a média por lote, e não a latência de um fluxo isolado medida na Seção~\ref{sec:comparacao-rf-svm-bin}; a velocidade definida no painel controla apenas o ritmo da apresentação.

\subsection{Resultados e análises}
\label{sec:resultados-analises}

\subsubsection{Explicabilidade com SHAP}

O cálculo do binário levou 40~min~53~s (notebook 14) e 11~min~52~s (notebook 15), e o do multiclasse, 3~h~01~min e 12~min, com um processo por vez e 1,0 a 1,8~GiB por processo, o que limitou o paralelismo. A aditividade foi reproduzida com erro máximo de $6{,}3\times10^{-7}$ (binário) e $3{,}8\times10^{-7}$ (multiclasse), e as nove verificações de cada execução foram aprovadas. A Tabela~\ref{tab:shap-principais} apresenta os atributos mais importantes.

\begin{table}
\centering
\caption{Atributo mais importante e participação de \texttt{avg\_fwd\_segment\_size} na importância global (\%)}
\label{tab:shap-principais}
\footnotesize
\resizebox{\textwidth}{!}{%
\begin{tabular}{llllcc}
\hline
Problema & Modelo & Mais importante no treino & Mais importante no teste & \multicolumn{2}{c}{\texttt{avg\_fwd\_segment\_size}} \\
 & & & & Treino & Teste \\
\hline
Binário & RF ajustado  & \texttt{avg\_fwd\_segment\_size} (8,37) & \texttt{avg\_fwd\_segment\_size} (9,25) & 8,37 & 9,25 \\
Binário & RF padrão    & \texttt{avg\_fwd\_segment\_size} (8,77) & \texttt{avg\_fwd\_segment\_size} (9,41) & 8,77 & 9,41 \\
Binário & SVM ajustado & \texttt{avg\_fwd\_segment\_size} (6,90) & \texttt{flow\_bytes\_s} (7,47) & 6,90 & 5,46 \\
Binário & SVM padrão   & \texttt{protocol} (21,97) & \texttt{protocol} (19,65) & 4,16 & 3,09 \\
\hline
Multiclasse & RF ajustado  & \texttt{avg\_fwd\_segment\_size} (10,99) & \texttt{avg\_fwd\_segment\_size} (12,43) & 10,99 & 12,43 \\
Multiclasse & RF padrão    & \texttt{avg\_fwd\_segment\_size} (10,30) & \texttt{avg\_fwd\_segment\_size} (11,44) & 10,30 & 11,44 \\
Multiclasse & SVM ajustado & \texttt{avg\_fwd\_segment\_size} (12,07) & \texttt{avg\_fwd\_segment\_size} (10,82) & 12,07 & 10,82 \\
Multiclasse & SVM padrão   & \texttt{avg\_fwd\_segment\_size} (12,74) & \texttt{fwd\_packet\_length\_std} (12,92) & 12,74 & 12,29 \\
\hline
\end{tabular}}
\par\smallskip
Fonte: Elaborado pelo autor (2026).
\end{table}

\textbf{O que os modelos usam.} O \texttt{avg\_fwd\_segment\_size} (tamanho médio do segmento no sentido de ida) foi o mais importante nos RFs, nos dois dias e nos dois problemas, e no SVM ajustado multiclasse; no SVM binário, ficou em 1º lugar no treino do modelo ajustado e entre a 6ª e a 9ª posições nos demais casos. No multiclasse, aparecem também \texttt{fwd\_packet\_length\_max} e outros atributos de tamanho de pacote. A importância do SVM ajustado é mais distribuída (no máximo 7,5\% no binário), ao passo que o SVM padrão binário concentra 22\% (treino) e 20\% (teste) em \texttt{protocol}, atributo de poucos valores distintos. Essa concentração é um fato dos dados; que ela explique o recall de 0,746 do SVM padrão no teste binário é uma hipótese, que não foi testada.

\textbf{Estabilidade.} As importâncias são estáveis dentro de cada dia: a correlação entre metades da amostra foi de 0,999 a 1,000 nos RFs e de 0,991 a 0,993 nos SVMs no binário, e de 0,998 a 0,999 e de 0,919 a 0,964 no multiclasse, respectivamente. Entre os dias, a ordem dos atributos também se manteve, mais no RF (correlação de 0,96 a 0,99) do que no SVM (0,86 a 0,95). Esse resultado exige cuidado: a ordem de importância se mantém mesmo quando o desempenho cai no teste, o que indica que os modelos usam os mesmos atributos nos dois dias. A estabilidade da explicação, portanto, não implica estabilidade do desempenho; a causa da queda não foi investigada com o SHAP.

\textbf{RF e SVM se apoiam em atributos diferentes.} A correlação de postos entre os dois algoritmos foi baixa: de 0,19 a 0,31 no binário e de 0,38 a 0,55 no multiclasse, com explicadores diferentes. Isso poderia ser um efeito do explicador, e por isso o notebook 15 aplicou o mesmo método (\texttt{PermutationExplainer}) aos dois. Trocar o explicador do RF elevou a correlação em 0,06 a 0,15 no binário e em 0,01 a 0,05 no multiclasse, mas ela permaneceu parcial (de 0,29 a 0,43 e de 0,40 a 0,59). A diferença entre os algoritmos, portanto, não é só um efeito do método de explicação. Já a escolha do explicador altera pouco o RF: a correlação entre \texttt{TreeExplainer} e \texttt{PermutationExplainer} foi de 0,88 a 0,90 no binário (7 dos 10 principais atributos em comum) e de 0,98 a 0,99 no multiclasse (9 dos 10).

\textbf{Valores negativos.} Os atributos \texttt{fwd\_header\_length} e \texttt{min\_seg\_size\_forward} têm valores negativos em parte das linhas (valores da ordem de $-2\times10^{9}$ em \texttt{fwd\_header\_length}), o que sugere um problema de captura, não verificado. No SVM, as linhas com esses valores mantiveram vetores de suporte próximos e não foram classificadas só pelo intercepto. No multiclasse, a taxa de erro do RF ajustado nessas linhas varia por classe: no MSSQL, foi de 92,82\% entre as linhas com \texttt{fwd\_header\_length} negativo e de 15,78\% entre as demais; no LDAP, a relação é inversa (0,23\% contra 5,56\%). Como o sentido da associação muda de uma classe para outra, ela não sustenta uma explicação causal e indica apenas onde investigar a captura dos dados.

\begin{figure}
\centering
\caption{Importância global dos atributos no RF ajustado binário}
\label{fig:shap-rf-bin}
\includegraphics[width=\textwidth]{assets/shap_bin_importancia_global_rf_ajustado.png}
\legend{Fonte: Elaborado pelo autor (2026).}
\end{figure}

\begin{figure}
\centering
\caption{Importância global dos atributos no SVM ajustado binário}
\label{fig:shap-svm-bin}
\includegraphics[width=\textwidth]{assets/shap_bin_importancia_global_svm_ajustado.png}
\legend{Fonte: Elaborado pelo autor (2026).}
\end{figure}

\begin{figure}
\centering
\caption{Importância global dos atributos no RF ajustado multiclasse}
\label{fig:shap-rf-multi}
\includegraphics[width=\textwidth]{assets/shap_multi_importancia_global_rf_ajustado.png}
\legend{Fonte: Elaborado pelo autor (2026).}
\end{figure}

\begin{figure}
\centering
\caption{Importância global dos atributos no SVM ajustado multiclasse}
\label{fig:shap-svm-multi}
\includegraphics[width=\textwidth]{assets/shap_multi_importancia_global_svm_ajustado.png}
\legend{Fonte: Elaborado pelo autor (2026).}
\end{figure}

\subsubsection{Demonstração com a simulação}

A simulação (Seção~\ref{sec:experimentos}) ilustra o uso dos modelos, e não produz resultados para a comparação. % Após inserir as figuras abaixo, acrescentar: A Figura~\ref{fig:sim-painel} mostra o painel no cenário de fases de ataque, em que a linha do tempo permite acompanhar a classificação de cada fase, e a Figura~\ref{fig:sim-comparacao} mostra a comparação entre RF e SVM sobre os mesmos fluxos, com as matrizes de confusão e o custo de inferência.
O painel permite acompanhar a classificação de cada fase em uma linha do tempo e comparar RF e SVM sobre os mesmos fluxos, com matrizes de confusão e custo de inferência. Os valores exibidos pertencem à amostra de 3.600 fluxos e à ordem construída do cenário; por isso, não são comparados aos do teste completo.

% Inserir as capturas de tela depois de executar a simulação e salvá-las em assets/:
% \begin{figure}[htb]
% \centering
% \caption{Painel da simulação no cenário de fases de ataque (tarefa multiclasse)}
% \label{fig:sim-painel}
% \includegraphics[width=\textwidth]{assets/simulacao_painel_multiclasse.png}
% \legend{Fonte: Elaborado pelo autor (2026).}
% \end{figure}
%
% \begin{figure}[htb]
% \centering
% \caption{Comparação entre RF e SVM sobre os mesmos fluxos na simulação}
% \label{fig:sim-comparacao}
% \includegraphics[width=\textwidth]{assets/simulacao_comparacao_rf_svm.png}
% \legend{Fonte: Elaborado pelo autor (2026).}
% \end{figure}

\subsubsection{Análise crítica e discussão}

\textbf{O que os resultados mostram em conjunto.} A comparação entre RF e SVM depende do critério e do problema, e nenhum modelo foi superior em tudo. No binário, o SVM ajustado teve melhor discriminação e melhor resultado no limiar natural (MCC de 0,82075 contra 0,36144), mas o RF foi muito melhor no ponto de 0,1\% de falso alarme (MCC de 0,96682 contra 0,10496) e processou o teste em 52,3~s contra 1.227,6~s (Seção~\ref{sec:comparacao-rf-svm-bin}). No multiclasse, o SVM padrão obteve os melhores F1 macro (0,6119) e MCC (0,5347), contra 0,5626 e 0,4324 do RF padrão, e esse RF oscila muito entre sementes (média de 0,4834 e 0,4022); em contrapartida, o RF padrão previu o teste em 61~s e o SVM padrão em cerca de 6~h~37~min (Seção~\ref{sec:comparacao-rf-svm-multi}). Ao contrário do que a validação indicava, o RF ajustado, o melhor na validação, teve o pior F1 macro no teste (0,4148).

\textbf{Ganhos e perdas.} O SVM ganhou em estabilidade (não há sorteio entre sementes) e em desempenho agregado no multiclasse; perdeu em custo: o ajuste final foi cerca de 6 vezes mais lento no binário (126,8~s contra 20,9~s) e 76 vezes no multiclasse, e a previsão do teste foi de 23 a cerca de 390 vezes mais lenta (cálculos a partir dos tempos). O RF ganhou em custo e em vazão em lote; perdeu em robustez à semente. O ajuste de hiperparâmetros, escolhido pela validação cruzada, teve no teste um efeito grande no SVM binário (MCC de 0,1477 para 0,82075 no limiar natural), pequeno no RF binário (0,36144 contra 0,35521) e negativo no multiclasse; no RF binário, aumentou o custo de treino em cerca de 2,9 vezes, ao passo que no multiclasse o reduziu (de 3.483.364 para 1.045.332 nós). Isso é coerente com a busca ter usado dados do mesmo dia do treino, o que não permitia corrigir uma diferença entre dias (Seção~\ref{sec:rf-bin-teste}).

\textbf{Em relação aos trabalhos correlatos.} A Tabela~\ref{tab:correlatos} resume três estudos que compararam RF e SVM, com os dados conforme descritos nos artigos.

\begin{table}
\centering
\caption{Comparação com os trabalhos correlatos}
\label{tab:correlatos}
\footnotesize
\resizebox{\textwidth}{!}{%
\begin{tabular}{p{3.2cm}p{3.4cm}p{2.6cm}p{5.6cm}}
\hline
Estudo & Dados & Divisão & Resultado (RF $\times$ SVM) \\
\hline
\citeonline{rastogi2022} & NSL-KDD, binário e multiclasse & 75\%/25\% & Acurácia de 98,70\% $\times$ 96,70\% (binário) e 98,47\% $\times$ 95,25\% (multiclasse); treino de 3,5~s $\times$ 232,3~s e teste de 0,13~s $\times$ 11,4~s (binário) \\
\citeonline{abiramasundari2025} & CICIDS2017, CICIDS2018 e CICDDoS2019 & 70\%/30\% & No CICDDoS2019, RF com precisão, recall e F1 de 0,98 e SVM de 0,95, 0,97 e 0,96 \\
\citeonline{hirsi2024} & Conjunto próprio em SDN (referência: CICDDoS2019) & 80\%/20\% & Acurácia de 98,97\% $\times$ 96,52\% e falso alarme de 0,023 $\times$ 0,036; no CICDDoS2019, o RF teve acurácia de 97,62\% e F1 de 97,69\% \\
Este trabalho (teste, binário) & CICDDoS2019, dias distintos & 12/01/2019 e 11/03/2019 & Limiar natural: recall de 0,95259 $\times$ 0,99642 e MCC de 0,36144 $\times$ 0,82075; treino final de 20,9~s $\times$ 126,8~s \\
\hline
\end{tabular}}
\par\smallskip
Fonte: Elaborado pelo autor (2026).
\end{table}

Há coerência no custo: \citeonline{rastogi2022}, único dos três a informar tempos nos trechos consultados, encontraram o RF muito mais rápido no treino (cerca de 66 vezes, cálculo a partir dos tempos) e na previsão, como aqui (cerca de 6 vezes no binário e 76 vezes no multiclasse, também calculado). A divergência está no desempenho: nos estudos, o RF obteve resultados iguais ou superiores aos do SVM, enquanto aqui o SVM superou o RF no limiar natural do binário e no multiclasse com a configuração padrão. Uma explicação possível, não testada, é a forma de avaliação: nos trechos consultados, os artigos dividem os dados em proporções fixas e não mencionam capturas de dias distintos, e informam principalmente acurácia, precisão, recall e F1, sem MCC. Neste trabalho, o mesmo tipo de avaliação (validação do mesmo dia) também deu resultados altos (MCC sem ponderação de 0,997 no RF e 0,988 no SVM, Seções~\ref{sec:rf-bin-treino} e \ref{sec:comparacao-rf-svm-bin}), e a queda apareceu só no dia de teste. Os números, portanto, não são diretamente comparáveis, por diferenças de conjunto, divisão e métricas, e a comparação é qualitativa.

\textbf{Relação com a hipótese.} A hipótese previa (i) diferenças de desempenho entre RF e SVM, nas métricas e no custo computacional; (ii) que essas diferenças variassem entre a classificação binária e a multiclasse; e (iii) que a otimização de hiperparâmetros melhorasse os resultados, com efeito variável entre os modelos. As três previsões foram sustentadas, com ressalvas.
\begin{enumerate}
\item \textit{Diferenças de desempenho e de custo.} Existem, mas dependem da métrica e do ponto de operação. As métricas que não dependem da proporção entre as classes (recall, falso alarme, MCC) separaram os modelos melhor que a acurácia e a precisão dos ataques, que ficam próximas de 1 mesmo para modelos com muitos erros na classe BENIGN. No custo, o RF foi mais rápido no treino e na previsão em lote nos dois problemas; o SVM teve menor latência por fluxo isolado no binário (1,410~ms contra 29,693~ms).
\item \textit{Variação entre binário e multiclasse.} Confirmada. No binário, o SVM ajustado superou o RF no limiar natural (MCC de 0,82075 contra 0,36144), e o RF foi superior no ponto de 0,1\% de falso alarme; no multiclasse, o SVM padrão teve os melhores F1 macro e MCC, e o RF ajustado, o pior desempenho.
\item \textit{Efeito da otimização de hiperparâmetros.} Confirmado o efeito variável, e não o ganho generalizado: a configuração foi escolhida pela validação cruzada, mas, no teste, o ganho foi grande no SVM binário (MCC de 0,1477 para 0,82075 no limiar natural), pequeno no RF binário (0,35521 para 0,36144) e negativo no RF multiclasse (F1 macro de 0,5626 para 0,4148) e no SVM multiclasse (0,6119 para 0,6022). A otimização foi feita com dados do dia de treino e, por isso, não corrige a diferença entre os dias.
\end{enumerate}
\textbf{Relação com os objetivos.} O objetivo principal, comparar RF e SVM na identificação de ataques DDoS, foi atendido, e os específicos foram cumpridos como segue.
\begin{enumerate}
\item \textit{Implementar e treinar RF e SVM.} Atendido: os dois algoritmos foram treinados nos problemas binário e multiclasse, em configuração padrão e ajustada, sobre os mesmos dados.
\item \textit{Desempenho antes e depois da otimização.} Atendido: a busca com validação cruzada foi feita nos quatro casos e comparada no teste; o efeito variou entre os modelos (item 3 da hipótese).
\item \textit{Binário e multiclasse, com custo computacional.} Atendido: as duas tarefas foram medidas com tempos de treino, de previsão em lote e de latência por fluxo.
\item \textit{Comparação com o estado da arte.} Atendido em parte: a comparação com os três estudos é qualitativa, porque eles usam outros conjuntos, divisões e métricas, e nenhum, nos trechos consultados, reporta MCC.
\end{enumerate}
Considerando o critério de sucesso, os objetivos 1 a 3 foram plenamente alcançados e o 4, de forma parcial; por isso, a comparação com a literatura deve ser lida como contextual, e não como superação ou equivalência numérica.

Quanto às vantagens e limitações, o SVM mostrou melhor desempenho agregado e estabilidade entre sementes, ao custo de tempos de previsão muito maiores; o RF mostrou menor custo, mas resultados mais sensíveis à semente. Esses resultados valem para este conjunto de dados, não permitem afirmar que um algoritmo seja superior em geral, e dependem de o teste ser uma única partição, reutilizada em várias etapas.

\textbf{Limitações.} O teste não foi usado uma única vez e deixa de ser uma avaliação independente a partir da primeira avaliação (Seção~\ref{sec:rf-bin-sementes}); os limiares da validação não se transferiram para o teste; os conflitos entre rótulos limitam o desempenho de qualquer modelo no multiclasse; o SHAP foi calculado em amostras (de 420 a 6.300 fluxos por modelo), com o SVM em amostras menores que as do RF. A simulação de IDS, por fim, é uma demonstração com 3.600 fluxos (2.800 no cenário balanceado) e não deve ser comparada ao teste completo.

% ---------------------------------------------------------------
% CONCLUSÃO
% ---------------------------------------------------------------
\chapter{Conclusão}

\section{Retomada do problema e dos objetivos}
Retome o problema e destaque o que foi atingido.

\section{Contribuições para a Computação}
Explique o que a pesquisa trouxe de novo.

\section{Limitações e trabalhos futuros}
Liste limites e sugira próximos passos coerentes.

% ---------------------------------------------------------------
% OBSERVAÇÕES DO AUTOR
% ---------------------------------------------------------------
\chapter*{Observações importantes}
\addcontentsline{toc}{chapter}{Observações importantes}

Para graduação, Wazlawick aceita trabalhos aplicados desde que:
\begin{itemize}
    \item O problema seja preciso;
    \item A solução esteja fundamentada;
    \item Haja análise crítica;
    \item Não seja apenas entrega funcional.
\end{itemize}

% ---------------------------------------------------------------
% APRESENTAÇÃO ORAL (BANCA)
% ---------------------------------------------------------------
% \chapter*{Apresentação oral (banca)}
% \addcontentsline{toc}{chapter}{Apresentação oral (banca)}

% Além do documento escrito, o TCC2 é avaliado por uma \textbf{apresentação oral} final perante a banca. Um template pronto em \texttt{Beamer} está disponível no arquivo \texttt{apresentacao\_tcc2.tex}, na mesma pasta deste documento.

% \noindent\fbox{\parbox{\dimexpr\textwidth-2\fboxsep-2\fboxrule}{%
% \textbf{Regras da apresentação do TCC2}

% \vspace{0.2cm}
% \begin{itemize}
%     \item \textbf{Duração: 30 minutos.} Cronometre os ensaios; ultrapassar o tempo é penalizado. A arguição da banca ($\approx$ 5 min) ocorre após a apresentação, separada desse total.
%     \item \textbf{Cubra o trabalho completo}, seguindo a estrutura do documento escrito: introdução e problema, objetivos e método, desenvolvimento, \textbf{resultados e análise} e conclusão/contribuições.
%     \item \textbf{Dê peso aos resultados:} grande parte do tempo deve ser dedicada a apresentar e \textbf{analisar criticamente} os resultados, não apenas descrevê-los.
%     \item \textbf{Ritmo:} cerca de 1 a 1,5 minuto por slide, totalizando aproximadamente 20 a 26 slides de conteúdo.
%     \item \textbf{Logo do SENAC:} grande no primeiro slide (capa) e pequena no canto superior direito dos demais slides --- já configurado no template.
%     \item \textbf{Demonstração:} se houver demo ao vivo, tenha um plano B com capturas de tela ou vídeo curto.
%     \item \textbf{Slides de apoio (backup):} inclua slides extras após o ``Obrigado'' com detalhes técnicos, trechos de código e tabelas completas para a arguição.
% \end{itemize}

% \vspace{0.2cm}
% \textbf{Distribuição de tempo sugerida (30 min):} introdução/problema $\approx$ 4 min; objetivos/método $\approx$ 5 min; desenvolvimento $\approx$ 8 min; resultados/análise $\approx$ 8 min; conclusão $\approx$ 3 min; demonstração $\approx$ 2 min.
% }}

% ---------------------------------------------------------------
% ELEMENTOS PÓS-TEXTUAIS
% ---------------------------------------------------------------
\postextual

% ===============================================================
% REFERÊNCIAS BIBLIOGRÁFICAS
% ===============================================================
% As referências bibliográficas são geradas automaticamente a partir
% do arquivo bibliografia_tcc2.bib, seguindo o padrão ABNT (NBR 6023).
%
% COMO CITAR no texto:
%   \cite{chave}         -> (SOBRENOME, ano) — citação indireta
%   \citeonline{chave}   -> Sobrenome (ano) — citação direta no texto
%
% TIPOS DE REFERÊNCIA no arquivo .bib:
%   @book        -> Livros
%   @article     -> Artigos de periódicos
%   @inproceedings -> Artigos em anais de eventos
%   @mastersthesis -> Dissertações de mestrado
%   @phdthesis   -> Teses de doutorado
%   @misc        -> Sites, vídeos e outros materiais
%   @manual      -> Documentação técnica
% ---------------------------------------------------------------
\bibliography{bibliografia_tcc2}

% ---------------------------------------------------------------
% APÊNDICES (se necessário)
% ---------------------------------------------------------------
% \begin{apendicesenv}
% \partapendices
%
% \chapter{Nome do Apêndice}
% \label{ap:apendice1}
%
% Conteúdo do apêndice (material elaborado pelo próprio autor).
%
% \end{apendicesenv}

% ---------------------------------------------------------------
% ANEXOS (se necessário)
% ---------------------------------------------------------------
% \begin{anexosenv}
% \partanexos
%
% \chapter{Nome do Anexo}
% \label{an:anexo1}
%
% Conteúdo do anexo (material de terceiros).
%
% \end{anexosenv}

\end{document}
